% Lectura y comentario : oficios, mundo del trabajo, derechos y deberes de los trabajadores — Espagnol Tle A — Cours signé OnBuch+
% Programme officiel MINESEC. Compiler avec : tectonic E20-oficios-derechos-deberes-trabajadores.tex
\newcommand{\DOCMATIERE}{Espagnol}
\newcommand{\DOCNIVEAU}{Tle A}
\newcommand{\DOCMODULE}{Módulo 4 : Activités économiques et monde du travail}
\newcommand{\DOCLECON}{E20}
\newcommand{\DOCTITRE}{Lectura y comentario : oficios, mundo del trabajo, derechos y deberes de los trabajadores}
% OnBuch+ — préambule « cours signés » ENRICHI (compilé avec tectonic / XeLaTeX)
% Système visuel « Pop » d'OnBuch+ : fond crème (jamais blanc), bordures encre
% épaisses, ombres « dures » décalées, titres Archivo Black en capitales.
% Version MATHÉMATIQUES : graphes (pgfplots/tikz), boîtes pédagogiques
% (définition, propriété, méthode, attention, le savais-tu, exercice résolu,
% exercice, TP, bilan), page de garde et signature OnBuch+ ancrée en bas.
%
% Macros attendues dans main.tex AVANT \input{preamble} :
%   \DOCMATIERE  \DOCNIVEAU  \DOCMODULE  \DOCLECON (numéro)  \DOCTITRE
\documentclass[11pt]{article}

\usepackage{fontspec}
\usepackage[spanish,french]{babel} % français = langue principale ; l'espagnol via \esp{..} / espagnol
\usepackage[a4paper,margin=2cm,top=3.4cm,bottom=2.8cm,headheight=1.4cm,footskip=1.3cm]{geometry}
\usepackage{amsmath,amssymb,amsfonts}
\usepackage{mathtools} % \xmapsto, \coloneqq... souvent utilisés par les modèles
\usepackage{cancel} % simplifications barrées (bilans de forces)
% \abs{z} (module, valeur absolue) et \norm{u} (norme), utilisés par les modèles
\providecommand{\abs}{}\let\abs\relax\DeclarePairedDelimiter{\abs}{\lvert}{\rvert}
\providecommand{\norm}{}\let\norm\relax\DeclarePairedDelimiter{\norm}{\lVert}{\rVert}
\usepackage[table]{xcolor} % [table] : fournit aussi \rowcolors (alternance de lignes)
\usepackage{pagecolor}
\usepackage{tikz}
\usetikzlibrary{arrows.meta,positioning,calc,shapes.geometric,shapes.misc,shapes.arrows,decorations.pathmorphing,decorations.pathreplacing,decorations.markings,patterns,shadows,mindmap,trees,fit,backgrounds,matrix,intersections,angles,quotes}
\usepackage{pgfplots}
\usepackage[version=4]{mhchem} % équations nucléaires \ce{^{235}_{92}U}
\usepackage[european,siunitx]{circuitikz} % schémas électriques
\pgfplotsset{compat=1.18}
\usepgfplotslibrary{fillbetween,groupplots} % aires entre courbes (calcul intégral)
\usepackage{enumitem}
\usepackage{fancyhdr}
% [most] charge toutes les bibliothèques tcolorbox (listings, minted, raster,
% documentation...) et gonfle inutilement la mémoire TeX sur un document long
% (~300 boîtes) : on ne charge que ce qui est réellement utilisé.
\usepackage[skins,breakable]{tcolorbox}
\usepackage{graphicx}
\usepackage{colortbl}
\usepackage{booktabs}
\usepackage{tabularx}
\usepackage{diagbox} % case d'en-tête diagonale des tableaux à double entrée
% Colonnes extensibles centrée / gauche / droite (C, L, R), souvent utilisées par les modèles
\newcolumntype{C}{>{\centering\arraybackslash}X}
\newcolumntype{L}{>{\raggedright\arraybackslash}X}
\newcolumntype{R}{>{\raggedleft\arraybackslash}X}
\usepackage{array}
\usepackage{multicol}
\usepackage{siunitx}
% parse-numbers=false : accepte aussi \SI{2,5 \times 10^{-3}}{...}
\sisetup{locale=FR,output-decimal-marker={,},group-minimum-digits=4,parse-numbers=false}
\DeclareSIUnit\ohm{\ensuremath{\symup{\Omega}}} % U+2126 absent de la police
\DeclareSIUnit\division{div} % réglages d'oscilloscope (ms/div, V/div)
\DeclareSIUnit\tr{tr} % tours (vitesse de rotation, ex. \SI{1500}{\tr\per\minute})
\DeclareSIUnit\molar{M} % concentration molaire (ex. \SI{3}{\molar}, \SI{1}{\micro\molar})
\DeclareSIUnit\an{an} % année (le modèle confond parfois avec \year, un compteur TeX) — ex. \SI{5}{\kilo\meter\per\an}
\DeclareSIUnit\FCFA{FCFA} % franc CFA (ex. \SI{500}{\FCFA\per\kilo\gram})
\DeclareSIUnit\degreeBrix{\textdegree Brix} % teneur en sucre d'un sirop/jus (réfractométrie)
\DeclareSIUnit\month{mois} % le modèle utilise parfois \month, un compteur TeX, comme unité de durée
\usepackage{qrcode}
% \degree : commande fréquemment utilisée par les modèles pour un angle ou une
% température en dehors de \SI{}{} (non fournie par les paquets chargés).
\providecommand{\degree}{^{\circ}}
% \qed (fin de démonstration), utilisé par les modèles sans amsthm
\providecommand{\qed}{\hfill$\square$}
% \barre : double barre des valeurs interdites dans un tableau de variations (tkz-tab)
\providecommand{\barre}{\ensuremath{\|}}
% environnement proof (amsthm non chargé) utilisé par les modèles
\ifdefined\proof\else\newenvironment{proof}[1][Démonstration]{\par\noindent\textit{#1.}\ }{\qed\par}\fi
\usepackage{lastpage}
\usepackage{hyperref}
\hypersetup{hidelinks}

% ============================================================
% Palette « Pop » OnBuch+ — mêmes hex que la classe P (plus_ui.dart)
% ============================================================
\definecolor{poporange}{HTML}{F59321}
\definecolor{popdark}{HTML}{E07A0C}
\definecolor{popcream}{HTML}{FFF6E8}
\definecolor{popink}{HTML}{1C1714}
\definecolor{poppurple}{HTML}{7A5AE0}
\definecolor{popgreen}{HTML}{2E9E6A}
\definecolor{poppink}{HTML}{E0457B}
\definecolor{popblue}{HTML}{2F7DE1}
\definecolor{popgold}{HTML}{C98A12}
\definecolor{popmuted}{HTML}{8A7F76}
\definecolor{popline}{HTML}{EADFD2}
\definecolor{popgray}{HTML}{8A7F76}
\colorlet{popgrey}{popgray}
% Alias tolérants (le modèle nomme parfois les couleurs différemment)
\colorlet{popwhite}{white}
\colorlet{popblack}{popink}
\colorlet{popred}{poppink}
\colorlet{popyellow}{popgold}
% Teintes claires (fonds de tableaux/encadrés) — le modèle nomme parfois
% les couleurs avec un suffixe L (ex. popblueL) sans les avoir définies.
\colorlet{poporangeL}{poporange!20}
\colorlet{popdarkL}{popdark!20}
\colorlet{poppurpleL}{poppurple!20}
\colorlet{popgreenL}{popgreen!20}
\colorlet{poppinkL}{poppink!20}
\colorlet{popblueL}{popblue!20}
\colorlet{popgoldL}{popgold!20}
\colorlet{popredL}{popred!20}
\colorlet{popgrayL}{popgray!20}
% Teintes claires pour fonds de tableaux / zones
\colorlet{popblueL}{popblue!10!white}
\colorlet{poporangeL}{poporange!14!white}
\colorlet{popgreenL}{popgreen!12!white}
\colorlet{poppurpleL}{poppurple!12!white}
\colorlet{poppinkL}{poppink!12!white}
\colorlet{popgoldL}{popgold!14!white}

\pagecolor{popcream}

% ============================================================
% Typographie
% ============================================================
\defaultfontfeatures{Ligatures=TeX}
\setmainfont{PlusJakartaSans-Medium.ttf}[Path=../../../fonts/,BoldFont=PlusJakartaSans-Bold.ttf]
% Maths en sans-serif (Fira Math) avec chiffres et lettres droites de Plus
% Jakarta, pour que les formules se fondent dans le texte.
\usepackage[math-style=ISO,bold-style=ISO]{unicode-math}
\setmathfont{FiraMath-Regular.otf}[Path=../../../fonts/]
\setmathfont{PlusJakartaSans-Medium.ttf}[Path=../../../fonts/,range={up/num,up/{Latin,latin},\mathup}]
\usepackage{icomma} % virgule décimale sans espace en mode math
% Caractères mathématiques Unicode absents des polices de texte (Plus Jakarta,
% Archivo Black) : rendus par leur équivalent mathématique, en texte comme en
% formule — sinon ils s'affichent en carré vide (ex. « Arithmétique dans ℤ »).
\usepackage{newunicodechar}
\newunicodechar{ℝ}{\ensuremath{\mathbb{R}}}
\newunicodechar{ℤ}{\ensuremath{\mathbb{Z}}}
\newunicodechar{ℂ}{\ensuremath{\mathbb{C}}}
\newunicodechar{ℕ}{\ensuremath{\mathbb{N}}}
\newunicodechar{ℚ}{\ensuremath{\mathbb{Q}}}
\newunicodechar{𝔻}{\ensuremath{\mathbb{D}}}
\newunicodechar{α}{\ensuremath{\alpha}}
\newunicodechar{β}{\ensuremath{\beta}}
\newunicodechar{γ}{\ensuremath{\gamma}}
\newunicodechar{δ}{\ensuremath{\delta}}
\newunicodechar{ε}{\ensuremath{\varepsilon}}
\newunicodechar{θ}{\ensuremath{\theta}}
\newunicodechar{λ}{\ensuremath{\lambda}}
\newunicodechar{μ}{\ensuremath{\mu}}
\newunicodechar{σ}{\ensuremath{\sigma}}
\newunicodechar{τ}{\ensuremath{\tau}}
\newunicodechar{φ}{\ensuremath{\varphi}}
\newunicodechar{ψ}{\ensuremath{\psi}}
\newunicodechar{ω}{\ensuremath{\omega}}
\newunicodechar{Σ}{\ensuremath{\Sigma}}
% majuscules produites par \MakeUppercase dans les titres (α-aminés -> Α)
\newunicodechar{Α}{\ensuremath{\alpha}}
\newunicodechar{Β}{\ensuremath{\beta}}
\newunicodechar{Γ}{\ensuremath{\Gamma}}
\newunicodechar{Δ}{\ensuremath{\Delta}}
\newunicodechar{Ε}{\ensuremath{\varepsilon}}
\newunicodechar{Θ}{\ensuremath{\Theta}}
\newunicodechar{Λ}{\ensuremath{\Lambda}}
\newunicodechar{Μ}{\ensuremath{\mu}}
\newunicodechar{Τ}{\ensuremath{\tau}}
\newunicodechar{Φ}{\ensuremath{\Phi}}
\newunicodechar{Ψ}{\ensuremath{\Psi}}
\newunicodechar{Ω}{\ensuremath{\Omega}}
\newunicodechar{↦}{\ensuremath{\mapsto}}
\newunicodechar{∈}{\ensuremath{\in}}
\newunicodechar{∉}{\ensuremath{\notin}}
\newunicodechar{∧}{\ensuremath{\wedge}}
\newunicodechar{⇔}{\ensuremath{\Leftrightarrow}}
\newunicodechar{⟺}{\ensuremath{\Longleftrightarrow}}
\newunicodechar{⇒}{\ensuremath{\Rightarrow}}
\newunicodechar{⟹}{\ensuremath{\Longrightarrow}}
\newunicodechar{∘}{\ensuremath{\circ}}
\newunicodechar{∪}{\ensuremath{\cup}}
\newunicodechar{∩}{\ensuremath{\cap}}
\newunicodechar{≡}{\ensuremath{\equiv}}
\newunicodechar{⊂}{\ensuremath{\subset}}
\newunicodechar{⊥}{\ensuremath{\perp}}
\newunicodechar{∃}{\ensuremath{\exists}}
\newunicodechar{∀}{\ensuremath{\forall}}
\newunicodechar{∛}{\ensuremath{\sqrt[3]{\,}}}
\newunicodechar{∜}{\ensuremath{\sqrt[4]{\,}}}
\newunicodechar{⃗}{\ensuremath{{}^{\to}}}
\newunicodechar{ⁿ}{\ifmmode^{n}\else\textsuperscript{\ensuremath{n}}\fi}
\newunicodechar{ˣ}{\ifmmode^{x}\else\textsuperscript{\ensuremath{x}}\fi}
\newunicodechar{ᵇ}{\ifmmode^{b}\else\textsuperscript{\ensuremath{b}}\fi}
\newunicodechar{ʳ}{\ifmmode^{r}\else\textsuperscript{\ensuremath{r}}\fi}
\newunicodechar{ᵘ}{\ifmmode^{u}\else\textsuperscript{\ensuremath{u}}\fi}
\newunicodechar{ʸ}{\ifmmode^{y}\else\textsuperscript{\ensuremath{y}}\fi}
\newunicodechar{ᵃ}{\ifmmode^{a}\else\textsuperscript{\ensuremath{a}}\fi}
\newunicodechar{ᵉ}{\ifmmode^{e}\else\textsuperscript{\ensuremath{e}}\fi}
\newunicodechar{ᵏ}{\ifmmode^{k}\else\textsuperscript{\ensuremath{k}}\fi}
\newunicodechar{ᵖ}{\ifmmode^{p}\else\textsuperscript{\ensuremath{p}}\fi}
\newunicodechar{ᴺ}{\ifmmode^{N}\else\textsuperscript{\ensuremath{N}}\fi}
\newunicodechar{ᶜ}{\ifmmode^{c}\else\textsuperscript{\ensuremath{c}}\fi}
\newunicodechar{ʰ}{\ifmmode^{h}\else\textsuperscript{\ensuremath{h}}\fi}
\newunicodechar{ᵠ}{\ifmmode^{q}\else\textsuperscript{\ensuremath{q}}\fi}
\newunicodechar{ₙ}{\ifmmode_{n}\else\textsubscript{\ensuremath{n}}\fi}
\newunicodechar{ₐ}{\ifmmode_{a}\else\textsubscript{\ensuremath{a}}\fi}
\newunicodechar{ᵢ}{\ifmmode_{i}\else\textsubscript{\ensuremath{i}}\fi}
\newunicodechar{ᵣ}{\ifmmode_{r}\else\textsubscript{\ensuremath{r}}\fi}
\newunicodechar{ₖ}{\ifmmode_{k}\else\textsubscript{\ensuremath{k}}\fi}
\newunicodechar{ᵦ}{\ifmmode_{\beta}\else\textsubscript{\ensuremath{\beta}}\fi}
\newunicodechar{ₓ}{\ifmmode_{x}\else\textsubscript{\ensuremath{x}}\fi}
\newfontfamily\popdisplay{ArchivoBlack-Regular.ttf}[Path=../../../fonts/]
\newfontfamily\poplogo{SpaceGrotesk-Bold.ttf}[Path=../../../fonts/]
\newfontfamily\popsemi{PlusJakartaSans-SemiBold.ttf}[Path=../../../fonts/]
\newfontfamily\popxbold{PlusJakartaSans-ExtraBold.ttf}[Path=../../../fonts/]
\newfontfamily\popxboldls{PlusJakartaSans-ExtraBold.ttf}[Path=../../../fonts/,LetterSpace=3.0]

\setlist[enumerate]{leftmargin=1.7em,itemsep=5pt,topsep=4pt}
\setlist[itemize]{leftmargin=1.7em,itemsep=4pt,topsep=4pt,label={\color{poporange}\textbullet}}
\setlength{\parindent}{0pt}
\setlength{\parskip}{5pt}
\renewcommand{\arraystretch}{1.25}

% pgfplots : style Pop par défaut
\pgfplotsset{
  every axis/.append style={axis line style={popink,line width=1pt},tick style={popink},
    grid style={popline,line width=0.5pt},label style={font=\small},tick label style={font=\footnotesize}},
  popcurve/.style={poporange,line width=1.8pt,no markers,smooth},
}
% Même style disponible dans un \draw TikZ simple (hors pgfplots)
\tikzset{popcurve/.style={poporange,line width=1.8pt}}
\tikzset{popgrid/.style={popline,very thin}}
\colorlet{popcurve}{poporange} % le nom du style est parfois utilisé comme couleur

% ============================================================
% Pill / squiggle
% ============================================================
\newcommand{\poppill}[3]{%
  \tikz[baseline=-0.55ex]\node[draw=popink,line width=1.2pt,rounded corners=2.6pt,%
    fill=#2,inner xsep=6.5pt,inner ysep=3pt]{%
    \popxboldls\fontsize{7}{7}\selectfont\color{#3}\MakeUppercase{#1}};%
}
\newcommand{\popsquiggle}[1]{%
  \tikz[baseline=-0.4ex]{
    \draw[#1,line width=2.6pt,line cap=round]
      plot [smooth] coordinates {(0,0.09) (0.34,0.01) (0.68,0.21) (1.02,0.01) (1.36,0.10)};
  }%
}
% Mot-clé surligné (marqueur orange sous le texte)
\newcommand{\cle}[1]{{\popxbold\color{popdark}#1}}

% ============================================================
% En-tête / pied
% ============================================================
\pagestyle{fancy}
\fancyhf{}
\renewcommand{\headrulewidth}{0pt}
\renewcommand{\footrulewidth}{0pt}
\lhead{\raisebox{-0.15ex}{\poplogo\fontsize{13}{13}\selectfont\color{popink}OnBuch\color{poporange}+}}
\rhead{\poppill{\DOCMATIERE\ \textbullet\ \DOCNIVEAU\ \textbullet\ Leçon \DOCLECON}{white}{popink}}
\lfoot{\popxbold\fontsize{7.6}{7.6}\selectfont\color{popmuted}\MakeUppercase{Cours signé OnBuch+}\ \textbullet\ {\popsemi\DOCTITRE}}
\rfoot{\tikz[baseline=-0.6ex]\node[rounded corners=8.5pt,fill=popink,minimum height=17pt,minimum width=17pt,inner xsep=5pt,inner sep=0]{\popxbold\fontsize{8}{8}\selectfont\color{popcream}\thepage\,/\,\pageref*{LastPage}};}

% ============================================================
% Titres
% ============================================================
\newcommand{\chaptitre}[2]{%
  \begin{tcolorbox}[enhanced,colback=poporange,colframe=popink,boxrule=2.6pt,arc=11pt,
      drop shadow=popink,left=15pt,right=15pt,top=12pt,bottom=11pt,before skip=3pt,after skip=18pt]
    \poppill{#2}{popink}{popcream}\par\vspace{9pt}
    {\raggedright\hyphenpenalty=10000\exhyphenpenalty=10000\popdisplay\fontsize{22}{24}\selectfont\color{popcream}\MakeUppercase{#1}\par}
  \end{tcolorbox}
}
% Section numérotée : pastille orange avec le numéro + titre Archivo
\newcounter{popsec}
\newcommand{\coursec}[1]{%
  \stepcounter{popsec}\setcounter{popsub}{0}%
  \par\vspace{16pt}\noindent
  \begin{minipage}{\linewidth}
  \tikz[baseline=-0.7ex]\node[draw=popink,line width=1.6pt,fill=poporange,rounded corners=4pt,minimum size=22pt,inner sep=0]{\popdisplay\fontsize{12}{12}\selectfont\color{popcream}\thepopsec};%
  \hspace{8pt}{\popdisplay\fontsize{15}{17}\selectfont\color{popink}\MakeUppercase{#1}}\par
  \vspace{4pt}\noindent{\color{poporange}\rule{36pt}{3.6pt}}
  \end{minipage}\par\nopagebreak\vspace{8pt}
}
\newcounter{popsub}
\newcommand{\courssub}[1]{%
  \stepcounter{popsub}%
  \par\vspace{9pt}\noindent{\popxbold\fontsize{11.5}{13}\selectfont\color{popink}{\color{poporange}\thepopsec.\thepopsub}\hspace{6pt}#1}\par\nopagebreak\vspace{4pt}
}

% ============================================================
% Boîtes pédagogiques — même recette : fond blanc, bordure épaisse colorée,
% ombre dure encre, badge plein. Titre optionnel [#1].
% ============================================================
\tcbset{popbase/.style={breakable,enhanced,colback=white,boxrule=2.3pt,arc=9pt,
  shadow={3pt}{-3pt}{0pt}{popink},left=14pt,right=14pt,top=11pt,bottom=11pt,before skip=11pt,after skip=15pt,
  fontupper=\popsemi\fontsize{10.6}{14.5}\selectfont\color{popink},
  fontlower=\popsemi\fontsize{10.6}{14.5}\selectfont\color{popink}}}
% Coche dessinée (le glyphe ✓ n'existe pas dans les polices du document)
\newcommand{\popcheck}[1]{\tikz[baseline=-0.5ex]\draw[#1,line width=1.6pt,line cap=round,line join=round](-0.13,0.0)--(-0.04,-0.09)--(0.13,0.1);}
\providecommand{\coche}{\popcheck{popgreen}}
\providecommand{\resultat}[1]{\par\smallskip\noindent\fcolorbox{popgreen}{popgreenL}{\parbox{\dimexpr\linewidth-2\fboxsep-2\fboxrule\relax}{\textbf{Résultat.} #1}}\par\smallskip}
\newcommand{\popboxhead}[3]{\poppill{#1}{#2}{white}\ifx\relax#3\relax\else\hspace{6pt}{\popxbold\fontsize{10.6}{12}\selectfont\color{popink}#3}\fi\par\vspace{7pt}}

\newtcolorbox{aretenir}[1][]{popbase,colframe=popgreen,before upper={\popboxhead{À retenir}{popgreen}{#1}}}
% Texte en espagnol (document de lecture, dialogue, poème, extrait) : cadre
% d'étude. Un titre optionnel (source, auteur) s'affiche dans l'en-tête.
\newtcolorbox{textebox}[1][]{popbase,colframe=popblue,before upper={\popboxhead{Texto}{popblue}{#1}\begin{otherlanguage}{spanish}},after upper={\end{otherlanguage}}}
% Marques ✓ / ✗ (forme correcte / fautive) souvent utilisées par les modèles
\providecommand{\cross}{\textcolor{poppink}{\ensuremath{\times}}}
\providecommand{\xmark}{\cross}\providecommand{\crossmark}{\cross}\providecommand{\cmark}{\ensuremath{\checkmark}}
% Ligne de réponse / justification à compléter (inventée par les modèles)
\providecommand{\justification}[1]{\par\noindent\textit{Justification~:} \dotfill\par}
% \textipa (package tipa non chargé) : la police ne couvre pas l'API -> simple passage
\providecommand{\textipa}[1]{#1}
% Soulignement (ulem non chargé) : \uline, \ul
\providecommand{\uline}[1]{\underline{#1}}\providecommand{\ul}[1]{\underline{#1}}
% Mot ou phrase en espagnol à l'intérieur d'un paragraphe français
\newcommand{\esp}[1]{\ifmmode\text{\textit{#1}}\else\begingroup\selectlanguage{spanish}\itshape #1\endgroup\fi}
\newtcolorbox{exemplebox}[1][]{popbase,colframe=poporange,before upper={\popboxhead{Exemple}{poporange}{#1}}}
\newtcolorbox{definition}[1][]{popbase,colframe=popblue,before upper={\popboxhead{Définition}{popblue}{#1}}}
\newtcolorbox{propriete}[1][]{popbase,colframe=poppurple,before upper={\popboxhead{Propriété}{poppurple}{#1}}}
\newtcolorbox{methode}[1][]{popbase,colframe=popgold,before upper={\popboxhead{Méthode}{popgold}{#1}}}
\newtcolorbox{attention}[1][]{popbase,colframe=poppink,before upper={\popboxhead{Attention}{poppink}{#1}}}
\newtcolorbox{experience}[1][]{popbase,colframe=popblue,colback=popblueL,before upper={\popboxhead{Expérience}{popblue}{#1}}}
% « Le savais-tu ? » : carte violette pleine (contexte, culture, Cameroun)
\newtcolorbox{savaistu}[1][]{popbase,colback=poppurple,colframe=popink,
  fontupper=\popsemi\fontsize{10.4}{14.2}\selectfont\color{white},
  before upper={\poppill{Le savais-tu\,?}{white}{poppurple}\ifx\relax#1\relax\else\hspace{6pt}{\popxbold\color{white}#1}\fi\par\vspace{7pt}}}
% Exercice résolu : énoncé en haut, \tcblower puis la solution sur fond crème
\newcounter{popexo}\newcounter{popexr}
\newtcolorbox{exoresolu}[1][]{popbase,colframe=poporange,colbacklower=popcream,
  segmentation style={popink,line width=1.2pt,dashed},
  before upper={\stepcounter{popexr}\popboxhead{Exercice résolu \thepopexr}{poporange}{#1}},
  before lower={\poppill{Solution}{popink}{popcream}\par\vspace{6pt}}}
% Exercice (non résolu en place) — la correction est dans la partie Corrigés
\newtcolorbox{exercice}[1][]{popbase,colframe=popink,
  before upper={\stepcounter{popexo}\popboxhead{Exercice \thepopexo}{popink}{#1}}}
% Corrigé
\newtcolorbox{corrige}[1][]{popbase,colframe=popgreen,colback=popgreenL,
  before upper={\popboxhead{Corrigé}{popgreen}{#1}}}
% Objectifs / prérequis / bilan
\newtcolorbox{objectifs}{popbase,colframe=popink,colback=poporangeL,
  before upper={\popboxhead{Objectifs de la leçon}{popink}{}}}
\newtcolorbox{prerequis}{popbase,colframe=popink,colback=popblueL,
  before upper={\popboxhead{Prérequis}{popblue}{}}}
\newtcolorbox{bilan}{popbase,colframe=popink,colback=white,boxrule=2.8pt,
  before upper={\popboxhead{Fiche bilan}{poporange}{L'essentiel à connaître pour l'examen}}}
% Figure encadrée avec légende
\newtcolorbox{popfigure}[1][]{enhanced,breakable=false,colback=white,colframe=popline,boxrule=1.4pt,arc=8pt,
  left=10pt,right=10pt,top=10pt,bottom=8pt,before skip=10pt,after skip=12pt,halign=center,
  lower separated=false,halign lower=center,
  fontlower=\popsemi\fontsize{9}{11.5}\selectfont\color{popmuted},
  #1}
\newcommand{\legende}[1]{\par\vspace{6pt}{\popsemi\fontsize{9}{11.5}\selectfont\color{popmuted}#1}}

% ============================================================
% Signature OnBuch+
% ============================================================
\newcommand{\poplogoinline}[1]{{\poplogo\fontsize{#1}{#1}\selectfont\color{popink}OnBuch\color{poporange}+}}
\newcommand{\signatureonbuch}{%
  \begin{tcolorbox}[enhanced,colback=popink,colframe=popink,boxrule=2.2pt,arc=10pt,
      drop shadow=poporange,left=14pt,right=14pt,top=10pt,bottom=10pt,before skip=0pt,after skip=0pt]
    \tikz[baseline=-0.7ex]\node[circle,fill=popgreen,draw=popcream,line width=1.3pt,minimum size=22pt,inner sep=0]{\popcheck{white}};%
    \hspace{9pt}\begin{minipage}[c]{0.62\linewidth}
      {\popxbold\fontsize{12}{13}\selectfont\color{popcream}Signé\ \textbullet\ {\poplogo OnBuch\color{poporange}+}}\par\vspace{2pt}
      {\raggedright\popsemi\fontsize{8}{10.5}\selectfont\color{popline}\MakeUppercase{Équipe pédagogique OnBuch+}\\\MakeUppercase{Conforme au programme officiel MINESEC}\par}
    \end{minipage}\hfill
    {\poplogo\fontsize{22}{22}\selectfont\color{popcream}OnBuch\color{poporange}+}
  \end{tcolorbox}%
}

% ============================================================
% Page de garde
% #1 titre · #2 matière · #3 niveau · #4 module · #5 n° leçon · #6 durée ·
% #7 description / objectifs (texte ou itemize)
% ============================================================
\newcommand{\pagegarde}[7]{%
  \thispagestyle{empty}
  \newgeometry{a4paper,margin=1.7cm,top=1.6cm,bottom=1.4cm}%
  \begin{tikzpicture}[remember picture,overlay]
    \fill[poporange!18] ([xshift=-3.2cm,yshift=-3.6cm]current page.north east) circle (4.6cm);
    \fill[poppurple!14] ([xshift=2.4cm,yshift=7.6cm]current page.south west) circle (3.8cm);
  \end{tikzpicture}%
  {\poplogo\fontsize{30}{30}\selectfont\color{popink}OnBuch\color{poporange}+}\hfill
  \raisebox{0.6ex}{\poppill{Cours signé \textbullet\ Premium}{popink}{popcream}}\par
  \vspace{1pt}\popsquiggle{poppurple}\par
  \vspace{8pt}
  \begin{tcolorbox}[enhanced,colback=poporange,colframe=popink,boxrule=2.8pt,arc=13pt,
      drop shadow=popink,left=18pt,right=18pt,top=12pt,bottom=13pt,after skip=10pt]
    \poppill{#2 \textbullet\ #3}{popink}{popcream}\hspace{5pt}\poppill{#4}{white}{popink}\par\vspace{8pt}
    {\popdisplay\fontsize{14}{14}\selectfont\color{popink}LEÇON #5}\par\vspace{4pt}
    {\raggedright\hyphenpenalty=10000\exhyphenpenalty=10000\popdisplay\fontsize{25}{27}\selectfont\color{popcream}\MakeUppercase{#1}\par}
    \vspace{8pt}
    {\popxbold\fontsize{9.5}{9.5}\selectfont\color{popink}\MakeUppercase{Durée conseillée : #6}}
  \end{tcolorbox}
  \begin{tcolorbox}[enhanced,colback=white,colframe=popink,boxrule=2.2pt,arc=10pt,
      drop shadow=popink,left=16pt,right=16pt,top=10pt,bottom=10pt,after skip=8pt]
    \poppill{Dans cette leçon}{poppurple}{white}\par\vspace{5pt}
    \popsemi\fontsize{9.3}{12.6}\selectfont\color{popink}#7
  \end{tcolorbox}
  \noindent\begin{minipage}[t]{0.487\linewidth}
    \begin{tcolorbox}[enhanced,colback=popgreen,colframe=popink,boxrule=2.2pt,arc=10pt,
        drop shadow=popink,left=11pt,right=11pt,top=10pt,bottom=10pt,height=3.1cm,valign=center]
      \tcbox[colback=white,colframe=popink,boxrule=1.3pt,arc=5pt,boxsep=0pt,nobeforeafter,
        left=3pt,right=3pt,top=3pt,bottom=3pt,valign=center]{\qrcode[height=1.9cm]{https://play.google.com/store/apps/details?id=com.luvvix.onbuch&hl=fr}}%
      \hspace{8pt}\begin{minipage}[c]{0.5\linewidth}
        {\popdisplay\fontsize{11}{12.5}\selectfont\color{white}\MakeUppercase{Télécharge OnBuch}}\par\vspace{4pt}
        {\raggedright\popsemi\fontsize{8.4}{11}\selectfont\color{white}Gratuit \textbullet\ Google Play\\Tuteur IA, annales, concours, résultats\par}
      \end{minipage}
    \end{tcolorbox}
  \end{minipage}\hfill
  \begin{minipage}[t]{0.487\linewidth}
    \begin{tcolorbox}[enhanced,colback=poppurple,colframe=popink,boxrule=2.2pt,arc=10pt,
        drop shadow=popink,left=11pt,right=11pt,top=10pt,bottom=10pt,height=3.1cm,valign=center]
      \tikz[baseline=-0.7ex]\node[circle,fill=white,draw=popink,line width=1.3pt,minimum size=1.6cm,inner sep=0]{\popdisplay\fontsize{24}{24}\selectfont\color{poppurple}+};%
      \hspace{8pt}\begin{minipage}[c]{0.6\linewidth}
        {\popdisplay\fontsize{11}{12.5}\selectfont\color{white}\MakeUppercase{Passe à OnBuch+}}\par\vspace{4pt}
        {\raggedright\popsemi\fontsize{8.4}{11}\selectfont\color{white}Tous les cours signés, Tuteur IA illimité, fiches PDF — onglet OnBuch+ de l'app\par}
      \end{minipage}
    \end{tcolorbox}
  \end{minipage}
  \vspace{8pt}
  \popcontenu
  \vfill
  \signatureonbuch
  \restoregeometry
  \newpage
}

% Rangée « Ce cours contient » (page de garde)
\newcommand{\poptile}[3]{% #1 couleur #2 gros texte #3 légende
  \begin{tcolorbox}[enhanced,colback=#1,colframe=popink,boxrule=2pt,arc=9pt,shadow={2.5pt}{-2.5pt}{0pt}{popink},
      left=6pt,right=6pt,top=9pt,bottom=9pt,halign=center,valign=center,height=1.75cm,width=\linewidth,nobeforeafter]
    {\popdisplay\fontsize{12.5}{13}\selectfont\color{white}\MakeUppercase{#2}}\par\vspace{3pt}
    {\centering\popsemi\fontsize{7.8}{9.5}\selectfont\color{white}#3\par}
  \end{tcolorbox}}
\newcommand{\popcontenu}{%
  \par\noindent{\popxboldls\fontsize{7.5}{7.5}\selectfont\color{popmuted}CE COURS SIGNÉ CONTIENT}\par\vspace{7pt}
  \noindent\begin{minipage}[t]{0.235\linewidth}\poptile{popblue}{Cours}{complet, illustré, pas à pas}\end{minipage}\hfill
  \begin{minipage}[t]{0.235\linewidth}\poptile{poporange}{Méthodes}{et exercices résolus}\end{minipage}\hfill
  \begin{minipage}[t]{0.235\linewidth}\poptile{poppink}{Exercices}{type examen + corrigés}\end{minipage}\hfill
  \begin{minipage}[t]{0.235\linewidth}\poptile{popgreen}{Fiche bilan}{l'essentiel à retenir}\end{minipage}\par}

% Sommaire
\newtcolorbox{sommaire}{enhanced,colback=white,colframe=popink,boxrule=2.2pt,arc=10pt,
  drop shadow=popink,left=18pt,right=18pt,top=13pt,bottom=13pt,before skip=6pt,after skip=20pt,
  before upper={\poppill{Sommaire}{poppurple}{white}\par\vspace{9pt}\popsemi\fontsize{10.6}{14.5}\selectfont\color{popink}}}

% Signature de fin de cours — poussée tout en bas de la dernière page
\newcommand{\findecours}{%
  \par\vfill\nopagebreak
  \begin{center}\popsquiggle{poporange}\end{center}\vspace{4pt}
  \signatureonbuch
}
\AtBeginDocument{\def\llbracket{\mathopen{[\mkern-3.5mu[}}\def\rrbracket{\mathclose{]\mkern-3.5mu]}}} % intervalles d'entiers (glyphe ⟦ absent de la police)
% Glyphes absents de Fira Math : redéfinis à partir de glyphes disponibles
\AtBeginDocument{%
  \def\setminus{\mathbin{\backslash}}\let\smallsetminus\setminus
  \def\vdots{\mathord{\vbox{\baselineskip3.2pt\lineskiplimit0pt\kern4pt\hbox{.}\hbox{.}\hbox{.}}}}%
  \def\ddots{\mathinner{\mkern1mu\raise6.5pt\hbox{.}\mkern2mu\raise3.6pt\hbox{.}\mkern2mu\raise0.7pt\hbox{.}\mkern1mu}}%
  \def\triangle{\mathord{\scalebox{0.9}{$\symup{\Delta}$}}}%
  \def\nabla{\mathord{\raisebox{\depth}{\scalebox{1}[-1]{$\symup{\Delta}$}}}}%
  \def\checkmark{\popcheck{popdark}}%
  \def\star{\mathbin{\ast}}\def\bigstar{\mathord{\ast}}%
  \def\diamond{\mathbin{\scalebox{0.6}{$\Diamond$}}}%
  \def\bigcup{\mathop{\mathchoice{\raisebox{-0.3em}{\scalebox{1.7}{$\cup$}}}{\scalebox{1.25}{$\cup$}}{\cup}{\cup}}\displaylimits}%
  \def\bigcap{\mathop{\mathchoice{\raisebox{-0.3em}{\scalebox{1.7}{$\cap$}}}{\scalebox{1.25}{$\cap$}}{\cap}{\cap}}\displaylimits}%
  \def\lesseqqgtr{\mathrel{\lessgtr}}\def\bigoplus{\mathop{\oplus}\displaylimits}%
}
\providecommand{\ssi}{si et seulement si} % abréviation fréquente
\AtBeginDocument{\providecommand{\wideparen}{\overparen}} % arc de cercle

%% FIGFIX-BEGIN v3 — correctifs globaux des figures (docs/onbuch-generation/tools/figfix)
\usepackage{adjustbox}\usepackage{etoolbox}
% toute figure TikZ/pgfplots plus large que la ligne est réduite à la largeur disponible
\BeforeBeginEnvironment{tikzpicture}{\begin{adjustbox}{max width=\linewidth}}
\AfterEndEnvironment{tikzpicture}{\end{adjustbox}}
% légendes pgfplots lisibles par-dessus les courbes
\pgfplotsset{every axis/.append style={legend style={fill=white,fill opacity=0.92,text opacity=1,draw=black!15,inner sep=3pt}}}
% étiquettes d'arêtes (syntaxe edge["..."]) avec fond blanc
\tikzset{every edge quotes/.append style={fill=white,inner sep=1.2pt}}
% bulles de cartes mentales : texte à la ligne dans la bulle, bulle un peu plus grande
\tikzset{every concept/.append style={text width=2.0cm,align=center,minimum size=2.3cm,inner sep=2pt}}
%% FIGFIX-END
\begin{document}

\pagegarde{Lectura y comentario : oficios, mundo del trabajo, derechos y deberes de los trabajadores}{Espagnol}{Tle A}{Módulo 4 : Activités économiques et monde du travail}{E20}{2 h}{Cette leçon de lecture et commentaire propose un texte original sur les métiers et le monde du travail, enrichi du lexique essentiel (empleo, paro, contrato, sindicato, huelga, jubilación). Elle entraîne l'élève à comprendre, commenter et réutiliser ce vocabulaire pour présenter les droits et devoirs des travailleurs.
\vspace{4pt}
\begin{multicols}{2}\small
\begin{itemize}
\item À la fin de cette leçon, je sais lire et comprendre globalement un texte espagnol sur les métiers et le monde du travail.
\item Je sais dégager l'idée principale et les idées secondaires d'un texte.
\item Je sais employer correctement le lexique du travail : oficio, profesión, empleo, paro, sueldo, contrato, sindicato, huelga, jubilación.
\item Je sais distinguer oficio, profesión et empleo et les utiliser à bon escient.
\item Je sais présenter les droits et les devoirs des travailleurs en espagnol.
\item Je sais répondre à des questions de compréhension en phrases complètes.
\end{itemize}\end{multicols}}

\begin{sommaire}
\begin{multicols}{2}
\begin{enumerate}
\item Texte support : « El mundo del trabajo hoy »
\item Compréhension du texte
\item Lexique thématique : el mundo del trabajo
\item Droits et devoirs des travailleurs
\item Méthode du commentaire de texte
\item Production guidée
\item Activité d'intégration
\item Méthodes et exercices résolus
\item Exercices
\item Corrigés des exercices
\item Fiche bilan
\end{enumerate}
\end{multicols}
\end{sommaire}

\begin{objectifs}
\begin{itemize}
\item À la fin de cette leçon, je sais lire et comprendre globalement un texte espagnol sur les métiers et le monde du travail.
\item Je sais dégager l'idée principale et les idées secondaires d'un texte.
\item Je sais employer correctement le lexique du travail : oficio, profesión, empleo, paro, sueldo, contrato, sindicato, huelga, jubilación.
\item Je sais distinguer oficio, profesión et empleo et les utiliser à bon escient.
\item Je sais présenter les droits et les devoirs des travailleurs en espagnol.
\item Je sais répondre à des questions de compréhension en phrases complètes.
\item Je sais rédiger un commentaire structuré (introduction, développement, conclusion) sur un texte.
\item Je sais produire un court texte personnel sur mon projet professionnel.
\end{itemize}
\end{objectifs}

\begin{prerequis}
\begin{itemize}
\item Le présent de l'indicatif et les verbes fréquents (trabajar, ganar, tener, deber, poder).
\item Le vocabulaire général de la vie quotidienne et des études (Première).
\item Les connecteurs simples : pero, porque, además, sin embargo, por eso.
\item La structure d'un paragraphe argumentatif vue en Première.
\item Les pronoms personnels et les adjectifs possessifs.
\end{itemize}
\end{prerequis}

\newpage

\coursec{Texte support : « El mundo del trabajo hoy »}

À Yaoundé, comme à Madrid ou à Mexico, les jeunes terminent leurs études avec une question en tête : \esp{¿qué oficio voy a ejercer?} Entre le menuisier du quartier Briqueterie, l'infirmière de l'hôpital central et l'informaticien \emph{freelance} qui travaille depuis un cybercafé de Douala, le monde du travail est partout autour de nous. Mais travailler, ce n'est pas seulement gagner un \esp{sueldo} : c'est aussi connaître ses \cle{droits} (\esp{contrato}, \esp{sindicato}, \esp{jubilación}) et ses \cle{devoirs}. Au baccalauréat, les textes sur le travail et les métiers sont fréquents : il faut savoir les lire, les comprendre et les commenter.

\textbf{Problématique :} comment lire efficacement un texte espagnol sur le monde du travail, deviner le sens des mots nouveaux grâce au contexte, et identifier le thème, le type de texte et le point de vue du narrateur ?

\courssub{Situation de lecture : deux portraits de travailleurs}

Avant de lire le texte, observons cette phrase entendue dans une cour de lycée à Douala :

\esp{Mi tío es carpintero en Douala. Tiene un oficio que aprendió con su padre. Trabaja mucho, pero no tiene contrato escrito ni sindicato que lo defienda. Mi prima, en cambio, es enfermera en Madrid : firmó un contrato, cobra un sueldo cada mes y piensa en su jubilación. Los dos trabajan con dignidad, pero no tienen los mismos derechos.}

\begin{textebox}[Phrase d'accroche : deux situations de travail]
Mi tío es carpintero en Douala. Tiene un oficio que aprendió con su padre. Trabaja mucho, pero no tiene contrato escrito ni sindicato que lo defienda. Mi prima, en cambio, es enfermera en Madrid : firmó un contrato, cobra un sueldo cada mes y piensa en su jubilación. Los dos trabajan con dignidad, pero no tienen los mismos derechos.
\end{textebox}

\emph{Traduction :} « Mon oncle est menuisier à Douala. Il a un métier qu'il a appris avec son père. Il travaille beaucoup, mais il n'a ni contrat écrit ni syndicat pour le défendre. Ma cousine, en revanche, est infirmière à Madrid : elle a signé un contrat, elle touche un salaire chaque mois et elle pense à sa retraite. Tous les deux travaillent avec dignité, mais ils n'ont pas les mêmes droits. »

Cette phrase annonce le thème du texte : le contraste entre deux manières de travailler, l'une artisanale et informelle, l'autre salariée et protégée. Retenons trois mots clés : \esp{oficio} (métier), \esp{contrato} (contrat), \esp{derechos} (droits).

\courssub{Méthode : comment lire un texte en espagnol}

\begin{methode}[Les cinq étapes de la lecture efficace]
\begin{enumerate}
\item \textbf{Lire le titre et la source} : ils annoncent le thème. Ici, \esp{El mundo del trabajo hoy} annonce un texte sur le travail aujourd'hui.
\item \textbf{Première lecture globale, sans dictionnaire} : lire le texte en entier, en silence, pour saisir l'idée générale. Ne pas s'arrêter sur les mots inconnus.
\item \textbf{Repérer les repères : ¿quién? ¿qué? ¿dónde? ¿cuándo?} (qui ? quoi ? où ? quand ?) : identifier les personnages, les lieux, la situation.
\item \textbf{Deuxième lecture avec le glossaire} : deviner d'abord le sens des mots nouveaux grâce au contexte, puis vérifier dans le glossaire.
\item \textbf{Reformuler} : résumer le texte à voix haute en 3 ou 4 phrases, avec ses propres mots.
\end{enumerate}
\end{methode}

\courssub{Première lecture : découverte du texte}

Lisons le texte en silence, puis à voix haute, en respectant la prononciation et l'intonation. Essayons de deviner le sens des mots soulignés par le contexte avant de consulter le glossaire.

\begin{textebox}[El mundo del trabajo hoy]
Mi tío Samuel es carpintero en Douala, la gran ciudad económica de Camerún. Tiene un taller pequeño cerca del mercado, donde fabrica mesas, sillas y armarios. Aprendió su oficio con su padre, que también era carpintero : no fue a la escuela de oficios, sino que aprendió trabajando, día tras día. Samuel se levanta a las seis de la mañana y trabaja hasta la noche. Sin embargo, no tiene contrato escrito : sus clientes le pagan por cada mueble, y si un día no hay trabajo, no hay sueldo. Tampoco pertenece a ningún sindicato que defienda sus derechos. « El paro es mi mayor miedo », dice. « Si me pongo enfermo, ¿quién pagará a mi familia ? »

Mi prima Lucía, en cambio, vive en Madrid y es enfermera en un hospital público. Para ejercer su profesión, estudió cuatro años en la universidad y aprobó un examen oficial. Firmó un contrato de trabajo con el hospital : trabaja cuarenta horas por semana y cobra un sueldo fijo cada mes. Está afiliada a un sindicato que defiende a los trabajadores de la salud. El año pasado, cuando el hospital quiso reducir el personal, el sindicato convocó una huelga y los enfermeros consiguieron mantener sus empleos. Lucía piensa también en el futuro : cada mes, una parte de su sueldo se guarda para su jubilación.

Los dos, Samuel y Lucía, trabajan con dignidad y amor por su oficio. Pero no tienen los mismos derechos ni la misma seguridad. El mundo del trabajo hoy es así : lleno de contrastes entre quienes tienen un contrato y quienes trabajan sin protección.
\end{textebox}

\emph{Traduction du texte :} « Mon oncle Samuel est menuisier à Douala, la grande ville économique du Cameroun. Il a un petit atelier près du marché, où il fabrique des tables, des chaises et des armoires. Il a appris son métier avec son père, qui était menuisier lui aussi : il n'est pas allé à l'école des métiers, mais il a appris en travaillant, jour après jour. Samuel se lève à six heures du matin et travaille jusqu'au soir. Cependant, il n'a pas de contrat écrit : ses clients le paient pour chaque meuble, et si un jour il n'y a pas de travail, il n'y a pas de salaire. Il n'appartient non plus à aucun syndicat qui défende ses droits. "Le chômage est ma plus grande peur", dit-il. "Si je tombe malade, qui paiera pour ma famille ?"

Ma cousine Lucía, en revanche, vit à Madrid et elle est infirmière dans un hôpital public. Pour exercer sa profession, elle a étudié quatre ans à l'université et a réussi un examen officiel. Elle a signé un contrat de travail avec l'hôpital : elle travaille quarante heures par semaine et touche un salaire fixe chaque mois. Elle est affiliée à un syndicat qui défend les travailleurs de la santé. L'année dernière, quand l'hôpital a voulu réduire le personnel, le syndicat a appelé à la grève et les infirmiers ont réussi à conserver leurs emplois. Lucía pense aussi à l'avenir : chaque mois, une partie de son salaire est mise de côté pour sa retraite.

Tous les deux, Samuel et Lucía, travaillent avec dignité et amour de leur métier. Mais ils n'ont pas les mêmes droits ni la même sécurité. Le monde du travail aujourd'hui est ainsi : plein de contrastes entre ceux qui ont un contrat et ceux qui travaillent sans protection. »

\courssub{Repérage des mots nouveaux grâce au contexte}

Avant de consulter le glossaire, devinons le sens des mots nouveaux. C'est une compétence essentielle au baccalauréat : le contexte donne presque toujours des indices.

\begin{exemplebox}[Deviner le sens grâce au contexte]
\begin{itemize}
\item \esp{Tiene un taller pequeño cerca del mercado, donde fabrica mesas, sillas y armarios.} --- Un lieu où l'on fabrique des meubles : \esp{taller} = atelier.
\item \esp{Si un día no hay trabajo, no hay sueldo.} --- Ce qu'on reçoit quand on travaille : \esp{sueldo} = salaire.
\item \esp{El paro es mi mayor miedo. Si me pongo enfermo...} --- La peur de celui qui n'a pas de travail stable : \esp{paro} = chômage.
\item \esp{El sindicato convocó una huelga.} --- Ce qu'un syndicat organise pour protester : \esp{huelga} = grève.
\item \esp{Una parte de su sueldo se guarda para su jubilación.} --- Ce qu'on prépare pour la fin de la vie active : \esp{jubilación} = retraite.
\end{itemize}
\end{exemplebox}

\courssub{Glossaire bilingue des mots difficiles}

\begin{definition}[Lexique du monde du travail]
\begin{itemize}
\item \esp{el oficio} : métier (surtout manuel ou artisanal, appris par la pratique).
\item \esp{la profesión} : profession (métier exigeant des études ou un diplôme).
\item \esp{el empleo} : emploi, poste de travail.
\item \esp{el paro} : chômage (en Espagne ; en Amérique latine : \esp{el desempleo}).
\item \esp{el sueldo} : salaire (synonyme : \esp{el salario}).
\item \esp{el contrato} : contrat (\esp{firmar un contrato} = signer un contrat).
\item \esp{el sindicato} : syndicat (\esp{estar afiliado a un sindicato} = être affilié à un syndicat).
\item \esp{la huelga} : grève (\esp{convocar una huelga} = appeler à la grève ; \esp{hacer huelga} = faire grève).
\item \esp{la jubilación} : retraite (\esp{jubilarse} = prendre sa retraite).
\item \esp{los derechos} : droits ; \esp{los deberes} : devoirs.
\item \esp{el empresario / la empresaria} : employeur, patron, chef d'entreprise.
\item \esp{el taller} : atelier ; \esp{el cliente} : client ; \esp{fabricar} : fabriquer.
\end{itemize}
\end{definition}

\begin{attention}[Faux amis et pièges de traduction]
\begin{itemize}
\item \esp{El paro} signifie « le chômage » en Espagne, jamais « l'arrêt ». Pour dire « arrêt », on utilise \esp{la parada} (arrêt de bus) ou \esp{la detención}. En Amérique latine, on préfère \esp{el desempleo}. \esp{Estar en paro} = être au chômage.
\item \esp{El oficio} n'est pas « l'office » : c'est le métier, surtout artisanal. \esp{El carpintero aprendió su oficio con su padre.} « Le menuisier a appris son métier avec son père. »
\item \esp{El empresario} est le patron ou l'employeur, pas « l'entrepreneur » au sens de créateur de start-up uniquement.
\item \esp{Jubilarse} = prendre sa retraite ; ne pas confondre avec le français « jubiler » !
\end{itemize}
\end{attention}

\courssub{Tableau récapitulatif : oficio, profesión, empleo}

Ces trois mots désignent le « travail », mais ils ne sont pas interchangeables. Observons le texte : Samuel a un \esp{oficio} (appris avec son père, sans diplôme), Lucía a une \esp{profesión} (quatre ans d'université, examen officiel), et les infirmiers de Madrid ont défendu leurs \esp{empleos} (leurs postes concrets à l'hôpital).

\begin{center}
\begin{tabularx}{\linewidth}{|X|X|X|}
\hline
\rowcolor{popblueL} Español & Français & Remarque \\
\hline
\esp{el oficio} & le métier (manuel, artisanal) & appris par la pratique : \esp{carpintero, zapatero, sastre} \\
\hline
\esp{la profesión} & la profession & exige des études ou un diplôme : \esp{enfermera, médico, abogado} \\
\hline
\esp{el empleo} & l'emploi, le poste & le poste concret occupé : \esp{perder el empleo} = perdre son emploi \\
\hline
\esp{el trabajo} & le travail & terme général : \esp{buscar trabajo} = chercher du travail \\
\hline
\esp{el sueldo / el salario} & le salaire & \esp{cobrar un sueldo} = toucher un salaire \\
\hline
\esp{la huelga} & la grève & \esp{hacer huelga} = faire grève \\
\hline
\esp{la jubilación} & la retraite & \esp{jubilarse} = prendre sa retraite \\
\hline
\esp{el paro / el desempleo} & le chômage & \esp{paro} en Espagne, \esp{desempleo} en Amérique latine \\
\hline
\end{tabularx}
\end{center}

\begin{definition}[Oficio, profesión, empleo : la nuance essentielle]
\esp{Oficio} désigne un métier manuel ou artisanal appris par la pratique, souvent transmis de père en fils ; \esp{profesión} désigne un métier exigeant des études ou un diplôme officiel ; \esp{empleo} désigne le poste de travail concret occupé par une personne. Exemple : \esp{Lucía estudió para la profesión de enfermera y consiguió un empleo en un hospital.} « Lucía a étudié pour la profession d'infirmière et a obtenu un emploi dans un hôpital. »
\end{definition}

\courssub{Droits et devoirs des travailleurs : synthèse}

Le texte oppose deux statuts : l'artisan informel (Samuel) et la salariée protégée (Lucía). Le programme exige de connaître les droits et devoirs fondamentaux.

\begin{aretenir}[Droits et devoirs fondamentaux du travailleur]
\begin{itemize}
\item \textbf{Derechos (Droits) :}
      \begin{itemize}
      \item \esp{Derecho a un contrato de trabajo} (écrit, clair, précisant horaires, salaire, tâches).
      \item \esp{Derecho a un sueldo / salario justo} (payé régulièrement, \esp{por transferencia} ou \esp{en efectivo} avec reçu).
      \item \esp{Derecho a la seguridad social y a la jubilación} (cotisations maladie, maternité, vieillesse).
      \item \esp{Derecho a la sindicación} (créer ou rejoindre un \esp{sindicato}, droit de grève \esp{huelga}).
      \item \esp{Derecho al descanso} (heures sup payées, congés payés \esp{vacaciones}, jours fériés).
      \item \esp{Derecho a la no discriminación} (sexe, origine, âge, handicap).
      \end{itemize}
\item \textbf{Deberes (Devoirs) :}
      \begin{itemize}
      \item \esp{Deber de cumplir las tareas} propres à son poste avec diligence.
      \item \esp{Deber de respetar horarios} et calendrier de l'entreprise.
      \item \esp{Deber de obedecer órdenes} légitimes de l'employeur (\esp{empresario}).
      \item \esp{Deber de guardar secreto} professionnel / confidentialité.
      \item \esp{Deber de no concurrir} (clause de non-concurrence souvent présente).
      \item \esp{Deber de cuidar los bienes} et outils de travail.
      \end{itemize}
\end{itemize}
\end{aretenir}

\begin{attention}[Contraste Samuel / Lucía : ce que le texte illustre]
\begin{center}
\begin{tabularx}{\linewidth}{|X|X|X|}
\hline
\rowcolor{popblueL} Aspecto & Samuel (Oficio informal) & Lucía (Profesión con contrato) \\
\hline
Contrato & \esp{No tiene contrato escrito} & \esp{Contrato de trabajo firmado} \\
\hline
Sueldo & \esp{Por pieza, variable, sin garantía} & \esp{Sueldo fijo mensual} \\
\hline
Protección social & \esp{Ninguna (ni enfermedad, ni jubilación)} & \esp{Seguridad social, jubilación} \\
\hline
Defensa colectiva & \esp{Sin sindicato} & \esp{Afiliada a sindicato, huelga efectiva} \\
\hline
Miedo principal & \esp{El paro, la enfermedad} & \esp{Reducción de plantilla (gestionada)} \\
\hline
\end{tabularx}
\end{center}
\end{attention}

\courssub{Carte mentale du vocabulaire}

\begin{popfigure}
\begin{tikzpicture}[mindmap, grow cyclic, every node/.style={concept}, concept color=popblue!60, text=white,
level 1/.append style={level distance=3.2cm, sibling angle=90},
level 2/.append style={level distance=2.2cm, sibling angle=45}]
\node[concept color=popblue!80, font=\large\bfseries]{El mundo del trabajo}
  child[concept color=popgreen!70]{ node {Tipos de trabajo}
    child{ node {Oficios} child { node {carpintero} } child { node {zapatero} } }
    child{ node {Profesiones} child { node {enfermera} } child { node {médico} } }
  }
  child[concept color=poporange!80]{ node {Derechos}
    child{ node {contrato} }
    child{ node {sueldo} }
    child{ node {jubilación} }
    child{ node {sindicato} }
  }
  child[concept color=poppurple!70]{ node {Deberes}
    child{ node {puntualidad} }
    child{ node {cumplir tareas} }
    child{ node {respetar normas} }
  }
  child[concept color=poppink!80]{ node {Conflictos}
    child{ node {paro} }
    child{ node {huelga} }
  };
\end{tikzpicture}
\legende{Carte mentale : le vocabulaire du monde du travail structuré en quatre branches (types de travail, droits, devoirs, conflits).}
\end{popfigure}

\courssub{Identification : thème, type de texte, point de vue}

Après la deuxième lecture, répondons aux trois questions fondamentales du commentaire.

\begin{aretenir}[Fiche d'identification du texte]
\begin{itemize}
\item \textbf{Thème :} le monde du travail aujourd'hui, avec le contraste entre le travail artisanal informel (Samuel, à Douala) et le travail salarié protégé (Lucía, à Madrid).
\item \textbf{Type de texte :} un récit-témoignage (\esp{relato-testimonio}) : le narrateur raconte la vie de deux membres de sa famille pour illustrer une réalité sociale.
\item \textbf{Point de vue :} narration à la première personne (\esp{mi tío}, \esp{mi prima}) : le narrateur est témoin, ce qui rend le récit personnel et crédible. Les paroles de Samuel sont rapportées au discours direct (\esp{« El paro es mi mayor miedo »}), ce qui donne la parole au travailleur lui-même.
\item \textbf{Idée principale :} \esp{Los dos trabajan con dignidad, pero no tienen los mismos derechos.} La dignité du travail est universelle, mais la protection des travailleurs varie selon les pays et les statuts.
\end{itemize}
\end{aretenir}

\courssub{Méthode : commentaire de texte écrit (épreuve du Bac)}

Le commentaire de texte suit une structure rigoureuse en trois parties. C'est l'exercice type de l'écrit (compréhension/expression) et de l'oral.

\begin{methode}[Structure du commentaire de texte]
\begin{enumerate}
\item \textbf{Introduction (1 paragraphe) :}
      \begin{itemize}
      \item \emph{Accroche} : phrase d'ouverture sur le thème général (ex: \esp{El mundo laboral presenta grandes desigualdades...}).
      \item \emph{Présentation du texte} : titre (entre guillemets), auteur (ici : anonyme / manuel), nature (récit-témoignage), source si connue.
      \item \emph{Problématique} : question à laquelle le texte répond (ex: \esp{¿Hasta qué punto el tipo de contrato determina la seguridad del trabajador?}).
      \item \emph{Annonce du plan} : \esp{En primer lugar...; En segundo lugar...; Finalmente...}
      \end{itemize}
\item \textbf{Développement (2 ou 3 paragraphes) :} Chaque paragraphe = une idée force appuyée sur le texte.
      \begin{itemize}
      \item \emph{Idée 1 :} Deux modèles de travail opposés (informel vs formel). Citer : \esp{oficio} vs \esp{profesión}, \esp{taller} vs \esp{hospital}.
      \item \emph{Idée 2 :} L'impact sur les droits et la sécurité (contrat, sueldo, sindicato, jubilación). Citer : \esp{no tiene contrato} / \esp{firmó un contrato}.
      \item \emph{Idée 3 (si 3 parties) :} La dignité commune mais l'inégalité structurelle. Citer : \esp{trabajan con dignidad... pero no tienen los mismos derechos}.
      \end{itemize}
\item \textbf{Conclusion (1 paragraphe) :}
      \begin{itemize}
      \item \emph{Bilan} : réponse synthétique à la problématique.
      \item \emph{Ouverture} : lien avec l'actualité (Cameroun, Guinée équatoriale, Espagne), perspective personnelle ou citation.
      \end{itemize}
\end{enumerate}
\end{methode}

\courssub{Exercice résolu : reformulation orale}

\begin{exoresolu}[Reformuler le texte en 3-4 phrases]
Énoncé : reformule le texte \esp{El mundo del trabajo hoy} en trois ou quatre phrases, avec tes propres mots, sans recopier le texte.
\tcblower
Solution détaillée : une bonne reformulation reprend l'idée principale, les deux personnages, le contraste et la conclusion, avec un vocabulaire personnel. Modèle de réponse :

\esp{El texto presenta a dos trabajadores muy diferentes : Samuel, un carpintero de Douala que trabaja sin contrato ni protección, y Lucía, una enfermera de Madrid que tiene un empleo seguro con sueldo fijo y sindicato. Aunque los dos aman su trabajo, sus condiciones de vida son muy distintas : Samuel teme el paro y la enfermedad, mientras que Lucía puede preparar su jubilación con tranquilidad. El texto muestra así los contrastes del mundo del trabajo actual y la importancia de los derechos de los trabajadores.}

« Le texte présente deux travailleurs très différents : Samuel, un menuisier de Douala qui travaille sans contrat ni protection, et Lucía, une infirmière de Madrid qui a un emploi sûr avec un salaire fixe et un syndicat. Bien que tous les deux aiment leur travail, leurs conditions de vie sont très différentes : Samuel craint le chômage et la maladie, tandis que Lucía peut préparer sa retraite tranquillement. Le texte montre ainsi les contrastes du monde du travail actuel et l'importance des droits des travailleurs. »
\end{exoresolu}

\courssub{Exercices de compréhension et d'expression}

\begin{exercice}[Compréhension détaillée]
Réponds en espagnol, avec des phrases complètes :
\begin{enumerate}
\item ¿Qué oficio tiene Samuel y dónde trabaja ?
\item ¿Por qué dice Samuel que el paro es su mayor miedo ?
\item ¿Qué derechos tiene Lucía que Samuel no tiene ?
\item ¿Qué hizo el sindicato cuando el hospital quiso reducir el personal ?
\item ¿Y tú? ¿Qué oficio o profesión quieres ejercer en el futuro ?
\end{enumerate}
\end{exercice}

\begin{corrige}[Exercice : compréhension du texte]
\begin{enumerate}
\item \esp{Samuel es carpintero y trabaja en un taller pequeño cerca del mercado de Douala.} (Samuel est menuisier et travaille dans un petit atelier près du marché de Douala.)
\item \esp{Porque no tiene contrato escrito: si un día no hay trabajo, no hay sueldo, y si se pone enfermo, nadie pagará a su familia.} (Parce qu'il n'a pas de contrat écrit : si un jour il n'y a pas de travail, il n'y a pas de salaire, et s'il tombe malade, personne ne paiera pour sa famille.)
\item \esp{Lucía tiene un contrato de trabajo, un sueldo fijo cada mes, un sindicato que la defiende y una jubilación para el futuro.} (Lucía a un contrat de travail, un salaire fixe chaque mois, un syndicat qui la défend et une retraite pour l'avenir.)
\item \esp{El sindicato convocó una huelga y los enfermeros consiguieron mantener sus empleos.} (Le syndicat a appelé à la grève et les infirmiers ont réussi à conserver leurs emplois.)
\item Réponse libre, par exemple : \esp{Quiero ser médico porque me gusta ayudar a la gente.} (Je veux devenir médecin parce que j'aime aider les gens.)
\end{enumerate}
\end{corrige}

\begin{exercice}[Production écrite guidée : article d'opinion]
Rédige un article court (80-100 mots) pour le journal de ton lycée sur le thème : \esp{« Trabajo y derechos : ¿qué esperamos los jóvenes? »}.
Utilise le vocabulaire de la leçon (\esp{contrato, sueldo, sindicato, jubilación, paro, oficio, profesión}) et structure ton texte :
\begin{enumerate}
\item Titre accrocheur.
\item Introduction : présentation du thème (chômage des jeunes, précarité).
\item Développement : ce que vous exigez (contrat, salaire décent, protection sociale, syndicat).
\item Conclusion : appel à l'action ou espoir.
\end{enumerate}
\end{exercice}

\begin{corrige}[Exercice : production écrite guidée (proposition de corrigé)]
\textbf{Título: \esp{Juventud y trabajo: exigimos un futuro con derechos}}

\esp{El paro juvenil es una realidad en Camerún y en todo el mundo. Muchos terminamos la carrera o el aprendizaje de un oficio y nos encontramos sin empleo ni contrato. No basta con tener ganas de trabajar; necesitamos garantías.}

\esp{Exigimos un contrato escrito que asegure un sueldo justo y cotización para la jubilación. Queremos poder afiliarnos a un sindicato que nos defienda ante el empresario. La huelga no es un capricho, es un derecho cuando no nos escuchan.}

\esp{Trabajar con dignidad, como Samuel y Lucía, debe significar tener los mismos derechos, estemos en Douala o en Madrid. ¡El futuro se construye con trabajo protegido!}
\end{corrige}

\begin{savaistu}[Le droit de grève en Espagne et au Cameroun]
En Espagne, la \esp{huelga} est un droit constitutionnel : l'article 28 de la Constitution de 1978 reconnaît le droit de grève (\esp{el derecho a la huelga}) à tous les travailleurs. La journée de travail légale y est de 40 heures par semaine, comme le montre le contrat de Lucía dans notre texte. La \esp{huelga general} (grève générale) est un moment fort de la vie sociale espagnole : les syndicats les plus connus sont \esp{UGT} (\esp{Unión General de Trabajadores}) et \esp{CCOO} (\esp{Comisiones Obreras}).

Au Cameroun, le droit de grève existe aussi, garanti par la Constitution de 1996 (préambule), même si son exercice est encadré par la loi n°92/007 du 14 août 1992 (code du travail). Les centrales syndicales principales sont la \esp{CGT} (Confédération Générale du Travail) et l'\esp{USLC} (Union des Syndicats Libres du Cameroun).

En Guinée équatoriale, seul pays hispanophone d'Afrique, l'espagnol est langue officielle et le vocabulaire du travail y est le même qu'en Espagne (\esp{contrato, nómina, seguridad social, jubilación}) !
\end{savaistu}

\begin{aretenir}[L'essentiel de cette leçon]
\begin{itemize}
\item \textbf{Lecture en 5 étapes :} titre/source, lecture globale, repères (qui/quoi/où/quand), lecture + glossaire, reformulation.
\item \textbf{Deviner le sens par le contexte} : compétence clé pour le Bac (gain de temps, autonomie).
\item \textbf{Trois nuances de « travail » :} \esp{oficio} (artisanal, pratique), \esp{profesión} (diplômée, études), \esp{empleo} (poste concret).
\item \textbf{Droits clés :} contrat écrit, salaire fixe, sécurité sociale, syndicat, grève, retraite.
\item \textbf{Devoirs clés :} diligence, horaires, obéissance légitime, confidentialité, soin du matériel.
\item \textbf{Commentaire de texte :} Introduction (accroche, présentation, problématique, plan), Développement (arguments + citations), Conclusion (bilan + ouverture).
\item \textbf{Attention faux amis :} \esp{paro} = chômage (pas « arrêt ») ; \esp{oficio} = métier (pas « office ») ; \esp{jubilarse} = prendre sa retraite (pas « jubiler ») ; \esp{empresario} = employeur/patron.
\end{itemize}
\end{aretenir}

\coursec{Compréhension du texte}

Dans la section précédente, nous avons lu et observé le texte \esp{« El mundo del trabajo hoy »}, qui présente deux personnages aux réalités professionnelles très différentes : un oncle, ouvrier du bâtiment à Douala, et une cousine, infirmière à Madrid. Lire un texte ne suffit pas : au baccalauréat comme dans la vie, il faut \cle{comprender}, c'est-à-dire saisir le sens global, vérifier les détails, interpréter les intentions de l'auteur et réagir personnellement. Cette section vous donne la méthode complète pour répondre aux questions de compréhension, avec des exemples tirés du texte support.

\courssub{Les quatre étapes de la compréhension}

\begin{popfigure}
\begin{tikzpicture}[
  node distance=0.7cm,
  etape/.style={rectangle, rounded corners, draw=popblue, fill=popblueL, thick, align=center, minimum width=4.5cm, minimum height=1.2cm, font=\small},
  flecha/.style={-{Stealth[length=3mm]}, thick, poporange}
]
\node[etape, align=center] (e1){\textbf{1. Lectura global}\\ ¿De qué habla el texto?};
\node[etape, below=of e1, align=center] (e2){\textbf{2. Lectura detallada}\\ ¿Quién? ¿Dónde? ¿Cuándo?};
\node[etape, below=of e2, align=center] (e3){\textbf{3. Interpretación}\\ ¿Por qué dice el autor...?};
\node[etape, below=of e3, align=center] (e4){\textbf{4. Reacción personal}\\ ¿Y tú, qué opinas?};
\draw[flecha] (e1) -- (e2);
\draw[flecha] (e2) -- (e3);
\draw[flecha] (e3) -- (e4);
\end{tikzpicture}
\legende{Organigramme de la démarche de compréhension : de la lecture globale à la réaction personnelle.}
\end{popfigure}

\begin{propriete}[La démarche en quatre temps]
Pour comprendre un texte en espagnol, on procède toujours dans le même ordre :
\begin{enumerate}
\item \textbf{Lecture globale} : on lit une première fois pour identifier le \cle{tema} (le sujet), le type de texte et les personnages.
\item \textbf{Lecture détaillée} : on relit en soulignant les informations précises : lieux, métiers, chiffres, différences.
\item \textbf{Interprétation} : on cherche ce que l'auteur \emph{suggère} sans le dire directement.
\item \textbf{Réaction personnelle} : on donne son avis, en lien avec le texte.
\end{enumerate}
\end{propriete}

\courssub{Questions de compréhension globale}

Les questions globales portent sur l'ensemble du texte. Les plus fréquentes sont :

\begin{itemize}
\item \esp{¿De qué habla el texto?} — De quoi parle le texte ?
\item \esp{¿Quiénes son los personajes?} — Qui sont les personnages ?
\item \esp{¿Dónde y cuándo ocurre la historia?} — Où et quand se déroule l'histoire ?
\end{itemize}

\begin{exemplebox}[Réponses globales sur notre texte]
\esp{El texto habla del mundo del trabajo: presenta a dos trabajadores, un obrero de la construcción en Douala y una enfermera en Madrid.}

« Le texte parle du monde du travail : il présente deux travailleurs, un ouvrier du bâtiment à Douala et une infirmière à Madrid. »

\esp{Los personajes son el tío del narrador, que es albañil, y su prima, que es enfermera.}

« Les personnages sont l'oncle du narrateur, qui est maçon, et sa cousine, qui est infirmière. »
\end{exemplebox}

\courssub{Questions de compréhension détaillée}

Les questions détaillées exigent de retrouver une information précise dans le texte. La règle d'or : \cle{reprendre les mots de la question dans la réponse}.

\begin{methode}[Répondre par une phrase complète]
\begin{enumerate}
\item Lis la question et souligne les mots-clés : \esp{¿Dónde trabaja la prima?}
\item Reprends ces mots pour commencer ta réponse : \esp{La prima trabaja...}
\item Complète avec l'information du texte : \esp{La prima trabaja en Madrid, en un hospital.}
\item Vérifie : ta réponse est une phrase complète, avec sujet et verbe conjugué.
\end{enumerate}
\end{methode}

\begin{center}
\begin{tabularx}{\linewidth}{|X|X|}
\hline
\rowcolor{popblueL} \textbf{Pregunta} & \textbf{Respuesta completa} \\
\hline
\esp{¿Qué oficio tiene el tío?} & \esp{El tío es albañil: trabaja en la construcción.} \\
\hline
\esp{¿Dónde trabaja la prima?} & \esp{La prima trabaja en Madrid, como enfermera en un hospital.} \\
\hline
\esp{¿Qué diferencias hay entre los dos trabajadores?} & \esp{El tío no tiene contrato escrito y cobra por día; la prima tiene contrato, cobra un sueldo cada mes y piensa en su jubilación.} \\
\hline
\end{tabularx}
\end{center}

\begin{exemplebox}[Phrases clés du texte, traduites]
\esp{No tiene contrato escrito.} = Il n'a pas de contrat écrit.

\esp{Cobra un sueldo cada mes.} = Elle touche un salaire chaque mois.

\esp{Piensa en su jubilación.} = Elle pense à sa retraite.
\end{exemplebox}

\begin{attention}[Ne recopie pas de longs passages]
Erreur fréquente au baccalauréat : copier trois lignes du texte comme réponse. Le correcteur veut voir que \emph{toi}, tu as compris. Cite seulement les \cle{mots utiles} et \textbf{reformule}. Mauvais : recopier tout le paragraphe sur l'oncle. Bon : \esp{El tío trabaja sin contrato escrito, por eso su situación es insegura.}
\end{attention}

\courssub{Questions de vocabulaire en contexte}

On te demandera souvent de retrouver dans le texte un synonyme ou un contraire. Il faut chercher un mot \emph{du texte}, pas un mot de ta tête.

\begin{itemize}
\item \esp{Busca en el texto un sinónimo de « salario ».} → \esp{sueldo} (le texte dit : \esp{cobra un sueldo cada mes}).
\item \esp{Busca el contrario de « tener empleo ».} → \esp{estar en paro} / \esp{estar en el paro} (être au chômage).
\end{itemize}

\begin{center}
\begin{tabularx}{\linewidth}{|X|X|X|}
\hline
\rowcolor{popblueL} \textbf{Español} & \textbf{Français} & \textbf{Catégorie} \\
\hline
\esp{el sueldo / el salario} & le salaire & synonymes \\
\hline
\esp{cobrar} & toucher (un salaire) & verbe \\
\hline
\esp{el paro} & le chômage & contraire de \esp{empleo} \\
\hline
\esp{el contrato} & le contrat & vocabulaire du travail \\
\hline
\esp{la jubilación} & la retraite & vocabulaire du travail \\
\hline
\esp{la huelga} & la grève & vocabulaire social \\
\hline
\end{tabularx}
\end{center}

\courssub{Questions d'interprétation}

L'interprétation va plus loin que la lecture : il faut expliquer \emph{pourquoi} l'auteur écrit quelque chose. Exemple :

\esp{¿Por qué dice el autor que los dos trabajan « con dignidad »?}

Réponse attendue : \esp{El autor dice que los dos trabajan con dignidad porque, aunque tienen condiciones muy diferentes, los dos se esfuerzan, ganan su vida con honestidad y son útiles a la sociedad.}

« L'auteur dit que les deux travaillent avec dignité parce que, bien qu'ils aient des conditions très différentes, tous deux font des efforts, gagnent leur vie honnêtement et sont utiles à la société. »

\begin{propriete}[Justifier par le texte]
Toute interprétation doit s'appuyer sur un élément du texte. Formule utile : \esp{El autor dice esto porque en el texto se lee que...} — « L'auteur dit cela parce que dans le texte on lit que... »
\end{propriete}

\courssub{Question personnelle}

La dernière question t'invite à donner ton avis :

\esp{¿Qué oficio te gustaría ejercer y por qué?}

Modèle de réponse : \esp{Me gustaría ejercer el oficio de médico porque quiero ayudar a los enfermos de mi barrio en Yaoundé.} — « J'aimerais exercer le métier de médecin parce que je veux aider les malades de mon quartier à Yaoundé. »

Attention à la construction : \esp{me gustaría} (conditionnel de politesse et de souhait) + infinitif ; la profession \emph{sans article} après \esp{ser} : \esp{quiero ser médico}, jamais \esp{quiero ser un médico}.

\begin{exoresolu}[Questions sur un court texte]
Énoncé. Lis ce texte et réponds par des phrases complètes :

\esp{Mi hermano Carlos es taxista en Douala. Trabaja doce horas al día, pero no tiene contrato ni seguro. Sueña con comprar su propio coche y afiliarse a un sindicato para defender sus derechos.}

1. \esp{¿Qué oficio tiene Carlos?} 2. \esp{¿Dónde trabaja?} 3. \esp{Busca en el texto el contrario de « tener derechos protegidos ».} 4. \esp{¿Por qué quiere Carlos afiliarse a un sindicato?}
\tcblower
Solution détaillée :

1. \esp{Carlos es taxista.} (reprise des mots de la question, profession sans article après \esp{ser})

2. \esp{Carlos trabaja en Douala.}

3. Le contraire est exprimé par : \esp{no tiene contrato ni seguro} — sa situation n'est pas protégée.

4. \esp{Carlos quiere afiliarse a un sindicato para defender sus derechos.} — « Carlos veut adhérer à un syndicat pour défendre ses droits. » On reprend la structure de la question (\esp{¿por qué quiere...?} → \esp{quiere...}) et on ajoute la justification du texte (\esp{para defender sus derechos}).
\end{exoresolu}

\begin{attention}[Pièges des francophones]
\begin{itemize}
\item Ne traduis pas mot à mot : \esp{¿Qué oficio tiene el tío?} ne se répond pas par \esp{El tío tiene albañil} mais \esp{El tío es albañil} (en espagnol, on \emph{est} son métier avec \esp{ser}).
\item \esp{Ejercer} s'emploie pour une profession (\esp{ejercer el oficio de...}), pas pour un sport.
\item N'oublie jamais l'accent interrogatif : \esp{qué}, \esp{dónde}, \esp{quiénes}, \esp{por qué}, \esp{cuándo}, \esp{cómo}.
\item \esp{Por día} (par jour / à la journée) vs \esp{para defender} (dans le but de défendre) : distinguer \cle{por} (cause, durée, fréquence) et \cle{para} (but, destination).
\end{itemize}
\end{attention}

\begin{savaistu}[Au baccalauréat]
Au Bac, les questions de compréhension portent \textbf{toujours} sur le texte proposé : toute réponse doit pouvoir être justifiée par une ligne du texte. Même la question personnelle doit rester en lien avec le thème. Entraîne-toi à souligner, pendant la lecture, les mots qui répondent à \esp{quién}, \esp{dónde}, \esp{cuándo} et \esp{por qué} : tu gagneras un temps précieux le jour de l'épreuve.
\end{savaistu}

\begin{exercice}[À toi de jouer]
Reprends le texte \esp{« El mundo del trabajo hoy »} de la section 1 et réponds par des phrases complètes :
\begin{enumerate}
\item \esp{¿De qué habla el texto?}
\item \esp{¿Qué diferencia hay entre el sueldo del tío y el de la prima?}
\item \esp{Busca un sinónimo de « trabajo » en el texto.}
\item \esp{¿Qué oficio te gustaría ejercer y por qué?}
\end{enumerate}
\end{exercice}

\begin{corrige}[Exercice : Compréhension de « El mundo del trabajo hoy »]
\begin{enumerate}
\item \esp{El texto habla del mundo del trabajo y presenta dos realidades laborales diferentes: la de un obrero en Douala y la de una enfermera en Madrid.}
\item \esp{El tío cobra por día y no tiene contrato escrito, mientras que la prima cobra un sueldo cada mes, tiene contrato y piensa en su jubilación.}
\item \esp{Oficio} (ou \esp{empleo}, \esp{profesión} selon le passage exact du texte support).
\item \esp{Me gustaría ejercer el oficio de [profession choisie] porque [raison personnelle liée au contexte camerounais ou hispanophone].} \\
Exemple : \esp{Me gustaría ejercer el oficio de ingeniero agrónomo para mejorar la producción de cacao en mi región.}
\end{enumerate}
\end{corrige}

\coursec{Lexique thématique : el mundo del trabajo}

Dans la section précédente, nous avons lu et compris le texte \esp{El mundo del trabajo hoy} : nous savons désormais de quoi il parle et comment il est construit. Mais pour commenter un texte, il ne suffit pas de le comprendre globalement : il faut posséder le \cle{vocabulaire précis} du thème. C'est l'objet de cette section : organiser, champ par champ, tout le lexique du monde du travail, avec son emploi, ses pièges et ses équivalents français. Retenez la méthode : on observe d'abord le mot dans une phrase authentique, puis on le classe dans son champ lexical, puis on l'utilise dans une phrase personnelle.

\courssub{Champ 1 — Les métiers : oficio, profesión, empleo}

Observons d'abord ces trois phrases :

\esp{Mi tío aprendió el oficio de carpintero con su padre.} « Mon oncle a appris le métier de menuisier avec son père. »

\esp{La profesión de médico exige muchos años de estudios.} « La profession de médecin exige de longues années d'études. »

\esp{En Douala, muchos jóvenes buscan un empleo estable.} « À Douala, beaucoup de jeunes cherchent un emploi stable. »

\begin{definition}[Oficio, profesión, empleo, trabajo, puesto, carrera]
\esp{Oficio} : métier manuel ou artisanal, appris par la pratique (\esp{carpintero}, \esp{zapatero}, \esp{peluquero}). \esp{Profesión} : profession en général, souvent intellectuelle ou diplômée (\esp{médico}, \esp{abogado}, \esp{profesor}). \esp{Empleo} : emploi, poste occupé contre un salaire. \esp{Trabajo} : le travail en général (mot le plus large). \esp{Puesto} : poste, fonction précise dans une entreprise (\esp{un puesto de responsabilidad} = un poste de responsabilité). \esp{Carrera} : carrière, parcours professionnel (\esp{hacer carrera} = faire carrière).
\end{definition}

\begin{definition}[Les acteurs du monde du travail]
\esp{Artesano/a} : artisan(e) ; \esp{obrero/a} : ouvrier(-ière) ; \esp{empleado/a} : employé(e) ; \esp{empresario/a} : chef(fe) d'entreprise ; \esp{jefe/a} : chef(fe), supérieur(e) direct(e) ; \esp{compañero/a de trabajo} : collègue de travail.
\end{definition}

\begin{attention}[Faux amis et calques]
\esp{Carrera} ne signifie pas « carrière » au sens de « carrière de pierre » : c'est le parcours professionnel, ou les études universitaires (\esp{la carrera de Derecho} = les études de droit). Attention aussi : \esp{oficio} ne veut pas dire « office » religieux dans ce contexte. Enfin, « le patron » se dit \esp{el jefe} ou \esp{el empresario} en Espagne ; en Amérique latine, \esp{el patrón} est le terme courant pour « le patron, l'employeur ». Ne traduisez donc pas « patron » par un mot inventé : choisissez selon le contexte.
\end{attention}

\begin{popfigure}
\begin{center}
\begin{tabular}{|l|l|}
\hline
\rowcolor{popblueL} \textbf{Español} & \textbf{Français} \\
\hline
el oficio & le métier (manuel) \\
\hline
la profesión & la profession \\
\hline
el empleo & l'emploi \\
\hline
el trabajo & le travail \\
\hline
el puesto & le poste \\
\hline
la carrera & la carrière \\
\hline
el/la artesano/a & l'artisan(e) \\
\hline
el/la obrero/a & l'ouvrier(-ière) \\
\hline
el/la empleado/a & l'employé(e) \\
\hline
el/la empresario/a & le/la chef(fe) d'entreprise \\
\hline
el/la jefe/a & le chef, la cheffe \\
\hline
el/la compañero/a de trabajo & le/la collègue \\
\hline
\end{tabular}
\end{center}
\legende{Tableau lexical : les métiers et les acteurs du travail}
\end{popfigure}

\courssub{Champ 2 — L'embauche et le contrat}

Observons le parcours d'une candidate :

\esp{Ana envió su currículum, pasó una entrevista de trabajo y firmó un contrato con una jornada laboral de ocho horas.} « Ana a envoyé son CV, a passé un entretien d'embauche et a signé un contrat avec une journée de travail de huit heures. »

\begin{definition}[Embauche et contrat]
\esp{Contrato} : contrat ; \esp{contratar} : embaucher ; \esp{despedir} : licencier (verbe irrégulier : \esp{despido, despides, despide\ldots}) ; \esp{entrevista de trabajo} : entretien d'embauche ; \esp{currículum} (ou \esp{currículum vitae}) : CV ; \esp{jornada laboral} : journée de travail, durée du travail ; \esp{salario} et \esp{sueldo} : salaire ; \esp{aumento} : augmentation (\esp{un aumento de sueldo} = une augmentation de salaire).
\end{definition}

\begin{propriete}[Contratar et despedir : verbes contraires]
\esp{Contratar} (régulier en \esp{-ar}) = embaucher : \esp{La empresa contrató a veinte obreros.} « L'entreprise a embauché vingt ouvriers. »

\esp{Despedir} (e $\to$ i, comme \esp{pedir}) = licencier : \esp{Despidieron al empleado por llegar tarde.} « Ils ont licencié l'employé parce qu'il arrivait en retard. »

Noms correspondants : \esp{el despido} = le licenciement ; \esp{el contrato} = le contrat. Notez la préposition : \esp{contratar a alguien}, \esp{despedir a alguien} (a personnel devant la personne).
\end{propriete}

\begin{attention}[Despedir : deux sens à ne pas mélanger]
\esp{Despedir} signifie « licencier », mais le verbe pronominal \esp{despedirse (de)} signifie « dire au revoir, prendre congé de » : \esp{Me despedí de mis compañeros.} « J'ai dit au revoir à mes collègues. » Ne traduisez donc jamais \esp{se despidió} par « il a été licencié » sans vérifier le contexte !
\end{attention}

\begin{attention}[Sueldo ou suelo ?]
Piège classique du francophone : \esp{el sueldo} = le salaire ; \esp{el suelo} = le sol, le plancher. Une seule lettre de différence, deux mondes ! \esp{Recibo mi sueldo a fin de mes.} « Je reçois mon salaire à la fin du mois. » \esp{El vaso cayó al suelo.} « Le verre est tombé par terre. »
\end{attention}

\begin{popfigure}
\begin{center}
\begin{tabular}{|l|l|}
\hline
\rowcolor{popgreenL} \textbf{Español} & \textbf{Français} \\
\hline
el contrato & le contrat \\
\hline
contratar & embaucher \\
\hline
despedir & licencier \\
\hline
el despido & le licenciement \\
\hline
la entrevista de trabajo & l'entretien d'embauche \\
\hline
el currículum & le CV \\
\hline
la jornada laboral & la journée de travail \\
\hline
el salario / el sueldo & le salaire \\
\hline
el aumento & l'augmentation \\
\hline
\end{tabular}
\end{center}
\legende{Tableau lexical : l'embauche et le contrat}
\end{popfigure}

\courssub{Champ 3 — Le chômage et la crise}

\esp{Desde la crisis económica, muchos jóvenes están en paro y buscan trabajo sin éxito.} « Depuis la crise économique, beaucoup de jeunes sont au chômage et cherchent du travail sans succès. »

\begin{definition}[Paro et desempleo]
\esp{El paro} = le chômage (terme courant en Espagne) ; \esp{el desempleo} = le chômage (terme universel, utilisé partout). \esp{El/la parado/a} = le/la chômeur(-euse) (Espagne) ; \esp{el/la desempleado/a} = le/la chômeur(-euse) (partout). \esp{Buscar trabajo} = chercher du travail ; \esp{la crisis económica} = la crise économique ; \esp{la precariedad} = la précarité.
\end{definition}

\begin{exemplebox}[Exemples traduits]
\esp{Está en paro desde hace un año.} « Il est au chômage depuis un an. »

\esp{El desempleo juvenil preocupa a los gobiernos.} « Le chômage des jeunes inquiète les gouvernements. »

\esp{Muchos parados viven en la precariedad.} « Beaucoup de chômeurs vivent dans la précarité. »
\end{exemplebox}

\begin{attention}[Parado : attention au calque]
\esp{Parado} signifie « au chômage » en Espagne, mais aussi « arrêté, immobile » ou « debout » dans plusieurs pays d'Amérique. Pour éviter toute ambiguïté à l'examen, préférez \esp{estar en paro} ou \esp{estar desempleado}. Et ne dites jamais \esp{estar en desempleo} : les constructions correctes sont \esp{estar en paro} ou \esp{estar desempleado}. Notez aussi l'expression de la durée avec le présent : \esp{Lleva dos años en paro.} « Il est au chômage depuis deux ans. »
\end{attention}

\begin{popfigure}
\begin{center}
\begin{tabular}{|l|l|}
\hline
\rowcolor{poporangeL} \textbf{Español} & \textbf{Français} \\
\hline
el paro / el desempleo & le chômage \\
\hline
el/la parado/a & le/la chômeur(-euse) (Espagne) \\
\hline
el/la desempleado/a & le/la chômeur(-euse) \\
\hline
buscar trabajo & chercher du travail \\
\hline
la crisis económica & la crise économique \\
\hline
la precariedad & la précarité \\
\hline
\end{tabular}
\end{center}
\legende{Tableau lexical : le chômage et la crise}
\end{popfigure}

\courssub{Champ 4 — La défense des travailleurs}

\begin{definition}[Sindicato et huelga]
\esp{El sindicato} : association de travailleurs qui défend leurs droits (le syndicat). \esp{El/la sindicalista} : le/la syndicaliste, membre d'un syndicat. \esp{La huelga} : arrêt collectif du travail pour protester (la grève). \esp{La manifestación} : la manifestation ; \esp{la negociación} : la négociation ; \esp{el convenio colectivo} : la convention collective.
\end{definition}

\begin{exemplebox}[Exemples traduits]
\esp{Los obreros se pusieron en huelga.} « Les ouvriers se sont mis en grève. »

\esp{El sindicato negocia un nuevo convenio colectivo.} « Le syndicat négocie une nouvelle convention collective. »

\esp{Hubo una manifestación en el centro de la ciudad.} « Il y a eu une manifestation au centre-ville. »
\end{exemplebox}

\begin{propriete}[Les constructions avec huelga]
Être en grève : \esp{estar en huelga} ; se mettre en grève : \esp{ponerse en huelga} ; faire grève : \esp{hacer huelga}. Exemple : \esp{Los empleados hicieron huelga durante tres días.} « Les employés ont fait grève pendant trois jours. »
\end{propriete}

\begin{attention}[Manifestación et negociación : deux pièges]
\esp{Manifestación} désigne spécifiquement le défilé de protestation ; une « manifestation » sportive ou culturelle se dit plutôt \esp{un acto} ou \esp{un evento}. Autre piège : \esp{negociación} est le nom ; le verbe est \esp{negociar} (négocier), sans aucun rapport avec le français « négliger » !
\end{attention}

\begin{popfigure}
\begin{center}
\begin{tabular}{|l|l|}
\hline
\rowcolor{poppurpleL} \textbf{Español} & \textbf{Français} \\
\hline
el sindicato & le syndicat \\
\hline
el/la sindicalista & le/la syndicaliste \\
\hline
la huelga & la grève \\
\hline
la manifestación & la manifestation \\
\hline
la negociación & la négociation \\
\hline
negociar & négocier \\
\hline
el convenio colectivo & la convention collective \\
\hline
\end{tabular}
\end{center}
\legende{Tableau lexical : la défense des travailleurs}
\end{popfigure}

\courssub{Champ 5 — Droits et devoirs des travailleurs}

\begin{definition}[Los derechos de los trabajadores]
\esp{El derecho al trabajo} : le droit au travail ; \esp{el derecho a un salario justo} : le droit à un salaire juste ; \esp{el derecho a la seguridad social} : le droit à la sécurité sociale ; \esp{el derecho a la jubilación} : le droit à la retraite. Structure : \esp{tener derecho a + nom} ou \esp{tener derecho a + infinitif} : \esp{Todo trabajador tiene derecho a descansar.} « Tout travailleur a le droit de se reposer. »
\end{definition}

\begin{definition}[Los deberes de los trabajadores]
\esp{Cumplir el horario} : respecter l'horaire ; \esp{respetar al jefe y a los compañeros} : respecter le chef et les collègues ; \esp{trabajar con honradez} : travailler avec honnêteté. Structures : \esp{tener el deber de + infinitif} ou \esp{deber + infinitif} : \esp{Los empleados deben cumplir el horario.} « Les employés doivent respecter l'horaire. »
\end{definition}

\begin{exemplebox}[Exemples traduits]
\esp{Todo ciudadano tiene derecho al trabajo.} « Tout citoyen a droit au travail. »

\esp{Un buen empleado trabaja con honradez y respeta a sus compañeros.} « Un bon employé travaille avec honnêteté et respecte ses collègues. »
\end{exemplebox}

\begin{attention}[Derecho : un mot, plusieurs sens]
\esp{El derecho} est ici un nom masculin : le droit (\esp{los derechos del trabajador} = les droits du travailleur). Ne pas confondre avec l'adverbe \esp{derecho} (« tout droit » : \esp{siga derecho} = continuez tout droit) ni avec l'adjectif \esp{derecho/a} (« droit », opposé de gauche : \esp{la mano derecha} = la main droite). Pour le devoir professionnel, dites \esp{trabajar con honradez} (avec honnêteté, probité).
\end{attention}

\begin{popfigure}
\begin{center}
\begin{tabular}{|l|l|}
\hline
\rowcolor{popgoldL} \textbf{Derechos (droits)} & \textbf{Deberes (devoirs)} \\
\hline
derecho al trabajo & cumplir el horario \\
\hline
derecho a un salario justo & respetar al jefe \\
\hline
derecho a la seguridad social & respetar a los compañeros \\
\hline
derecho a la jubilación & trabajar con honradez \\
\hline
\end{tabular}
\end{center}
\legende{Tableau lexical : droits et devoirs des travailleurs}
\end{popfigure}

\courssub{Champ 6 — La retraite : jubilación}

\esp{Mi abuelo se jubiló a los sesenta años y ahora es un jubilado feliz que recibe su pensión cada mes.} « Mon grand-père a pris sa retraite à soixante ans et c'est maintenant un retraité heureux qui reçoit sa pension chaque mois. »

\begin{definition}[La famille de jubilación]
\esp{La jubilación} : la retraite (le fait, le moment) ; \esp{jubilarse} : prendre sa retraite (verbe pronominal régulier en \esp{-ar}) ; \esp{el/la jubilado/a} : le/la retraité(e) ; \esp{la pensión} : la pension (l'argent reçu).
\end{definition}

\begin{exemplebox}[Exemples traduits]
\esp{Se jubiló a los sesenta años.} « Il a pris sa retraite à soixante ans. »

\esp{Los jubilados reciben una pensión del Estado.} « Les retraités reçoivent une pension de l'État. »
\end{exemplebox}

\begin{attention}[Jubilación et pensión : ne pas confondre]
\esp{Jubilación} désigne l'événement (la cessation d'activité), \esp{pensión} désigne l'argent. On dit : \esp{pedir la jubilación} (demander sa retraite) mais \esp{cobrar la pensión} (toucher sa pension). Autre piège : \esp{pensión} signifie aussi « pension de famille » (logement) : c'est le contexte qui décide.
\end{attention}

\begin{popfigure}
\begin{center}
\begin{tikzpicture}[node distance=0.35cm and 0.4cm,
  etape/.style={rectangle, rounded corners, draw=popblue, fill=popblueL, align=center, minimum height=0.9cm, font=\small},
  fleche/.style={-{Stealth[length=2.5mm]}, thick, poporange}]
\node[etape] (e1) {estudios};
\node[etape, right=of e1] (e2) {entrevista};
\node[etape, right=of e2] (e3) {contrato};
\node[etape, right=of e3] (e4) {trabajo};
\node[etape, right=of e4] (e5) {aumento};
\node[etape, right=of e5] (e6) {jubilación};
\draw[fleche] (e1) -- (e2);
\draw[fleche] (e2) -- (e3);
\draw[fleche] (e3) -- (e4);
\draw[fleche] (e4) -- (e5);
\draw[fleche] (e5) -- (e6);
\end{tikzpicture}
\end{center}
\legende{Frise de la vie professionnelle : des études à la retraite}
\end{popfigure}

\courssub{Expressions utiles à mémoriser}

\begin{aretenir}[Expressions clés du monde du travail]
\begin{itemize}
\item \esp{ganarse la vida} = gagner sa vie : \esp{Me gano la vida como taxista.} « Je gagne ma vie comme chauffeur de taxi. »
\item \esp{trabajar a tiempo completo} = travailler à temps plein ; \esp{trabajar a tiempo parcial} = travailler à temps partiel.
\item \esp{estar en paro} = être au chômage.
\item \esp{ponerse en huelga} = se mettre en grève.
\item \esp{firmar un contrato} = signer un contrat.
\item \esp{tener derecho a + infinitif} = avoir le droit de ; \esp{deber + infinitif} = devoir.
\end{itemize}
\end{aretenir}

\begin{savaistu}[Ganarse la vida]
\esp{Ganarse la vida} (littéralement « se gagner la vie ») est l'expression la plus courante dans tout le monde hispanophone pour dire « gagner sa vie » : \esp{Se gana la vida vendiendo fruta en el mercado.} « Elle gagne sa vie en vendant des fruits au marché. » Au Cameroun, on pourrait dire d'un motocycliste : \esp{El mototaxista se gana la vida transportando pasajeros por Yaoundé.} « Le moto-taximan gagne sa vie en transportant des passagers dans Yaoundé. »
\end{savaistu}

\courssub{Texte d'application : Un día en la oficina de empleo}

\begin{textebox}[Un día en la oficina de empleo (texte original)]
María entró en la oficina de empleo de su barrio con su currículum bajo el brazo. Tenía veinticinco años y estaba en paro desde hacía ocho meses. Después de terminar sus estudios, había trabajado como empleada en una pequeña empresa, pero la crisis económica obligó al empresario a despedir a la mitad de los obreros y empleados.

En la sala de espera, María habló con un señor mayor, un artesano carpintero que buscaba trabajo porque su taller había cerrado. «Antes, los oficios se aprendían con los padres y duraban toda la vida», suspiró él. «Ahora todo cambia muy rápido.»

Por fin, llegó el turno de María. El entrevistador leyó su currículum con atención y le propuso una entrevista de trabajo en una fábrica de cacao de las afueras de la ciudad. «El contrato sería a tiempo completo, con una jornada laboral de ocho horas y un sueldo justo», explicó. «Además, tendría derecho a la seguridad social y, después de muchos años, a la jubilación.»

María sonrió por primera vez en meses. Sabía que un empleado también tiene deberes: cumplir el horario, respetar al jefe y a los compañeros y trabajar con honradez. Pero estaba dispuesta a todo con tal de ganarse la vida dignamente. Al salir, pensó en su tío, sindicalista jubilado, que siempre repetía: «Los derechos de los trabajadores no se regalan, se conquistan.» Ella, por ahora, solo quería firmar un contrato y empezar a trabajar.
\end{textebox}

\textbf{Glossaire :} \esp{bajo el brazo} = sous le bras ; \esp{la sala de espera} = la salle d'attente ; \esp{el taller} = l'atelier ; \esp{suspirar} = soupirer ; \esp{el turno} = le tour ; \esp{las afueras} = la périphérie ; \esp{dispuesta a todo} = prête à tout ; \esp{con tal de + infinitif} = pourvu de, afin de ; \esp{dignamente} = dignement ; \esp{no se regalan, se conquistan} = ne sont pas offerts, ils se conquièrent.

\textbf{Questions de compréhension :}
\begin{enumerate}
\item ¿Por qué está María en paro?
\item ¿Qué oficio tenía el señor mayor de la sala de espera?
\item ¿Qué le propone el entrevistador a María?
\item ¿Qué derechos ofrece el contrato de la fábrica?
\item ¿Qué deberes conoce María?
\end{enumerate}

\textbf{Éléments de réponses :} 1. Parce que la crisis económica obligó a su empresario a despedir a la mitad de los trabajadores. 2. Era artesano carpintero. 3. Le propone una entrevista de trabajo en una fábrica de cacao. 4. Un sueldo justo, la seguridad social y la jubilación. 5. Cumplir el horario, respetar al jefe y a los compañeros y trabajar con honradez.

\courssub{Exercice résolu}

\begin{exoresolu}[Complétez avec le mot qui convient]
Complétez les phrases avec : \esp{paro}, \esp{contrato}, \esp{huelga}, \esp{jubilación}, \esp{sueldo}, \esp{sindicato}.
\begin{enumerate}
\item Después de la entrevista, la empresa me ofreció un \ldots{} a tiempo completo.
\item Mi padre recibe su \ldots{} el último día del mes.
\item Los obreros se pusieron en \ldots{} para protestar contra los despidos.
\item El \ldots{} defiende los derechos de los trabajadores.
\item Después de cuarenta años de trabajo, pidió la \ldots{}.
\item Lleva dos años en \ldots{} y busca trabajo sin éxito.
\end{enumerate}
\tcblower
\textbf{Solution détaillée :}
\begin{enumerate}
\item \esp{contrato} : après un entretien, l'entreprise propose un contrat. \esp{Después de la entrevista, la empresa me ofreció un contrato a tiempo completo.} « Après l'entretien, l'entreprise m'a proposé un contrat à temps plein. »
\item \esp{sueldo} : on reçoit son salaire en fin de mois (attention : pas \esp{suelo}, qui signifie « le sol » !).
\item \esp{huelga} : \esp{ponerse en huelga} = se mettre en grève.
\item \esp{sindicato} : c'est le syndicat qui défend les droits des travailleurs.
\item \esp{jubilación} : \esp{pedir la jubilación} = demander sa retraite.
\item \esp{paro} : \esp{estar en paro} = être au chômage ; \esp{llevar + durée} exprime « depuis ».
\end{enumerate}
\end{exoresolu}

\begin{exercice}[À vous de jouer]
Traduisez en espagnol : 1. Il est au chômage depuis un an. 2. Les ouvriers se sont mis en grève. 3. Elle a signé un contrat à temps partiel. 4. Mon grand-père a pris sa retraite à soixante ans. 5. Tout travailleur a droit à un salaire juste.
\end{exercice}

\begin{corrige}[Exercice « À vous de jouer »]
1. \esp{Está en paro desde hace un año.} (ou \esp{Lleva un año en paro.}) 2. \esp{Los obreros se pusieron en huelga.} 3. \esp{Firmó un contrato a tiempo parcial.} 4. \esp{Mi abuelo se jubiló a los sesenta años.} 5. \esp{Todo trabajador tiene derecho a un salario justo.}
\end{corrige}

Vous disposez maintenant de tout le lexique nécessaire : métiers, embauche, chômage, syndicats, droits, devoirs et retraite. La section suivante vous montrera comment mobiliser ce vocabulaire pour présenter les droits et devoirs des travailleurs dans un commentaire de texte.

\coursec{Droits et devoirs des travailleurs}

Dans la section précédente, nous avons construit le lexique essentiel du monde du travail : \esp{oficio}, \esp{profesión}, \esp{empleo}, \esp{paro}, \esp{sueldo}, \esp{contrato}, \esp{sindicato}, \esp{huelga}, \esp{jubilación}. Mais travailler, ce n'est pas seulement exercer un métier : c'est aussi être protégé par des \cle{droits} (\esp{derechos}) et être tenu à des \cle{devoirs} (\esp{deberes}). C'est précisément ce vocabulaire et ces structures grammaticales que l'examinateur du Baccalauréat attend dans un essai, un commentaire ou une traduction sur le thème du travail. Observons d'abord un texte authentique qui met en scène ces notions.

\courssub{Texte d'observation : derechos y deberes}

\begin{textebox}[Un día en la fábrica de cacao]
Carmen trabaja desde hace diez años en una fábrica de chocolate cerca de Douala. Como todos los trabajadores, tiene derecho a un contrato escrito, a un salario justo y a la seguridad social. Cada año, disfruta de vacaciones pagadas y puede afiliarse a un sindicato para defender sus intereses. Si un día las condiciones de trabajo se vuelven injustas, la ley reconoce su derecho a la huelga. Cuando sea mayor, tendrá derecho a una jubilación digna después de tantos años de esfuerzo.

Pero Carmen sabe que los derechos van acompañados de deberes. Debe cumplir el horario: llega a la fábrica a las siete de la mañana y nunca llega tarde. Tiene que realizar bien su trabajo y respetar las normas de seguridad, porque las máquinas pueden ser peligrosas. Además, está obligada a respetar a sus compañeros y a ser honrada en todo momento.

El mes pasado, los obreros de la fábrica pidieron un aumento de salario. El sindicato negoció con el patrón durante dos semanas. Al final, consiguieron una mejora sin necesidad de hacer huelga. «Los derechos se conquistan con diálogo y con trabajo serio», dice Carmen con una sonrisa.
\end{textebox}

\begin{itemize}
\item \esp{seguridad social} = sécurité sociale
\item \esp{vacaciones pagadas} = congés payés
\item \esp{afiliarse a un sindicato} = adhérer à un syndicat
\item \esp{jubilación digna} = retraite digne
\item \esp{cumplir el horario} = respecter l'horaire
\item \esp{normas de seguridad} = règles de sécurité
\item \esp{compañeros} = collègues
\item \esp{patrón} = patron, employeur
\item \esp{hacer huelga} = faire grève
\item \esp{aumento de salario} = augmentation de salaire
\item \esp{conseguir una mejora} = obtenir une amélioration
\item \esp{sin necesidad de} = sans avoir besoin de
\end{itemize}

\textbf{Questions de compréhension :}
\begin{enumerate}
\item ¿Qué derechos tiene Carmen en la fábrica ?
\item ¿Cuáles son sus deberes como trabajadora ?
\item ¿Cómo consiguieron los obreros el aumento de salario ?
\item ¿Qué frase del texto muestra que derechos y deberes van juntos ?
\end{enumerate}

\courssub{Les droits fondamentaux des travailleurs}

\begin{definition}[Los derechos del trabajador]
Un \esp{derecho} est un droit : ce que la loi garantit au travailleur et que l'employeur doit respecter. Les droits fondamentaux sont :
\begin{itemize}
\item \esp{el derecho al trabajo} : le droit au travail ;
\item \esp{el derecho a un salario justo} : le droit à un salaire juste ;
\item \esp{el derecho a un contrato} : le droit à un contrat (écrit) ;
\item \esp{el derecho a la seguridad social} : le droit à la sécurité sociale ;
\item \esp{el derecho a las vacaciones pagadas} : le droit aux congés payés ;
\item \esp{el derecho a afiliarse a un sindicato} : le droit d'adhérer à un syndicat ;
\item \esp{el derecho a la huelga} : le droit de grève ;
\item \esp{el derecho a la jubilación} : le droit à la retraite.
\end{itemize}
\end{definition}

Observez les exemples tirés du texte : \esp{tiene derecho a un contrato}, \esp{derecho a la huelga}, \esp{tendrá derecho a una jubilación}. Quelle préposition revient toujours après \esp{derecho} ? La préposition \esp{a}.

\begin{propriete}[La structure « tener derecho a »]
\esp{Tener derecho a} se construit avec la préposition \esp{a} + nom ou + infinitif :
\begin{itemize}
\item \esp{Tengo derecho a descansar.} = J'ai le droit de me reposer.
\item \esp{Los trabajadores tienen derecho a un salario justo.} = Les travailleurs ont droit à un salaire juste.
\item \esp{Todo ciudadano tiene derecho a expresar su opinión.} = Tout citoyen a le droit d'exprimer son opinion.
\end{itemize}
On ne dit jamais \esp{*tener derecho de} : la préposition \esp{a} est obligatoire, même quand le français dit « droit \emph{de} ».
\end{propriete}

\begin{attention}[Piège de la préposition : \cle{derecho a}]
Le francophone traduit « le droit de faire grève » par \esp{*el derecho de hacer huelga} : c'est une faute grave. En espagnol, on dit toujours \esp{el derecho a hacer huelga}. Même logique que \esp{soñar con}, \esp{acordarse de}, \esp{depender de} : chaque nom ou verbe impose sa préposition, et il faut l'apprendre avec lui. Retenezonce pour toutes : \esp{derecho a}, \esp{tener derecho a}.
\end{attention}

\courssub{Les devoirs des travailleurs}

\begin{definition}[Los deberes del trabajador]
Un \esp{deber} est un devoir : ce que le travailleur doit faire en échange de ses droits. Les principaux devoirs sont :
\begin{itemize}
\item \esp{cumplir el horario} : respecter l'horaire ;
\item \esp{realizar bien el trabajo} : bien faire son travail ;
\item \esp{respetar las normas de seguridad} : respecter les règles de sécurité ;
\item \esp{ser puntual y honrado} : être ponctuel et honnête ;
\item \esp{respetar a los compañeros} : respecter ses collègues (personne \rightarrow \esp{a}).
\end{itemize}
\end{definition}

Dans le texte, trois structures expriment le devoir : \esp{debe cumplir}, \esp{tiene que realizar}, \esp{está obligada a respetar}. Comparons leurs nuances.

\begin{propriete}[Exprimer le devoir : deber, tener que, estar obligado a]
Trois structures principales, toujours suivies de l'infinitif :
\begin{itemize}
\item \esp{deber + infinitif} : le devoir moral, éthique ou légal (obligation interne ou norme). \esp{El empleado debe cumplir el horario.} = L'employé doit (a le devoir de) respecter l'horaire.
\item \esp{tener que + infinitif} : l'obligation concrète, matérielle, souvent imposée par les circonstances. \esp{Tiene que trabajar ocho horas al día.} = Il doit (il est obligé de par les faits) travailler huit heures par jour.
\item \esp{estar obligado a + infinitif} : l'obligation forte, formelle, imposée par une règle, un contrat ou la loi (accord du participe \esp{obligado/a/os/as} avec le sujet). \esp{Los obreros están obligados a llevar casco.} = Les ouvriers sont obligés de porter un casque.
\end{itemize}
\end{propriete}

\begin{exemplebox}[Exemples traduits]
\esp{Todo trabajador tiene derecho a un contrato.} = Tout travailleur a droit à un contrat.\\
\esp{Los empleados deben respetar el horario.} = Les employés doivent respecter l'horaire.\\
\esp{Tiene que trabajar ocho horas al día.} = Il doit travailler huit heures par jour.\\
\esp{Tenemos derecho a vacaciones pagadas.} = Nous avons droit à des congés payés.\\
\esp{Debes ser puntual y honrado en el trabajo.} = Tu dois être ponctuel et honnête au travail.\\
\esp{El patrón está obligado a pagar la seguridad social.} = Le patron est obligé de payer la sécurité sociale.
\end{exemplebox}

\begin{attention}[Deber ou deber de ? La nuance capitale pour le Bac]
\esp{Deber + infinitif} exprime le \cle{devoir / l'obligation} : \esp{Debe trabajar más.} = Il doit travailler davantage (obligation).\\
Mais \esp{deber de + infinitif} exprime la \cle{supposition / la probabilité} : \esp{Debe de estar cansado.} = Il doit être fatigué (je suppose qu'il l'est).\\
Cette nuance est systématiquement testée au Bac dans les exercices de traduction (thème/version) :
\begin{itemize}
\item « Elle doit être malade » (supposition) $\rightarrow$ \esp{Debe de estar enferma.}
\item « Elle doit être malade » (obligation absurde ici, mais possible : « Elle a pour devoir d'être malade ») $\rightarrow$ \esp{Debe estar enferma.}
\end{itemize}
\emph{Astuce} : Si vous pouvez remplacer par « il est probable que » + subjonctif (\esp{es probable que esté enfermo}), c'est \esp{deber de}.
\end{attention}

\begin{center}
\begin{tabularx}{\linewidth}{|X|X|X|}
\hline
\rowcolor{popblueL} Français & Español & Remarque \\
\hline
avoir droit à / le droit de & tener derecho a & toujours la préposition \esp{a} (jamais \esp{de}) \\
\hline
devoir (obligation morale/légale) & deber + inf. & \esp{Debes estudiar.} \\
\hline
devoir (nécessité concrète) & tener que + inf. & \esp{Tengo que irme ya.} \\
\hline
être obligé de & estar obligado a + inf. & accord de \esp{obligado} ; \esp{Estamos obligados a firmar.} \\
\hline
devoir (supposition) & deber de + inf. & \esp{Debe de llover.} = Il doit pleuvoir. \\
\hline
faire grève & hacer huelga & activité : sans article ; événement : \esp{convocar una huelga} \\
\hline
prendre sa retraite & jubilarse & verbe pronominal : \esp{Me jubilo el año que viene.} \\
\hline
\end{tabularx}
\end{center}

\begin{popfigure}
\begin{tikzpicture}[node distance=1.5cm]
\node[draw=popblue, fill=popblueL, rounded corners, thick, minimum width=7cm, align=center, font=\bfseries] (der) {Derechos del trabajador};
\node[draw=poporange, fill=poporangeL, rounded corners, thick, minimum width=7cm, align=center, font=\bfseries, right=of der] (deb) {Deberes del trabajador};
\node[draw=popblue, fill=white, rounded corners, align=left, font=\small, text width=6.5cm, below=of der] (derl) {
\textbullet\ derecho al trabajo\\
\textbullet\ derecho a un salario justo\\
\textbullet\ derecho a un contrato\\
\textbullet\ derecho a vacaciones pagadas\\
\textbullet\ derecho a la huelga y a la jubilación};
\node[draw=poporange, fill=white, rounded corners, align=left, font=\small, text width=6.5cm, below=of deb] (debl) {
\textbullet\ cumplir el horario\\
\textbullet\ realizar bien el trabajo\\
\textbullet\ respetar las normas de seguridad\\
\textbullet\ ser puntual y honrado\\
\textbullet\ respetar a los compañeros};
\draw[<->, thick, popgreen] (der.east) -- node[above, font=\small\itshape, popgreen]{equilibrio} (deb.west);
\end{tikzpicture}
\legende{Droits et devoirs : les deux faces indissociables du monde du travail.}
\end{popfigure}

\courssub{Exercice résolu : thème et structures}

\begin{exoresolu}[Traduction et choix des structures]
Traduisez en espagnol :
1) Tout travailleur a droit à un salaire juste.
2) Les employés doivent respecter les règles de sécurité.
3) Elle doit être fatiguée après le travail (supposition).
4) Nous sommes obligés de porter un uniforme.
\tcblower
1) \esp{Todo trabajador tiene derecho a un salario justo.} — Préposition \esp{a} obligatoire après \esp{derecho}. \esp{Salario justo} (collocation fréquente au Bac).\\
2) \esp{Los empleados deben respetar las normas de seguridad.} — \esp{Deber} + infinitif pour exprimer le devoir légal / la norme. \esp{Normas de seguridad} (lexique imposé).\\
3) \esp{Debe de estar cansada después del trabajo.} — \esp{Deber de} + infinitif car il s'agit d'une \emph{supposition} (probabilité), non d'une obligation. Accord de \esp{cansada} (féminin).\\
4) \esp{Estamos obligados a llevar un uniforme.} — \esp{Estar obligado a} + infinitif ; \esp{obligados} accordé au sujet \emph{nosotros} (masculin pluriel par défaut ou mixte). \esp{Llevar} (porter/vêtir) et non \esp{poner}.
\end{exoresolu}

\courssub{Méthode : Commenter un texte sur le monde du travail}

\begin{methode}[Commentaire de texte type Bac — 4 étapes]
\textbf{Étape 1 : Présentation et problématisation (5 min)}
\begin{enumerate}
\item Identifier la nature du document (article de presse, témoignage, extrait littéraire, dialogue), sa source, sa date, son auteur.
\item Définir le thème général (ici : droits/devoirs, conditions de travail, rôle syndical).
\item Formuler une \cle{problématique} : \emph{« En quoi ce texte illustre-t-il l'équilibre entre droits et devoirs dans le monde du travail actuel ? »} ou \emph{« Comment le dialogue social permet-il de faire respecter les droits des travailleurs ? »}
\end{enumerate}

\textbf{Étape 2 : Analyse structurée du contenu (10 min)}
\begin{enumerate}
\item Repérer les \cle{idées principales} (souvent 2 ou 3) : ex. 1. L'énumération des droits (contrat, salaire, sécurité, syndicat, grève, retraite). 2. L'énumération des devoirs (horaire, qualité, sécurité, honnêteté). 3. La résolution du conflit par la négociation (rôle du syndicat).
\item Citer des \cle{arguments} et \cle{exemples} précis du texte (citations courtes intégrées à la phrase).
\item Noter le \cle{ton} (réaliste, revendicatif, optimiste, critique) et le \cle{point de vue} (ici : focalisation interne sur Carmen).
\end{enumerate}

\textbf{Étape 3 : Analyse linguistique — Le « commentaire de langue » (10 min)}
\begin{enumerate}
\item \textbf{Lexique thématique} : relever les mots du champ lexical (trabajador, fábrica, contrato, salario, sindicato, huelga, jubilación, obreros, patrón, aumento).
\item \textbf{Grammaire / Structures} : identifier et commenter l'usage des structures ciblées :
\begin{itemize}
\item \esp{Tener derecho a} + nom/infinitif (droits).
\item \esp{Deber / Tener que / Estar obligado a} + infinitif (devoirs).
\item \esp{Deber de} + infinitif (si supposition présente).
\item Connecteurs logiques : \esp{pero, además, porque, al final, sin necesidad de}.
\item Temps verbaux : présent (réalité actuelle), passé (récit : \esp{pidieron, negoció, consiguieron}), futur (projection : \esp{tendrá}), subjonctif (antériorité future : \esp{cuando sea}).
\end{itemize}
\end{enumerate}

\textbf{Étape 4 : Synthèse et ouverture personnelle (5 min)}
\begin{enumerate}
\item Répondre à la problématique en résumant l'apport du texte.
\item Ouvrir sur une comparaison (Cameroun / Espagne / Amérique latine), une expérience personnelle ou une réflexion citoyenne : \emph{« Ce texte montre que les droits ne sont pas acquis définitivement ; ils se conquièrent par le dialogue (syndicat) et s'accompagnent de responsabilités. Au Cameroun comme en Guinée équatoriale, le Code du travail prévoit ces mêmes garanties, mais leur application effective reste un défi quotidien. »}
\end{enumerate}
\end{methode}

\courssub{Comparaison : Cameroun et monde hispanophone}

Au Bac, les sujets d'essai ou de commentaire demandent souvent de confronter les réalités. Voici des repères structurés pour votre argumentation.

\begin{center}
\begin{tabularx}{\linewidth}{|X|X|X|}
\hline
\rowcolor{popblueL} Aspect & Cameroun & Monde hispanophone (Espagne / Am. Latine / Guinée équatoriale) \\
\hline
Cadre légal principal & Code du travail (Loi n°92/007) & \esp{Estatuto de los Trabajadores} (Espagne) ; Codes du travail nationaux \\
\hline
Salaire minimum & SMIG (Salaire Minimum Interprofessionnel Garanti) & \esp{Salario Mínimo Interprofesional (SMI)} (Espagne) ; variables selon pays \\
\hline
Syndicats majeurs & CSTC, USLC, CCT, etc. & \esp{UGT, CCOO} (Espagne) ; \esp{CTA} (Argentine) ; \esp{CUT} (Chili/Colombie) \\
\hline
Sécurité sociale / Retraite & CNPS (Caisse Nationale de Prévoyance Sociale) & \esp{Seguridad Social} (système contributif public) ; \esp{AFP} (privé, Chili) \\
\hline
Droit de grève & Reconnus (Art. 7 Code travail), préavis obligatoire & Droit constitutionnel (\esp{Art. 28 CE} Espagne), service minimum parfois imposé \\
\hline
Langue du droit & Français, Anglais & \esp{Espagnol} (langue officielle en Guinée équatoriale, atout pour l'emploi régional) \\
\hline
\end{tabularx}
\end{center}

\begin{savaistu}[Le droit au travail, un droit constitutionnel et universel]
La Constitution espagnole de 1978 (\esp{Constitución Española}) reconnaît explicitement le droit au travail (\esp{Art. 35}) et le droit de grève (\esp{Art. 28}). En Amérique latine, la fête du travail, \esp{el Día Internacional de los Trabajadores}, se célèbre le 1\textsuperscript{er} mai comme au Cameroun, souvent marquée par des manifestations (\esp{manifestaciones}) et des discours revendicatifs. En Guinée équatoriale, pays hispanophone voisin du Cameroun (membre de la CEMAC), l'espagnol est la langue de la législation, de l'administration et des contrats de travail : maîtriser le vocabulaire juridique espagnol (\esp{contrato indefinido, finiquito, despido, baja laboral}) est un atout professionnel majeur pour un Camerounais souhaitant travailler dans la sous-région.
\end{savaistu}

\courssub{Production orale guidée : Mini-débat}

\begin{experience}[¿Es justa la huelga? ¿Debe haber un salario mínimo?]
\textbf{Consigne :} Par groupes de 4 élèves. Préparez 2 arguments \emph{pour} et 2 arguments \emph{contre} pour chaque question. Vous \textbf{devez} utiliser les structures étudiées (\esp{tener derecho a, deber, tener que, estar obligado a, deber de}). Temps de préparation : 10 min. Présentation : 1 min/groupe.

\textbf{Expressions utiles (à adapter) :}
\begin{itemize}
\item \esp{Estoy a favor de la huelga porque los trabajadores tienen derecho a defender sus intereses.}
\item \esp{Estoy en contra de la huelga porque perjudica a la economía y a los clientes.}
\item \esp{En mi opinión, debe haber un salario mínimo para vivir dignamente.}
\item \esp{Los empresarios tienen que poder pagar los salarios sin quebrar.}
\item \esp{El gobierno está obligado a garantizar el derecho a la huelga.}
\item \esp{Debe de ser difícil encontrar un equilibrio entre derechos y productividad.}
\end{itemize}
\textbf{Évaluation :} Richesse du lexique (sindicato, convenio colectivo, patronal, piquete), correction des structures (derecho a, deber de), fluidité, interaction.
\end{experience}

\begin{aretenir}[L'essentiel pour le Bac — Fiche de révision]
\begin{itemize}
\item \textbf{Droits} : \esp{Tener derecho a} + nom / infinitif (\textbf{jamais} \esp{de}). Ex. : \esp{Tengo derecho a vacaciones.}
\item \textbf{Devoirs (obligation)} : \esp{Deber} (moral/légal), \esp{Tener que} (concret), \esp{Estar obligado a} (formel, accord du participe) + infinitif.
\item \textbf{Supposition} : \esp{Deber de} + infinitif = « il doit être... » (probabilité). \emph{Piège thème : « Il doit être fatigué » $\rightarrow$ \esp{Debe de estar cansado}.}
\item \textbf{Binôme indissociable} : \esp{Los derechos van acompañados de deberes.}
\item \textbf{Lexique « passe-partout » Bac} : \esp{salario justo, contrato indefinido/temporal, seguridad social, vacaciones pagadas, sindicato, huelga, jubilación, convenio colectivo, despido, baja laboral.}
\item \textbf{Contexte camerounais} : CNPS, SMIG, Code du travail, syndicats (CSTC, USLC), Guinée équatoriale (español lengua oficial).
\end{itemize}
\end{aretenir}

\coursec{Méthode du commentaire de texte}

Dans la section précédente, nous avons présenté les droits et les devoirs des travailleurs : le droit à un \esp{contrato}, à un \esp{sueldo} juste, à la protection d'un \esp{sindicato}, et le devoir de travailler avec honnêteté et ponctualité. Nous savons maintenant \emph{quoi} dire sur le monde du travail. Il reste à apprendre \emph{comment} le dire : c'est l'objet de cette section. À l'examen, on vous demandera souvent de \cle{commenter} un texte comme « El mundo del trabajo hoy ». Le commentaire de texte est un exercice codifié : il obéit à des règles précises que nous allons observer, énoncer puis appliquer.

\courssub{Qu'est-ce qu'un commentaire de texte ?}

\begin{definition}[Le commentaire de texte]
Le \cle{commentaire de texte} (\esp{comentario de texto}) est une \textbf{analyse organisée} qui \textbf{présente}, \textbf{explique} et \textbf{apprécie} un texte. Il ne s'agit ni de recopier le texte, ni de le résumer, ni de donner librement son avis : il s'agit de montrer que l'on a compris le texte, son vocabulaire, ses idées et le message de l'auteur, puis de porter un jugement personnel \emph{justifié}.
\end{definition}

Observons d'abord deux productions d'élèves sur le même texte :

\begin{exemplebox}[Deux débuts de commentaire : lequel est correct ?]
\textbf{Élève A :} \esp{El texto dice que hay un carpintero en Douala y una enfermera en Madrid. El carpintero trabaja mucho. La enfermera también trabaja. El texto es interesante.}

\textbf{Élève B :} \esp{Este texto es un testimonio sobre el mundo del trabajo. Presenta a dos personajes: un carpintero de Douala sin contrato y una enfermera de Madrid con un empleo estable. Analizaremos primero las ideas principales del texto, después el léxico del trabajo y finalmente el mensaje del autor.}
\end{exemplebox}

L'élève A \emph{paraphrase} : il répète le texte sans l'analyser, sans plan, sans distance. L'élève B \emph{commente} : il présente le texte (type, thème, personnages) et annonce un plan en trois parties. C'est le modèle B qu'il faut suivre.

\courssub{Les trois étapes du commentaire}

\begin{propriete}[Structure obligatoire du commentaire]
Tout commentaire de texte comporte \textbf{trois parties} :
\begin{enumerate}
\item \textbf{Introducción} : présenter le texte (type de texte, thème, auteur ou personnages) et \textbf{annoncer le plan}.
\item \textbf{Desarrollo} : développement en \textbf{2 ou 3 parties} — les idées principales, le lexique du travail utilisé, le message de l'auteur.
\item \textbf{Conclusión} : bilan de l'analyse + \textbf{ouverture} ou \textbf{opinion personnelle justifiée}.
\end{enumerate}
\end{propriete}

\begin{popfigure}
\begin{center}
\begin{tikzpicture}
\node[fill=popblueL, draw=popblue, rounded corners, text width=10.5cm, align=center, font=\small] (intro) at (0,0) {\textbf{INTRODUCCIÓN} — presentación del texto (tipo, tema, personajes) + anuncio del plan};
\node[fill=poporangeL, draw=poporange, rounded corners, text width=8.5cm, align=center, font=\small] (des) at (0,-1.8) {\textbf{DESARROLLO} — análisis en 2 o 3 partes : ideas principales, léxico del trabajo, mensaje del autor};
\node[fill=popgreenL, draw=popgreen, rounded corners, text width=6.5cm, align=center, font=\small] (conc) at (0,-3.6) {\textbf{CONCLUSIÓN} — síntesis + opinión personal justificada};
\draw[-{Stealth[length=3mm]}, thick, popblue] (intro) -- (des);
\draw[-{Stealth[length=3mm]}, thick, poporange] (des) -- (conc);
\end{tikzpicture}
\end{center}
\legende{Le commentaire en entonnoir : du général (présentation) vers l'analyse, puis vers la synthèse et l'opinion.}
\end{popfigure}

\textbf{Étape 1 — L'introduction.} Elle répond à trois questions : \emph{Quel type de texte ?} (un article, un témoignage, un dialogue, un récit), \emph{De quoi parle-t-il ?} (le thème), \emph{Qui ?} (l'auteur ou les personnages). Elle se termine par l'\textbf{annonce du plan}. Formules utiles :

\begin{itemize}
\item \esp{Este texto es un testimonio / un artículo / un diálogo sobre...} « Ce texte est un témoignage / un article / un dialogue sur... »
\item \esp{El texto trata de las diferencias entre dos trabajadores.} « Le texte traite des différences entre deux travailleurs. »
\item \esp{Presenta a dos personajes: un carpintero y una enfermera.} « Il présente deux personnages : un menuisier et une infirmière. »
\item \esp{Analizaremos primero..., después... y finalmente...} « Nous analyserons d'abord..., ensuite... et enfin... »
\end{itemize}

\textbf{Étape 2 — Le développement.} Il se construit en 2 ou 3 axes, chacun formant un paragraphe :
\begin{enumerate}
\item \textbf{Las ideas principales} : quelles idées le texte défend-il ? Exemple : \esp{El autor muestra que el trabajo dignifica, pero que los derechos no son iguales para todos.} « L'auteur montre que le travail dignifie, mais que les droits ne sont pas égaux pour tous. »
\item \textbf{El léxico del trabajo} : relever et expliquer le vocabulaire thématique (\esp{oficio, contrato, sueldo, paro, sindicato, huelga, jubilación}) et montrer à quoi il sert dans le texte.
\item \textbf{El mensaje del autor} : quelle leçon, quelle critique, quel appel le texte contient-il ?
\end{enumerate}

\textbf{Étape 3 — La conclusion.} Elle fait le \textbf{bilan} en une ou deux phrases, puis s'ouvre sur une \textbf{opinion personnelle justifiée} ou une ouverture (comparaison avec le Cameroun, question plus large). Exemple : \esp{En conclusión, este texto denuncia la desigualdad entre los trabajadores. En mi opinión, todos los trabajadores deben tener un contrato y la protección de un sindicato.} « En conclusion, ce texte dénonce l'inégalité entre les travailleurs. À mon avis, tous les travailleurs doivent avoir un contrat et la protection d'un syndicat. »

\courssub{Les connecteurs du commentaire}

Un commentaire sans connecteurs est une liste de phrases ; avec connecteurs, c'est un raisonnement. Voici les outils indispensables :

\begin{center}
\begin{tabularx}{\linewidth}{|X|X|X|}
\hline
\rowcolor{popblueL} \textbf{Español} & \textbf{Français} & \textbf{Emploi} \\
\hline
\esp{en primer lugar} & en premier lieu, d'abord & introduire le premier axe \\
\hline
\esp{además} & de plus, en outre & ajouter une idée \\
\hline
\esp{por otro lado} & d'autre part & passer à un autre aspect \\
\hline
\esp{sin embargo} & cependant, pourtant & opposer deux idées \\
\hline
\esp{por eso / por lo tanto} & c'est pourquoi / donc & conséquence \\
\hline
\esp{en conclusión} & en conclusion & ouvrir la conclusion \\
\hline
\esp{en mi opinión} & à mon avis & donner son avis \\
\hline
\esp{pienso que / creo que} & je pense que / je crois que & exprimer un jugement \\
\hline
\end{tabularx}
\end{center}

\begin{exemplebox}[Connecteurs en contexte]
\esp{En primer lugar, el texto presenta la situación del carpintero. Además, describe el trabajo de la enfermera. Por otro lado, denuncia la falta de contratos. Sin embargo, el mensaje final es optimista.} « En premier lieu, le texte présente la situation du menuisier. De plus, il décrit le travail de l'infirmière. D'autre part, il dénonce l'absence de contrats. Cependant, le message final est optimiste. »
\end{exemplebox}

\begin{attention}[Opinion : indicatif ou subjonctif ?]
\esp{En mi opinión}, \esp{pienso que} et \esp{creo que} sont suivis de l'\textbf{indicatif} : \esp{Pienso que todos los trabajadores \textbf{tienen} droits.} Mais à la \textbf{négation}, on passe au \textbf{subjonctif} : \esp{No pienso que la situación \textbf{sea} justa.} « Je ne pense pas que la situation soit juste. » Retenez : affirmation $\rightarrow$ indicatif ; négation de l'opinion $\rightarrow$ subjonctif.
\end{attention}

\courssub{Commentaire modèle sur « El mundo del trabajo hoy »}

Voici le commentaire modèle rédigé en classe sur le texte support lu et analysé dans les sections précédentes.

\begin{textebox}[Comentario modelo — « El mundo del trabajo hoy »]
Este texto es un testimonio sobre el mundo del trabajo. Presenta a dos personajes: un carpintero de Douala sin contrato y una enfermera de Madrid con un empleo estable. El autor muestra que el trabajo dignifica, pero que los derechos no son iguales para todos. En primer lugar, el texto describe la vida difícil del carpintero: trabaja mucho, pero no tiene ni contrato ni seguridad. Además, presenta la situación de la enfermera, que recibe un sueldo regular y está protegida por un sindicato. Por otro lado, el léxico del trabajo — oficio, empleo, paro, huelga, jubilación — organiza todo el texto y subraya la oposición entre los dos personajes. En conclusión, el autor denuncia la desigualdad entre los trabajadores del mundo. En mi opinión, todo trabajador debe tener un contrato y la protección de un sindicato, porque el trabajo digno es un derecho y no un privilegio.
\end{textebox}

\begin{definition}[Glossaire du commentaire modèle]
\esp{sin contrato} = sans contrat ; \esp{empleo estable} = emploi stable ; \esp{dignificar} = donner de la dignité ; \esp{seguridad} = sécurité (ici : protection sociale) ; \esp{subrayar} = souligner ; \esp{denunciar} = dénoncer ; \esp{derecho} = droit ; \esp{privilegio} = privilège.
\end{definition}

\textbf{Questions d'observation sur le modèle :}
\begin{enumerate}
\item Où commence et où finit l'introduction ? Quelle phrase annonce le plan ?
\item Relevez les quatre connecteurs du développement et indiquez leur fonction.
\item Quelle phrase exprime l'opinion personnelle ? Comment est-elle justifiée ?
\end{enumerate}

\courssub{La démarche en quatre pas}

\begin{methode}[Rédiger un commentaire de texte]
\begin{enumerate}
\item \textbf{Relever les idées} : à la lecture, soulignez les idées principales, le lexique du travail et les phrases qui révèlent le message de l'auteur.
\item \textbf{Classer en 2 ou 3 axes} : regroupez vos relevés en parties (idées / lexique / message) et rédigez une introduction qui les annonce.
\item \textbf{Rédiger avec connecteurs} : un paragraphe par axe, reliés par \esp{en primer lugar, además, por otro lado, sin embargo} ; terminez par \esp{en conclusión} et une opinion introduite par \esp{en mi opinión} et \emph{justifiée} par \esp{porque}.
\item \textbf{Relire} : vérifiez les accords (genre et nombre), les conjugaisons, les accents (\esp{opinión, jubilación, además}) et la ponctuation espagnole (¿...? ¡...!).
\end{enumerate}
\end{methode}

\begin{exoresolu}[Remettre un commentaire en ordre]
Énoncé — Voici cinq phrases d'un commentaire, dans le désordre. Remettez-les dans l'ordre logique et justifiez.
\begin{enumerate}
\item[\esp{a)}] \esp{En mi opinión, el gobierno debe proteger a los trabajadores sin contrato.}
\item[\esp{b)}] \esp{Este texto es un artículo sobre el paro y los derechos de los trabajadores.}
\item[\esp{c)}] \esp{En conclusión, el texto defiende la dignidad de todos los trabajadores.}
\item[\esp{d)}] \esp{En primer lugar, el autor describe la situación de los obreros de una fábrica de cacao.}
\item[\esp{e)}] \esp{Además, el léxico del trabajo — huelga, sindicato, sueldo — refuerza el mensaje del texto.}
\end{enumerate}
\tcblower
Solution détaillée — L'ordre correct est : \textbf{b} $\rightarrow$ \textbf{d} $\rightarrow$ \textbf{e} $\rightarrow$ \textbf{c} $\rightarrow$ \textbf{a}.
\begin{itemize}
\item \textbf{b} ouvre le commentaire : c'est la présentation du texte (type + thème), donc l'introduction.
\item \textbf{d} commence le développement, signalé par \esp{en primer lugar} (premier axe : les idées).
\item \textbf{e} poursuit le développement grâce à \esp{además} (deuxième axe : le lexique).
\item \textbf{c} amorce la conclusion avec \esp{en conclusión} : c'est le bilan.
\item \textbf{a} termine par l'opinion personnelle (\esp{en mi opinión}), placée après le bilan : on donne son avis \emph{à la fin}, jamais au début.
\end{itemize}
\end{exoresolu}

\courssub{Les erreurs à éviter}

\begin{attention}[Les trois fautes qui coûtent des points]
\begin{itemize}
\item \textbf{Paraphraser le texte} : répéter ce que dit le texte phrase par phrase n'est pas un commentaire. Il faut \emph{analyser} : dire \emph{comment} et \emph{pourquoi} l'auteur écrit.
\item \textbf{Oublier l'introduction} : un commentaire qui commence directement par l'analyse est incomplet. Présentez toujours le texte et annoncez le plan.
\item \textbf{Donner son avis sans le justifier} : après \esp{en mi opinión}, ajoutez toujours une justification avec \esp{porque}. Mauvais : \esp{En mi opinión, el texto es bueno.} Correct : \esp{En mi opinión, todo trabajador debe tener un contrato, porque el contrato garantiza sus derechos.}
\end{itemize}
Ajoutons deux pièges de langue : \esp{además} (de plus) prend un accent et ne se confond pas avec \esp{demás} (les autres) ; et on dit \esp{el texto trata \textbf{de}...}, jamais \esp{trata sobre} (calque de l'anglais \emph{treats about} / français \emph{traite de}), on préfère \esp{trata de} ou \esp{versa sobre}.
\end{attention}

\begin{exercice}[À vous de commenter]
En vous inspirant du commentaire modèle, rédigez en 8 à 10 lignes un commentaire du texte « El mundo del trabajo hoy » en insistant cette fois sur la situation des travailleurs au Cameroun. Consignes : une introduction avec annonce du plan ; deux axes de développement reliés par au moins trois connecteurs différents ; une conclusion avec une opinion personnelle justifiée par \esp{porque}.
\end{exercice}

\begin{corrige}[Exercice — éléments de correction]
Un bon commentaire contiendra : (1) une introduction du type \esp{Este texto es un testimonio sobre el mundo del trabajo en Camerún y en España. Analizaremos las ideas principales y el mensaje del autor.} ; (2) un premier axe sur les idées (\esp{En primer lugar, el texto muestra que muchos trabajadores cameruneses no tienen contrato...}) ; (3) un deuxième axe sur le lexique ou le message (\esp{Además, el léxico — oficio, sueldo, sindicato, paro — refuerza la denuncia...}) ; (4) une conclusion avec bilan et opinion justifiée (\esp{En conclusión... En mi opinión, el Estado debe controlar los contratos, porque sin contrato no hay derechos.}). Barème indicatif : introduction 2 pts, axes et connecteurs 4 pts, conclusion et opinion justifiée 2 pts, correction de la langue 2 pts.
\end{corrige}

\begin{aretenir}[L'essentiel de la méthode]
Le commentaire de texte suit toujours le même chemin : \textbf{introducción} (présentation + plan), \textbf{desarrollo} (2 ou 3 axes : idées, lexique, message), \textbf{conclusión} (bilan + opinion justifiée). Les connecteurs (\esp{en primer lugar, además, por otro lado, sin embargo, en conclusión, en mi opinión}) donnent sa colonne vertébrale au raisonnement. Trois interdits : paraphraser, sauter l'introduction, opiner sans justifier. Et un réflexe grammatical : \esp{pienso que} + indicatif, mais \esp{no pienso que} + subjonctif.
\end{aretenir}

\coursec{Production guidée}

Après avoir appris à \emph{lire} et à \emph{commenter} un texte sur le monde du travail, il est temps de passer à l'étape la plus personnelle de la leçon : \emph{produire}. Le commentaire de texte nous a appris à analyser les idées des autres ; la production guidée va nous apprendre à exprimer les nôtres, en espagnol correct, avec le lexique des métiers et des droits des travailleurs que nous avons accumulé. C'est exactement la compétence exigée au baccalauréat : rédiger un texte court, cohérent et riche en vocabulaire.

\courssub{Un modèle avant de produire : observation}

Avant d'écrire, observons comment un élève peut parler de son projet professionnel. Lis le texte suivant, rédigé par Amina, élève de Terminale à Yaoundé.

\begin{textebox}[El proyecto profesional de Amina]
Mi nombre es Amina y vivo en Yaoundé. Desde pequeña, siempre he soñado con tener una profesión útil para mi país. Quiero ser enfermera porque me gusta ayudar a la gente y cuidar a los enfermos. Mi madre trabaja en un hospital de Douala y ella me ha contado muchas historias sobre su oficio.

Sé que el mundo del trabajo no es fácil. En Camerún, como en muchos países, hay mucho paro y muchos jóvenes no encuentran empleo después de sus estudios. Por eso, quiero prepararme bien : estudiar en una buena escuela de enfermería y aprender idiomas, porque el español puede ayudarme a trabajar incluso en Guinea Ecuatorial.

Cuando termine mis estudios, espero encontrar un empleo estable. Espero firmar un buen contrato con un hospital público o privado, tener un sueldo digno y cotizar para mi jubilación. También quiero conocer mis derechos como trabajadora : el derecho a un salario justo, a vacaciones pagadas y a afiliarme a un sindicato si es necesario. Sé que también tendré deberes : ser puntual, respetar a los pacientes y trabajar con seriedad.

Mi padre siempre dice que un oficio bien hecho es una fuente de orgullo. Estoy de acuerdo con él. No quiero solo ganar dinero : quiero servir a mi comunidad. Si un día hay una huelga de enfermeros para reclamar mejores condiciones de trabajo, participaré, porque los derechos de los trabajadores hay que defenderlos.

En conclusión, mi proyecto profesional es claro : estudiar, trabajar con dignidad y ayudar a los demás. El trabajo es un derecho, pero también un deber para construir un país mejor.
\end{textebox}

\textbf{Glossaire :}

\begin{definition}[Vocabulaire du texte]
\begin{itemize}
\item \esp{enfermera} : infirmière ; \esp{cuidar a los enfermos} : soigner les malades.
\item \esp{paro} : chômage ; \esp{empleo estable} : emploi stable.
\item \esp{sueldo digno} : salaire décent ; \esp{cotizar} : cotiser (verser des cotisations sociales).
\item \esp{vacaciones pagadas} : congés payés ; \esp{afiliarse a un sindicato} : adhérer à un syndicat.
\item \esp{ser puntual} : être ponctuel ; \esp{una fuente de orgullo} : une source de fierté.
\item \esp{reclamar} : réclamer, revendiquer ; \esp{defender los derechos} : défendre les droits.
\end{itemize}
\end{definition}

\textbf{Questions d'observation :}

\begin{enumerate}
\item Combien de mots du lexique de la leçon (\esp{oficio, profesión, empleo, paro, sueldo, contrato, sindicato, huelga, jubilación}) retrouves-tu dans ce texte ? Souligne-les.
\item Comment Amina annonce-t-elle son projet ? Quelle expression utilise-t-elle pour le justifier ?
\item Quels droits et quels devoirs mentionne-t-elle ?
\end{enumerate}

\begin{corrige}[Questions d'observation]
\begin{enumerate}
\item On retrouve les neuf mots du lexique : \esp{profesión, oficio, paro, empleo, contrato, sueldo, jubilación, sindicato, huelga}.
\item Elle annonce son projet avec \esp{Quiero ser enfermera} et le justifie avec \esp{porque me gusta ayudar a la gente}.
\item Droits : salaire juste (\esp{salario justo}), congés payés (\esp{vacaciones pagadas}), adhésion à un syndicat (\esp{afiliarse a un sindicato}). Devoirs : ponctualité (\esp{ser puntual}), respect des patients (\esp{respetar a los pacientes}), sérieux au travail (\esp{trabajar con seriedad}).
\end{enumerate}
\end{corrige}

\courssub{La méthode de production écrite en cinq étapes}

Comment Amina a-t-elle construit son texte ? Elle n'a pas écrit au hasard : elle a suivi une démarche que tu vas reproduire.

\begin{methode}[Écrire un paragraphe « Mi proyecto profesional »]
\begin{enumerate}
\item \textbf{Idea} : choisis ton projet (un métier précis) et note une idée directrice : pourquoi ce métier ?
\item \textbf{Vocabulario} : écris d'abord \cle{5 mots-clés} du lexique de la leçon (par exemple : \esp{profesión, empleo, sueldo, contrato, jubilación}).
\item \textbf{Frases} : construis \cle{une phrase autour de chaque mot-clé}. Ne cherche pas encore la perfection : une phrase simple et correcte par mot.
\item \textbf{Párrafo} : relie tes phrases avec des connecteurs (\esp{porque, también, por eso, además, en conclusión}) pour former un paragraphe cohérent.
\item \textbf{Revisión} : relis-toi en vérifiant trois points : accords (genre et nombre), conjugaison, présence des mots du lexique.
\end{enumerate}
\end{methode}

\begin{popfigure}
\begin{center}
\begin{tikzpicture}[node distance=0.55cm,
  etape/.style={rectangle, rounded corners=4pt, draw=popblue, fill=popblueL, thick, minimum width=2.6cm, minimum height=0.9cm, align=center, font=\small\bfseries},
  fleche/.style={-{Stealth[length=3mm]}, thick, poporange}]
\node[etape] (a) {Idea};
\node[etape, right=of a] (b) {Vocabulario};
\node[etape, right=of b] (c) {Frases};
\node[etape, right=of c] (d) {Párrafo};
\node[etape, right=of d] (e) {Revisión};
\draw[fleche] (a) -- (b);
\draw[fleche] (b) -- (c);
\draw[fleche] (c) -- (d);
\draw[fleche] (d) -- (e);
\node[below=0.35cm of c, font=\itshape\small, text=popmuted, align=center] {De l'idée au texte corrigé : cinq étapes obligatoires};
\end{tikzpicture}
\end{center}
\legende{Organigramme de la production écrite : de l'idée au paragraphe relu.}
\end{popfigure}

\courssub{Grille d'aide à l'expression}

Pour construire tes phrases, utilise cette grille en trois temps : \cle{débuter}, \cle{justifier}, \cle{évoquer les conditions}.

\begin{center}
\begin{tabularx}{\linewidth}{|l|X|X|}
\hline
\rowcolor{popblueL} \textbf{Étape} & \textbf{Español} & \textbf{Français} \\
\hline
Débuter & \esp{Quiero ser médico / médica.} & Je veux devenir médecin. \\
\hline
 & \esp{Mi proyecto profesional es ser maestra.} & Mon projet professionnel est de devenir institutrice. \\
\hline
Justifier & \esp{porque me gusta ayudar a la gente.} & parce que j'aime aider les gens. \\
\hline
 & \esp{porque es un oficio útil para mi país.} & parce que c'est un métier utile à mon pays. \\
\hline
Conditions & \esp{Espero tener un buen sueldo y un contrato estable.} & J'espère avoir un bon salaire et un contrat stable. \\
\hline
 & \esp{Espero firmar un buen contrato.} & J'espère signer un bon contrat. \\
\hline
Droits et devoirs & \esp{Quiero conocer mis derechos y cumplir mis deberes.} & Je veux connaître mes droits et remplir mes devoirs. \\
\hline
Conclure & \esp{En conclusión, el trabajo es un derecho y un deber.} & En conclusion, le travail est un droit et un devoir. \\
\hline
\end{tabularx}
\end{center}

\begin{exemplebox}[Phrases modèles traduites]
\begin{itemize}
\item \esp{Quiero ser médico porque me gusta ayudar a la gente.} « Je veux devenir médecin parce que j'aime aider les gens. »
\item \esp{Espero firmar un buen contrato.} « J'espère signer un bon contrat. »
\item \esp{Mi tío tiene un oficio muy respetado : es carpintero.} « Mon oncle a un métier très respecté : il est menuisier. »
\item \esp{Después de trabajar cuarenta años, mi abuelo disfruta de su jubilación.} « Après avoir travaillé quarante ans, mon grand-père profite de sa retraite. »
\item \esp{Los trabajadores se afiliaron a un sindicato para defender sus derechos.} « Les travailleurs ont adhéré à un syndicat pour défendre leurs droits. »
\end{itemize}
\end{exemplebox}

\begin{attention}[Piège : pas d'article après « quiero ser »]
En français on dit « je veux devenir \emph{un} médecin », mais en espagnol, après \esp{ser} (ou \esp{quiero ser}), on ne met \textbf{pas d'article} devant le nom de métier :

\esp{Quiero ser médico.} (et non \emph{Quiero ser un médico})

\esp{Ella es profesora.} « Elle est professeur. »

\textbf{Exception} : on met l'article si le métier est \emph{qualifié} par un adjectif ou une précision : \esp{Quiero ser un médico famoso.} « Je veux devenir un médecin célèbre. » ; \esp{Es una profesora excelente.} « C'est un excellent professeur. »
\end{attention}

\begin{exoresolu}[Construire un paragraphe à partir de mots-clés]
Énoncé : à partir des cinq mots-clés \esp{profesión, empleo, sueldo, contrato, jubilación}, écris une phrase pour chacun, puis assemble-les en un paragraphe « Mi proyecto profesional ».
\tcblower
Solution détaillée :

\textbf{Étape 1 — une phrase par mot-clé :}
\begin{itemize}
\item \esp{profesión} : \esp{Quiero tener una profesión útil : ser ingeniero.} « Je veux avoir une profession utile : être ingénieur. »
\item \esp{empleo} : \esp{En mi país no es fácil encontrar empleo.} « Dans mon pays, il n'est pas facile de trouver un emploi. »
\item \esp{sueldo} : \esp{Espero tener un buen sueldo para ayudar a mi familia.} « J'espère avoir un bon salaire pour aider ma famille. »
\item \esp{contrato} : \esp{Quiero firmar un contrato estable en una empresa seria.} « Je veux signer un contrat stable dans une entreprise sérieuse. »
\item \esp{jubilación} : \esp{También quiero cotizar para mi jubilación.} « Je veux aussi cotiser pour ma retraite. »
\end{itemize}

\textbf{Étape 2 — le paragraphe assemblé avec connecteurs :}

\esp{Quiero tener una profesión útil : ser ingeniero, porque me gustan las matemáticas y la construcción. Sé que en mi país no es fácil encontrar empleo, por eso quiero estudiar mucho. Espero tener un buen sueldo para ayudar a mi familia y firmar un contrato estable en una empresa seria. Además, quiero cotizar para mi jubilación. En conclusión, mi proyecto es trabajar con dignidad.}

« Je veux avoir une profession utile : être ingénieur, parce que j'aime les mathématiques et la construction. Je sais que dans mon pays il n'est pas facile de trouver un emploi, c'est pourquoi je veux beaucoup étudier. J'espère avoir un bon salaire pour aider ma famille et signer un contrat stable dans une entreprise sérieuse. De plus, je veux cotiser pour ma retraite. En conclusion, mon projet est de travailler avec dignité. »

Vérification : 5 mots du lexique présents, accords corrects (\esp{una profesión útil}, \esp{un buen sueldo}), pas d'article après \esp{ser} (\esp{ser ingeniero}).
\end{exoresolu}

\courssub{Tâche 1 : production écrite « Mi proyecto profesional »}

\begin{exercice}[À toi d'écrire]
Rédige un paragraphe de 8 à 10 lignes intitulé \esp{Mi proyecto profesional}. Tu dois :
\begin{itemize}
\item employer \textbf{au moins 6 mots} du lexique de la leçon : \esp{oficio, profesión, empleo, paro, sueldo, contrato, sindicato, huelga, jubilación} ;
\item suivre la grille : débuter (\esp{Quiero ser...}), justifier (\esp{porque me gusta...}), évoquer les conditions (\esp{Espero tener un buen sueldo y un contrato...}) ;
\item mentionner au moins un droit et un devoir du travailleur ;
\item terminer par une phrase de conclusion (\esp{En conclusión...}).
\end{itemize}
Respecte les cinq étapes de la méthode : idée, vocabulaire, phrases, paragraphe, révision.
\end{exercice}

\courssub{Tâche 2 : production orale — la entrevista}

\begin{experience}[Interview en binôme : periodista y trabajador]
Par groupes de deux : un élève joue le rôle du \cle{journaliste} (\esp{el periodista}), l'autre celui d'un \cle{travailleur} (\esp{el trabajador / la trabajadora}) : chauffeur de bendskin, vendeuse au marché de Douala, enseignant, infirmière, agriculteur de cacao...

\textbf{Questions obligatoires du journaliste :}
\begin{itemize}
\item \esp{¿Cuál es su oficio?} « Quel est votre métier ? »
\item \esp{¿Tiene contrato?} « Avez-vous un contrat ? »
\item \esp{¿Cuáles son sus derechos?} « Quels sont vos droits ? »
\item \esp{¿Y sus deberes?} « Et vos devoirs ? »
\item \esp{¿Está afiliado a un sindicato?} « Êtes-vous affilié à un syndicat ? »
\end{itemize}

\textbf{Exemples de réponses du travailleur :}
\begin{itemize}
\item \esp{Mi oficio es conductor de mototaxi.} « Mon métier est conducteur de moto-taxi. »
\item \esp{No tengo contrato, pero espero firmar uno pronto.} « Je n'ai pas de contrat, mais j'espère en signer un bientôt. »
\item \esp{Tengo derecho a un sueldo justo y a vacaciones.} « J'ai droit à un salaire juste et à des congés. »
\item \esp{Mi deber es ser puntual y respetar a los clientes.} « Mon devoir est d'être ponctuel et de respecter les clients. »
\end{itemize}

Puis \cle{échangez les rôles}. Le professeur écoute et note les erreurs de lexique et de prononciation pour la correction collective.
\end{experience}

\courssub{Relecture collective et auto-évaluation}

\textbf{Relecture collective.} Le professeur choisit deux ou trois paragraphes d'élèves et les corrige avec toute la classe au tableau. On vérifie en priorité :

\begin{center}
\begin{tabularx}{\linewidth}{|X|X|X|}
\hline
\rowcolor{popblueL} \textbf{Point vérifié} & \textbf{Erreur fréquente} & \textbf{Correction} \\
\hline
Article après \esp{ser} + métier & \emph{Quiero ser un médico} & \esp{Quiero ser médico} \\
\hline
Accord de l'adjectif & \emph{un sueldo buena} & \esp{un buen sueldo} \\
\hline
Faux ami & \emph{paro} compris comme « arrêt » & \esp{paro} = chômage ; « arrêt » se dit \esp{parada} ou \esp{alto} \\
\hline
Après \esp{esperar que} & \emph{espero que tengo} & \esp{espero tener} (infinitif) ou \esp{espero que tenga} (subjonctif) \\
\hline
Lexique & moins de 6 mots de la liste & ajouter \esp{contrato, jubilación, sindicato}... \\
\hline
\end{tabularx}
\end{center}

\begin{attention}[Deux pièges de vocabulaire à corriger en priorité]
\begin{itemize}
\item \esp{el paro} signifie « le chômage » (faux ami : « l'arrêt » se dit \esp{la parada} ou \esp{el alto}). \esp{Hay mucho paro entre los jóvenes.} « Il y a beaucoup de chômage chez les jeunes. »
\item \esp{jubilación} = retraite (pas « jubilation » !). \esp{Mi abuelo está jubilado.} « Mon grand-père est à la retraite. »
\end{itemize}
\end{attention}

\textbf{Auto-évaluation.} Coche chaque objectif atteint :

\begin{itemize}
\item[$\square$] Je sais lire et comprendre un texte sur le monde du travail.
\item[$\square$] J'emploie correctement au moins 6 mots du lexique : \esp{oficio, profesión, empleo, paro, sueldo, contrato, sindicato, huelga, jubilación}.
\item[$\square$] Je présente mon projet professionnel avec \esp{Quiero ser...} sans article devant le métier.
\item[$\square$] Je sais citer au moins deux droits et deux devoirs des travailleurs.
\item[$\square$] Je sais mener une petite interview en espagnol sur le travail.
\item[$\square$] Je relis ma production en vérifiant accords, conjugaison et vocabulaire.
\end{itemize}

\begin{savaistu}[Comment est notée la production écrite au Bac ?]
Au baccalauréat camerounais, la production écrite d'espagnol est évaluée selon \cle{trois critères} : la \textbf{richesse du vocabulaire} (as-tu utilisé les mots de la leçon ?), la \textbf{correction grammaticale} (accords, conjugaisons, articles) et la \textbf{cohérence} (les idées sont-elles liées par des connecteurs ?). Avant de rendre ta copie, vérifie toujours ces trois points : c'est la « revisión » de notre organigramme, et elle peut te faire gagner plusieurs points !
\end{savaistu}

\begin{exercice}[Pour aller plus loin]
Transforme ton paragraphe « Mi proyecto profesional » en une petite lettre adressée à ton correspondant espagnol : commence par \esp{Querido amigo / Querida amiga}, présente ton projet en 10 lignes, et termine par \esp{Un abrazo} suivi de ton prénom. Souligne les 6 mots du lexique utilisés.
\end{exercice}

\newpage

\coursec{Activité d'intégration}

\coursec{Leçon E20 : Lectura y comentario — Oficios, mundo del trabajo, derechos y deberes}

\courssub{Texte support : « El mundo del trabajo hoy »}

\begin{textebox}[El mundo del trabajo hoy]
En Camerún, como en gran parte del mundo hispanohablante, el mundo del trabajo está en plena transformación. En Yaoundé y Douala, miles de jóvenes buscan su primer \esp{empleo} tras terminar sus estudios; algunos logran un \esp{contrato} fijo, otros aceptan trabajos precarios sin \esp{seguridad social}. El \esp{paro} juvenil sigue siendo un reto mayor: según la OIT, uno de cada cuatro jóvenes cameruneses está en \esp{paro} o subempleado.\par
\medskip
En el sector informal —vendedores en el mercado Mokolo, conductores de \esp{bendskin}, artesanos en Mvog-Mbi— la noción de \esp{derechos} (vacaciones pagadas, \esp{jubilación}, indemnización) a menudo es teórica. Sin embargo, los \esp{sindicatos} (como la CSTC o la USLC) organizan \esp{huelgas} para reivindicar un \esp{sueldo} digno y el respeto del \esp{horario}.\par
\medskip
Al otro lado del Atlántico, en España o en México, la \esp{profesión} de enfermero, profesor o ingeniero está regulada por convenios colectivos. La \esp{huelga} es un derecho constitucional. En Guinea Ecuatorial, único país africano de lengua oficial española, el Código de Trabajo garantiza la \esp{jubilación} a los 60 años y prohíbe el despido injustificado.\par
\medskip
En todas partes, el trabajador tiene \esp{deberes}: puntualidad, lealtad, respeto de las normas de seguridad. Y \esp{derechos}: a un \esp{sueldo} justo, a la formación, a la no discriminación, a sindicarse. Conocerlos es el primer paso para defenderlos.
\end{textebox}

\emph{Traduction :} « Le monde du travail aujourd'hui. Au Cameroun, comme dans une grande partie du monde hispanophone, le monde du travail est en pleine transformation. À Yaoundé et Douala, des milliers de jeunes cherchent leur premier emploi après leurs études ; certains obtiennent un contrat à durée indéterminée, d'autres acceptent des emplois précaires sans sécurité sociale. Le chômage des jeunes reste un défi majeur : selon l'OIT, un jeune Camerounais sur quatre est au chômage ou sous-employé. Dans le secteur informel — vendeurs au marché Mokolo, conducteurs de bendskin, artisans à Mvog-Mbi — la notion de droits (congés payés, retraite, indemnités) est souvent théorique. Cependant, les syndicats (comme la CSTC ou l'USLC) organisent des grèves pour réclamer un salaire décent et le respect des horaires. De l'autre côté de l'Atlantique, en Espagne ou au Mexique, la profession d'infirmier, professeur ou ingénieur est régie par des conventions collectives. La grève est un droit constitutionnel. En Guinée équatoriale, seul pays africain de langue officielle espagnole, le Code du travail garantit la retraite à 60 ans et interdit le licenciement abusif. Partout, le travailleur a des devoirs : ponctualité, loyauté, respect des normes de sécurité. Et des droits : à un salaire juste, à la formation, à la non-discrimination, à se syndiquer. Les connaître est la première étape pour les défendre. »

\courssub{Compréhension du texte}

\begin{exercice}[Questions de compréhension]
Répondez en espagnol, par des phrases complètes.
\begin{enumerate}
\item ¿Cuál es la situación del \esp{paro} juvenil en Camerún según el texto?
\item ¿Qué papel juegan los \esp{sindicatos} en el sector informal?
\item ¿Qué garantiza el Código de Trabajo en Guinea Ecuatorial respecto a la \esp{jubilación}?
\item Cite dos \esp{derechos} y dos \esp{deberes} del trabajador mencionados en el último párrafo.
\item Explique con sus palabras la diferencia entre \esp{empleo} formal e informal a partir del texto.
\end{enumerate}
\end{exercice}

\begin{corrige}[Exercice n°1]
\begin{enumerate}
\item Según el texto, uno de cada cuatro jóvenes cameruneses está en paro o subempleado.
\item Los sindicatos organizan huelgas para reivindicar un sueldo digno y el respeto del horario.
\item Garantiza la jubilación a los 60 años y prohíbe el despido injustificado.
\item Derechos: a un sueldo justo, a la formación, a la no discriminación, a sindicarse. Deberes: puntualidad, lealtad, respeto de las normas de seguridad.
\item El empleo formal tiene contrato, seguridad social, derechos garantizados (vacaciones, jubilación). El empleo informal (vendedores, bendskin, artesanos) carece a menudo de contrato y de protección social; los derechos son teóricos.
\end{enumerate}
\end{corrige}

\courssub{Lexique thématique : el mundo del trabajo}

\begin{definition}[Vocabulaire de base — Mots du programme officiel]
\begin{center}
\begin{tabularx}{\linewidth}{|X|X|}\hline
\rowcolor{popblueL} Español & Français \\ \hline
\esp{el oficio} & le métier (souvent manuel/artisanal) \\ \hline
\esp{la profesión} & la profession (exige une formation diplômante) \\ \hline
\esp{el empleo / el trabajo} & l'emploi / le travail \\ \hline
\esp{el paro / el desempleo} & le chômage \\ \hline
\esp{el sueldo / el salario} & le salaire \\ \hline
\esp{el contrato de trabajo} & le contrat de travail \\ \hline
\esp{el sindicato} & le syndicat \\ \hline
\esp{la huelga / estar en huelga} & la grève / être en grève \\ \hline
\esp{la jubilación / jubilarse} & la retraite / prendre sa retraite \\ \hline
\esp{los derechos / los deberes} & les droits / les devoirs \\ \hline
\esp{el horario / las vacaciones} & l'horaire / les congés \\ \hline
\esp{la seguridad social} & la sécurité sociale \\ \hline
\esp{el convenio colectivo} & la convention collective \\ \hline
\esp{el despido / despedir} & le licenciement / licencier \\ \hline
\esp{la indemnización} & l'indemnité (de licenciement) \\ \hline
\end{tabularx}
\end{center}
\end{definition}

\begin{attention}[Faux amis et pièges lexicaux]
\begin{itemize}
\item \esp{El paro} = le chômage (nom). Ne pas confondre avec le verbe \esp{parar} (s'arrêter, arrêter). « Être au chômage » = \esp{estar en paro} ou \esp{estar desempleado}.
\item \esp{El sueldo} et \esp{el salario} sont synonymes ; \esp{sueldo} est plus courant en Espagne, \esp{salario} en Amérique latine.
\item \esp{Oficio} vs \esp{profesión} : \esp{oficio} = métier manuel/artisanal (charpentier, cuisinier) ; \esp{profesión} = profession libérale ou diplômée (médecin, avocat, professeur). Un même travail peut être les deux selon le contexte.
\item Genre : \esp{\textbf{la} profesión}, \esp{\textbf{la} huelga}, \esp{\textbf{la} jubilación}, \esp{\textbf{la} seguridad social} ; mais \esp{\textbf{el} sindicato}, \esp{\textbf{el} sueldo}, \esp{\textbf{el} contrato}, \esp{\textbf{el} paro}, \esp{\textbf{el} empleo}.
\item \esp{Jubilarse} (verbe pronominal) = prendre sa retraite. Ne pas dire \esp{*retirarse} (qui signifie « se retirer » au sens général).
\end{itemize}
\end{attention}

\courssub{Droits et devoirs des travailleurs : structures grammaticales}

\begin{propriete}[Exprimer les droits : \esp{tener derecho a}]
\esp{Tener derecho a} + \textbf{nom} / \textbf{infinitif} \\
\emph{Signification :} « avoir droit à », « avoir le droit de ». \\
\emph{Construction :} toujours la préposition \textbf{a} (jamais \emph{de}). \\
\emph{Exemples :}
\begin{itemize}
\item \esp{Los trabajadores tienen derecho a un sueldo justo.} (nom) — Les travailleurs ont droit à un salaire juste.
\item \esp{Tengo derecho a afiliarme al sindicato.} (infinitif) — J'ai le droit de m'affilier au syndicat.
\item \esp{¿Tiene usted derecho a vacaciones pagadas?} — Avez-vous droit à des congés payés ?
\end{itemize}
\emph{Attention :} \esp{tener derecho a} s'accorde avec le sujet (\esp{tengo, tienes, tiene, tenemos, tenéis, tienen}).
\end{propriete}

\begin{propriete}[Exprimer les devoirs / obligations : trois structures]
\begin{center}
\begin{tabularx}{\linewidth}{|X|X|X|}\hline
\rowcolor{popblueL} Structure & Emploi & Exemple traduit \\ \hline
\esp{Tener que + infinitif} & Obligation extérieure, contrainte objective & \esp{Tengo que llegar a las 8.} (Je dois arriver à 8h.) \\ \hline
\esp{Deber + infinitif} & Obligation morale, devoir professionnel & \esp{Debo respetar a mis compañeros.} (Je dois respecter mes collègues.) \\ \hline
\esp{Estar obligado/a a + infinitif} & Obligation légale, contractuelle, forte & \esp{Estamos obligados a cumplir el horario.} (Nous sommes tenus de respecter l'horaire.) \\ \hline
\end{tabularx}
\end{center}
\emph{Remarque :} \esp{estar obligado} s'accorde en genre et nombre avec le sujet (\esp{obligado, obligada, obligados, obligadas}).
\end{propriete}

\courssub{Méthode du commentaire de texte}

\begin{methode}[Commenter un texte sur le monde du travail]
\textbf{Étape 1 — Lecture active.} Lisez deux fois : une fois pour le sens global, une fois en surlignant le lexique du travail (\esp{oficio, sueldo, huelga...}), les connecteurs logiques (\esp{sin embargo, además, en cambio}) et les temps verbaux (présent, passé, subjonctif).\par
\textbf{Étape 2 — Identification.} Notez : source (presse, littérature, discours), thème principal, thèse de l'auteur, structure du texte (introduction, arguments, conclusion).\par
\textbf{Étape 3 — Analyse linguistique.} Relevez : (a) le vocabulaire thématique ; (b) les structures de droits/devoirs (\esp{tener derecho a, deber, estar obligado a}) ; (c) les marqueurs d'opinion (\esp{es fundamental, considero que, a mi juicio}) ; (d) l'emploi du subjonctif après expressions de nécessité (\esp{es necesario que + subj.}).\par
\textbf{Étape 4 — Rédaction du commentaire (en espagnol).} Structure type :\par
\begin{enumerate}
\item \textbf{Introducción :} Présenter le document (titre, auteur, date, source, thème). Annoncer le plan.
\item \textbf{Desarrollo :} 2 ou 3 paragraphes arguments + exemples du texte. Utilisez : \esp{En primer lugar... ; Además... ; Por otro lado... ; En conclusión...}
\item \textbf{Conclusión :} Résumer l'idée force, donner un avis personnel ou ouvrir sur le contexte camerounais/hispanophone.
\end{enumerate}
\textbf{Étape 5 — Relecture.} Vérifiez : accords, accents, \esp{¿ ¡}, subjonctifs, \esp{por/para}, \esp{ser/estar}.
\end{methode}

\courssub{Production guidée : Interview radiophonique « Los oficios de nuestra ciudad »}

\courssub{Objectifs de l'activité}

À l'issue de cette activité d'intégration, l'élève sera capable de :
\begin{itemize}
\item réinvestir dans un échange oral authentique le \cle{lexique} du monde du travail étudié dans la leçon : \esp{oficio}, \esp{profesión}, \esp{empleo}, \esp{paro}, \esp{sueldo}, \esp{contrato}, \esp{sindicato}, \esp{huelga}, \esp{jubilación} ;
\item formuler des \cle{questions} ouvertes et fermées en espagnol correct (intonation, inversion, ¿...?) ;
\item exprimer les \cle{droits} et les \cle{devoirs} des travailleurs avec les structures adaptées : \esp{tener derecho a}, \esp{tener que}, \esp{deber}, \esp{estar obligado a} ;
\item écouter un échange oral pour en extraire des informations précises (compréhension de l'oral) ;
\item prendre la parole en public avec aisance, en respectant un rôle (journaliste / travailleur).
\end{itemize}

\courssub{Situation}

La radio scolaire du lycée de Yaoundé, \emph{Radio Joven}, prépare une émission spéciale pour la \esp{Fiesta del Trabajo} (le 1\ier{} mai). Le thème de l'émission : \esp{« Los oficios de nuestra ciudad »}. Chaque binôme d'élèves de Terminale A doit enregistrer une interview de trois à quatre minutes : un élève joue le rôle du journaliste, l'autre celui d'un travailleur camerounais ou hispanophone (charpentier de Mvog-Mbi, infirmière de l'hôpital central, chauffeur de taxi à Douala, professeur, vendeuse au marché Mokolo, cuisinier à Malabo...). Les meilleures interviews seront diffusées pendant la semaine culturelle du lycée.

Voici le document d'appui distribué par le professeur : la fiche de préparation de la radio.

\begin{textebox}[Radio Joven — Ficha de preparación de la entrevista]
\textbf{EMISIÓN ESPECIAL : « LOS OFICIOS DE NUESTRA CIUDAD »}\par
\medskip
Queridos alumnos de Terminale :\par
Para celebrar la Fiesta del Trabajo, nuestra radio escolar busca entrevistas auténticas con trabajadores de todos los oficios. Vuestra entrevista debe responder a estas preguntas :\par
¿Cuál es su oficio o profesión? ¿Desde cuándo trabaja usted?\par
¿Tiene un contrato de trabajo? ¿Cómo es su sueldo?\par
¿Cuáles son sus derechos como trabajador?\par
¿Cuáles son sus deberes?\par
¿Qué problemas conoce en su trabajo? ¿El paro? ¿La huelga? ¿Piensa en la jubilación?\par
\medskip
Duración : tres o cuatro minutos. ¡Buena suerte, periodistas!
\end{textebox}

\emph{Traduction du document :} « Émission spéciale : "Les métiers de notre ville". Chers élèves de Terminale : pour célébrer la Fête du travail, notre radio scolaire recherche des interviews authentiques de travailleurs de tous les métiers. Votre interview doit répondre à ces questions : Quel est votre métier ? Depuis quand travaillez-vous ? Avez-vous un contrat de travail ? Comment est votre salaire ? Quels sont vos droits en tant que travailleur ? Quels sont vos devoirs ? Quels problèmes connaissez-vous dans votre travail ? Le chômage ? La grève ? Pensez-vous à la retraite ? Durée : trois ou quatre minutes. Bonne chance, journalistes ! »

\courssub{Consignes}

\begin{enumerate}
\item \textbf{Choix des rôles et du métier.} Par binômes, choisissez qui sera le journaliste et qui sera le travailleur. Choisissez ensemble un métier précis et un lieu réaliste (Yaoundé, Douala, Bafoussam, Madrid, Malabo...). Le travailleur invente son identité : nom, âge, expérience.

\item \textbf{Préparation du questionnaire.} Le journaliste rédige au moins \textbf{six questions} en espagnol, en s'appuyant sur la fiche de Radio Joven. Utilisez les interrogatifs : \esp{¿Cuál?}, \esp{¿Desde cuándo?}, \esp{¿Cuánto?}, \esp{¿Por qué?}, \esp{¿Qué?}. Au moins deux questions doivent porter sur les droits et deux sur les devoirs.

\item \textbf{Préparation des réponses.} Le travailleur prépare ses réponses en employant \textbf{au moins huit mots du lexique de la leçon} (\esp{oficio}, \esp{profesión}, \esp{empleo}, \esp{paro}, \esp{sueldo}, \esp{contrato}, \esp{sindicato}, \esp{huelga}, \esp{jubilación}...). Il doit citer au moins \textbf{deux droits} (avec \esp{tener derecho a}) et \textbf{deux devoirs} (avec \esp{tener que} ou \esp{deber}).

\item \textbf{Répétition.} Répétez l'interview à voix haute. Attention à la prononciation, à l'intonation montante des questions et aux accents : \esp{profesión}, \esp{jubilación}, \esp{sindicato}.

\item \textbf{Représentation.} Jouez l'interview devant la classe (ou enregistrez-la). Pendant l'écoute, les autres élèves remplissent la fiche d'écoute ci-dessous : métier du personnage, deux droits cités, un devoir cité, un problème mentionné.

\item \textbf{Vote et débat.} La classe vote pour l'interview la plus riche en vocabulaire de la leçon et justifie son choix en espagnol : \esp{Voto por la entrevista de... porque usa muchas palabras del vocabulario:...}
\end{enumerate}

\begin{experience}[Fiche d'écoute — Ficha de escucha]
Pendant chaque interview, complétez en espagnol :\par
1. \esp{Oficio del trabajador : \ldots}\par
2. \esp{Un derecho citado : \ldots}\par
3. \esp{Otro derecho citado : \ldots}\par
4. \esp{Un deber citado : \ldots}\par
5. \esp{Un problema mencionado (paro, huelga, sueldo bajo...) : \ldots}\par
6. \esp{Palabras del vocabulario de la lección que he oído : \ldots}
\end{experience}

\courssub{Ressources à mobiliser}

\textbf{Lexique de la leçon (rappel synthétique).}

\begin{center}
\begin{tabularx}{\linewidth}{|X|X|}\hline
\rowcolor{popblueL} Español & Français \\ \hline
\esp{el oficio / la profesión} & le métier / la profession \\ \hline
\esp{el empleo / el trabajo} & l'emploi / le travail \\ \hline
\esp{el paro / el desempleo} & le chômage \\ \hline
\esp{el sueldo / el salario} & le salaire \\ \hline
\esp{el contrato de trabajo} & le contrat de travail \\ \hline
\esp{el sindicato} & le syndicat \\ \hline
\esp{la huelga / estar en huelga} & la grève / être en grève \\ \hline
\esp{la jubilación / jubilarse} & la retraite / prendre sa retraite \\ \hline
\esp{los derechos / los deberes} & les droits / les devoirs \\ \hline
\esp{el horario / las vacaciones} & l'horaire / les congés \\ \hline
\esp{la seguridad social} & la sécurité sociale \\ \hline
\end{tabularx}
\end{center}

\textbf{Structures utiles pour les droits et les devoirs.}
\begin{itemize}
\item \esp{Tener derecho a + infinitivo / nombre} : \esp{Los trabajadores tienen derecho a un sueldo justo.} « Les travailleurs ont droit à un salaire juste. »
\item \esp{Tener que + infinitivo} : \esp{Tengo que llegar puntual al trabajo.} « Je dois arriver à l'heure au travail. »
\item \esp{Deber + infinitivo} : \esp{Debo respetar a mis compañeros.} « Je dois respecter mes collègues. »
\item \esp{Estar obligado a + infinitivo} : \esp{Estamos obligados a cumplir el horario.} « Nous sommes obligés de respecter l'horaire. »
\end{itemize}

\textbf{Formules de l'interview.}
\begin{itemize}
\item Ouvrir : \esp{Buenos días, señor / señora..., gracias por aceptar esta entrevista.}
\item Enchaîner : \esp{Y ahora, otra pregunta...} ; \esp{Hablemos de sus derechos...}
\item Conclure : \esp{Muchas gracias por su tiempo. ¡Hasta pronto!}
\end{itemize}

\begin{attention}[Pièges à éviter — Check-list avant l'oral]
\begin{itemize}
\item Ne dites pas \esp{*tener derecho de} : la structure correcte est \esp{tener derecho \textbf{a}}.
\item \esp{El paro} signifie « le chômage » ; ne le confondez pas avec \esp{parar} (« arrêter »). « Être au chômage » se dit \esp{estar en paro} ou \esp{estar desempleado}.
\item Attention au genre : \esp{\textbf{la} profesión}, \esp{\textbf{la} huelga}, \esp{\textbf{la} jubilación}, \esp{\textbf{la} seguridad social}, mais \esp{\textbf{el} sindicato}, \esp{\textbf{el} sueldo}, \esp{\textbf{el} contrato}, \esp{\textbf{el} paro}, \esp{\textbf{el} empleo}.
\item N'oubliez pas le point d'interrogation ouvrant : ¿\esp{Cuánto gana usted?}
\item Pour le vouvoiement de l'interview, utilisez \esp{usted} : \esp{¿Tiene usted un contrato?} et non \esp{*¿Tienes tú...?} (trop familier pour une interview radiophonique).
\item \esp{Llegar puntual} $\rightarrow$ préférez \esp{llegar puntualmente} ou \esp{llegar a tiempo} (adverbe).
\end{itemize}
\end{attention}

\courssub{Proposition de production et corrigé commenté}

\begin{exoresolu}[Modèle d'interview : una enfermera de Yaoundé]
Produisez, par binômes, une interview radiophonique complète (douze à seize répliques) respectant les consignes de Radio Joven.
\tcblower
\textbf{Production modèle (enregistrée par deux élèves) :}\par
\medskip
\esp{— Buenos días, queridos oyentes de Radio Joven. Hoy entrevistamos a la señora Etoundi, enfermera en el Hospital Central de Yaoundé. Buenos días, señora Etoundi.}\par
\esp{— Buenos días, y gracias por la invitación.}\par
\esp{— ¿Cuál es su profesión y desde cuándo trabaja usted?}\par
\esp{— Soy enfermera. Trabajo en este hospital desde hace quince años. Es un oficio difícil pero muy noble.}\par
\esp{— ¿Tiene usted un contrato de trabajo?}\par
\esp{— Sí, tengo un contrato fijo. Es muy importante porque protege al trabajador.}\par
\esp{— Hablemos de su sueldo. ¿Está usted contenta?}\par
\esp{— Mi sueldo es regular, pero la vida es cara en Yaoundé. Muchos colegas piden un aumento.}\par
\esp{— ¿Cuáles son sus derechos como trabajadora?}\par
\esp{— Tengo derecho a un sueldo justo, a vacaciones pagadas y a la seguridad social. También tengo derecho a afiliarme a un sindicato.}\par
\esp{— ¿Y cuáles son sus deberes?}\par
\esp{— Tengo que llegar puntualmente, debo cuidar bien a los enfermos y estoy obligada a respetar el horario, incluso de noche.}\par
\esp{— ¿Qué problemas conoce en su trabajo?}\par
\esp{— A veces trabajamos demasiadas horas. El año pasado hubo una huelga de enfermeros para protestar contra las malas condiciones. Además, muchos jóvenes sufren el paro: no encuentran empleo después de sus estudios.}\par
\esp{— ¿Piensa usted en la jubilación?}\par
\esp{— ¡Todavía no! Pero espero jubilarme a los sesenta años con una buena pensión.}\par
\esp{— Muchas gracias, señora Etoundi, por su tiempo y sus respuestas.}\par
\esp{— Gracias a usted. ¡Hasta pronto!}\par
\medskip
\textbf{Commentaire des choix de langue :}
\begin{itemize}
\item \textbf{Lexique réinvesti} : douze mots de la leçon sont employés (\esp{profesión}, \esp{oficio}, \esp{contrato}, \esp{sueldo}, \esp{derecho}, \esp{sindicato}, \esp{deberes}, \esp{huelga}, \esp{paro}, \esp{empleo}, \esp{jubilación}, \esp{jubilarme}) : la consigne des huit mots minimum est largement remplie.
\item \textbf{Droits} : la structure \esp{tener derecho a} est utilisée quatre fois, avec un nom (\esp{a un sueldo justo}, \esp{a vacaciones pagadas}, \esp{a la seguridad social}) et avec un infinitif (\esp{a afiliarme}) : les deux constructions sont correctes.
\item \textbf{Devoirs} : variété des structures (\esp{tener que}, \esp{deber}, \esp{estar obligada a}) pour éviter les répétitions ; notez l'accord au féminin : \esp{obligad\textbf{a}} car la locutrice est une femme.
\item \textbf{Vouvoiement} : le journaliste utilise \esp{usted} (\esp{¿Tiene usted...?}, \esp{¿Está usted contenta?}), registre adapté à une interview formelle.
\item \textbf{Grammaire} : \esp{desde hace quince años} (durée qui continue, avec présent), \esp{hubo una huelga} (prétérit indéfini pour un fait passé ponctuel), \esp{espero jubilarme} (infinitif après \esp{esperar} quand le sujet est identique).
\item \textbf{Contexte camerounais} : l'Hospital Central de Yaoundé et le coût de la vie ancrent l'échange dans un cadre réaliste.
\item \textbf{Correction} : \esp{llegar puntualmente} (adverbe) remplace \esp{llegar puntual} (anglicisme/calque).
\end{itemize}
\end{exoresolu}

\courssub{Bilan de l'activité}

Cette activité d'intégration a permis de réunir tous les savoir-faire de la leçon dans une situation de communication authentique :
\begin{itemize}
\item \textbf{Compréhension de l'oral} : les auditeurs ont rempli une fiche d'écoute ciblée (métier, droits, devoirs, problèmes).
\item \textbf{Expression orale en interaction} : chaque binôme a conduit une interview complète, avec questions ouvertes, vouvoiement et formules de politesse radiophonique.
\item \textbf{Lexique} : le vocabulaire du monde du travail (\esp{oficio}, \esp{sueldo}, \esp{contrato}, \esp{sindicato}, \esp{huelga}, \esp{paro}, \esp{jubilación}) a été réemployé en contexte, ce qui en garantit la mémorisation.
\item \textbf{Grammaire en situation} : les structures d'expression des droits (\esp{tener derecho a}) et des devoirs (\esp{tener que}, \esp{deber}, \esp{estar obligado a}) ont été mobilisées spontanément.
\item \textbf{Dimension culturelle et citoyenne} : en évoquant la Fiesta del Trabajo, le chômage des jeunes ou les grèves d'infirmiers, les élèves ont pris conscience des réalités sociales du monde du travail, au Cameroun comme dans le monde hispanophone.
\end{itemize}

Le vote final, justifié en espagnol, clôt l'activité par une prise de parole argumentée : la classe valide ainsi collectivement les critères de réussite d'une production orale riche et correcte.

\begin{savaistu}[Le 1er mai dans le monde hispanophone et au Cameroun]
La \esp{Fiesta del Trabajo} (Fête du Travail) est fériée dans presque tous les pays hispanophones (Espagne, Mexique, Argentine, Colombie, Guinée équatoriale...). En Espagne, c'est jour de manifestations syndicales (\esp{manifestaciones}) organisées par les grands syndicats \esp{UGT} et \esp{CCOO}. Au Cameroun, le 1er mai est aussi férié : défilés à Yaoundé, discours du ministre du Travail, remise de médailles du travail. En Guinée équatoriale, le \esp{Día del Trabajador} donne lieu à des cérémonies officielles à Malabo et Bata. Le vocabulaire de la leçon (\esp{sindicato, huelga, contrato, jubilación}) est au cœur de ces célébrations : il nomme les droits conquis par les luttes ouvrières depuis le XIXe siècle.
\end{savaistu}

\newpage

\coursec{Méthodes et exercices résolus}

\courssub{Méthode 1 : Lire et comprendre un texte sur le monde du travail}

\begin{methode}[Comprendre un texte professionnel en 5 étapes]
\begin{enumerate}
\item \textbf{Lecture globale} : lire le texte une première fois sans dictionnaire. Repérer le type de texte (article de presse, reportage, témoignage), le thème et la thèse de l'auteur.
\item \textbf{Repérage du lexique clé} : souligner les mots du champ lexical du travail : \esp{oficio, profesión, empleo, paro, sueldo, contrato, sindicato, huelga, jubilación, jornada laboral, despido}. Ces mots guident la compréhension.
\item \textbf{Lecture analytique} : identifier le \cle{qui} (protagonistas), le \cle{quoi} (problema), le \cle{où/quand} (contexto), le \cle{pourquoi} (causas) et les \cle{solutions} proposées.
\item \textbf{Exploitation du contexte} : deviner le sens des mots inconnus grâce à la famille de mots (\esp{trabajar} $\rightarrow$ \esp{trabajador, trabajo}) et aux mots transparents (\esp{contrato} = contrat).
\item \textbf{Réponse aux questions} : répondre en \textbf{phrases complètes en espagnol}, en reprenant les termes de la question, sans recopier de longs passages du texte.
\end{enumerate}
\textbf{Structures à retenir pour répondre :} \esp{Según el texto...} (selon le texte), \esp{El autor afirma que...} (l'auteur affirme que), \esp{Se deduce que...} (on en déduit que), \esp{El problema principal es...} (le problème principal est).
\end{methode}

\begin{exoresolu}[Compréhension de texte : type Bac]
\textbf{Texte :} \esp{En muchos países africanos, los jóvenes titulados buscan un empleo fijo con un buen sueldo y un contrato estable. Sin embargo, el paro juvenil empuja a muchos a ejercer oficios informales : mecánico, costurera, taxista de moto. Estos trabajadores no tienen ni seguridad social ni derecho a la jubilación. Los sindicatos reclaman más protección para todos los trabajadores.}

\textbf{Questions :}
\begin{enumerate}
\item \esp{¿Qué buscan los jóvenes titulados ?}
\item \esp{¿Por qué ejercen muchos jóvenes oficios informales ?}
\item \esp{¿Qué derechos no tienen los trabajadores informales ?}
\item \esp{¿Qué reclaman los sindicatos ?}
\end{enumerate}
\tcblower
\textbf{Raisonnement :} chaque réponse se construit en reprenant la structure de la question et en la complétant avec l'information du texte. Attention à l'accord sujet-verbe et aux articles.

\textbf{Réponses rédigées :}
\begin{enumerate}
\item \esp{Los jóvenes titulados buscan un empleo fijo con un buen sueldo y un contrato estable.} « Les jeunes diplômés cherchent un emploi fixe avec un bon salaire et un contrat stable. » (On reprend le sujet de la question ; \esp{buscan} au présent, 3\up{e} personne du pluriel.)
\item \esp{Muchos jóvenes ejercen oficios informales porque el paro juvenil es muy alto.} « Beaucoup de jeunes exercent des métiers informels parce que le chômage des jeunes est très élevé. » (La question \esp{¿Por qué...?} exige une réponse introduite par \esp{porque}.)
\item \esp{Los trabajadores informales no tienen ni seguridad social ni derecho a la jubilación.} « Les travailleurs informels n'ont ni sécurité sociale ni droit à la retraite. » (Négation double espagnole : \esp{ni... ni...}.)
\item \esp{Los sindicatos reclaman más protección para todos los trabajadores.} « Les syndicats réclament plus de protection pour tous les travailleurs. »
\end{enumerate}
\textbf{Conclusion méthodologique :} une réponse complète, grammaticalement correcte et fidèle au texte vaut toujours plus qu'une phrase recopiée hors contexte.
\end{exoresolu}

\courssub{Méthode 2 : Employer le lexique des métiers et du travail}

\begin{methode}[Construire et utiliser le champ lexical du travail]
\begin{enumerate}
\item \textbf{Classer le vocabulaire par familles} :
\begin{itemize}
\item Personnes : \esp{el/la trabajador(a), el/la empleado(a), el/la obrero(a), el/la parado(a)} (chômeur), \esp{el/la jubilado(a)} (retraité), \esp{el/la empresario(a)} (chef d'entreprise).
\item Lieux et cadres : \esp{la empresa, la fábrica, la oficina, el taller} (atelier), \esp{el mercado}.
\item Contrat et rémunération : \esp{el contrato, el sueldo / el salario, el horario, la jornada laboral} (journée de travail).
\item Conflit et défense : \esp{el sindicato} (syndicat), \esp{la huelga} (grève), \esp{hacer huelga} (faire grève), \esp{reclamar / reivindicar} (revendiquer), \esp{el despido} (licenciement).
\end{itemize}
\item \textbf{Distinguer les mots proches} : \esp{oficio} = métier manuel appris par la pratique (menuisier, couturière) ; \esp{profesión} = profession exigeant des études (médecin, avocat) ; \esp{empleo} = poste de travail rémunéré ; \esp{paro} = chômage (Espagne) ; en Amérique latine on dit aussi \esp{desempleo}.
\item \textbf{Mémoriser les collocations} : \esp{buscar empleo} (chercher un emploi), \esp{firmar un contrato} (signer un contrat), \esp{estar en paro} (être au chômage), \esp{jubilarse} (prendre sa retraite), \esp{ganarse la vida} (gagner sa vie).
\item \textbf{Réutiliser le lexique dans des phrases personnelles} : un mot n'est acquis que s'il est employé dans une phrase complète, avec l'article et l'accord corrects.
\end{enumerate}
\end{methode}

\begin{exoresolu}[Lexique : compléter et employer]
\textbf{Énoncé :} Complétez les phrases avec le mot qui convient : \esp{huelga, sueldo, paro, sindicato, contrato, oficio}.
\begin{enumerate}
\item \esp{Mi tío aprendió el \dots{} de carpintero con su padre.}
\item \esp{Los obreros están en \dots{} porque quieren un aumento de salario.}
\item \esp{Después de la entrevista, María firmó un \dots{} de dos años.}
\item \esp{El \dots{} de los trabajadores defiende sus derechos.}
\item \esp{Con la crisis, muchos jóvenes de Douala están en el \dots{}.}
\item \esp{A fin de mes, cada empleado recibe su \dots{}.}
\end{enumerate}
\tcblower
\textbf{Raisonnement :} identifier l'indice sémantique de chaque phrase avant de choisir le mot.
\begin{enumerate}
\item \esp{Mi tío aprendió el \textbf{oficio} de carpintero con su padre.} → métier manuel transmis par la pratique = \esp{oficio} (et non \esp{profesión}, qui suppose des études).
\item \esp{Los obreros están en \textbf{huelga} porque quieren un aumento de salario.} → arrêt de travail collectif = \esp{estar en huelga}.
\item \esp{Después de la entrevista, María firmó un \textbf{contrato} de dos años.} → on « signe » un contrat : collocation \esp{firmar un contrato}.
\item \esp{El \textbf{sindicato} de los trabajadores defiende sus derechos.} → l'organisation qui défend les travailleurs.
\item \esp{Con la crisis, muchos jóvenes de Douala están en el \textbf{paro}.} → \esp{estar en (el) paro} = être au chômage.
\item \esp{A fin de mes, cada empleado recibe su \textbf{sueldo}.} → rémunération mensuelle.
\end{enumerate}
\textbf{Conclusion :} vérifier toujours la collocation (verbe + nom) : c'est elle qui confirme le choix du mot.
\end{exoresolu}

\courssub{Méthode 3 : Présenter les droits et devoirs des travailleurs}

\begin{methode}[Structurer un exposé sur droits y deberes]
\begin{enumerate}
\item \textbf{Introduction} : situer le sujet. \esp{El trabajo es un derecho fundamental, pero también implica deberes.}
\item \textbf{Énumérer les droits} avec des connecteurs : \esp{En primer lugar... Además... También... Por último...} Droits types : \esp{derecho a un sueldo justo, a un contrato, a la seguridad social, a vacaciones pagadas, a afiliarse a un sindicato, a la jubilación}.
\item \textbf{Énumérer les devoirs} : \esp{cumplir el horario, respetar las normas de la empresa, trabajar con responsabilidad, ser puntual, guardar el secreto profesional}.
\item \textbf{Employer les structures impersonnelles} : \esp{Se debe + infinitivo} (on doit), \esp{Hay que + infinitivo} (il faut), \esp{El trabajador tiene derecho a + nom/infinitivo}, \esp{Está prohibido + infinitivo}.
\item \textbf{Conclusion nuancée} : \esp{En conclusión, derechos y deberes son dos caras de la misma moneda.}
\end{enumerate}
\textbf{Attention au verbe} \esp{deber} : \esp{deber + infinitivo} = devoir (obligation) ; le nom \esp{el deber} = le devoir ; \esp{los deberes} = les devoirs (aussi « les devoirs à la maison »).
\end{methode}

\begin{exoresolu}[Thème : traduction vers l'espagnol]
\textbf{Énoncé :} Traduisez en espagnol.
\begin{enumerate}
\item Tout travailleur a droit à un salaire juste et à des vacances payées.
\item Les employés doivent respecter l'horaire de l'entreprise.
\item Au Cameroun, beaucoup de jeunes gagnent leur vie dans le secteur informel.
\item Il faut protéger les droits des travailleurs.
\end{enumerate}
\tcblower
\textbf{Solutions et justifications :}
\begin{enumerate}
\item \esp{Todo trabajador tiene derecho a un sueldo justo y a vacaciones pagadas.} → « avoir droit à » = \esp{tener derecho a} ; \esp{todo} + singulier = « tout » (chaque) ; \esp{vacaciones pagadas} : accord du participe passé au féminin pluriel.
\item \esp{Los empleados deben respetar el horario de la empresa.} → \esp{deber + infinitivo} pour l'obligation ; \esp{respetar} : verbe en -ar régulier ; attention à ne pas écrire \esp{*respectar}.
\item \esp{En Camerún, muchos jóvenes se ganan la vida en el sector informal.} → « gagner sa vie » = \esp{ganarse la vida} (verbe pronominal, \esp{se} accordé avec \esp{muchos jóvenes}) ; « Cameroun » = \esp{Camerún} avec accent sur le u.
\item \esp{Hay que proteger los derechos de los trabajadores.} → « il faut » impersonnel = \esp{hay que + infinitivo} ; \esp{proteger} : g devant e s'écrit \esp{g} (\esp{proteger}), mais \esp{protejo} à la 1\up{re} personne.
\end{enumerate}
\textbf{Conclusion :} en thème, chaque structure française doit être remplacée par sa structure espagnole authentique, jamais calquée mot à mot.
\end{exoresolu}

\courssub{Méthode 4 : Commenter un texte sur le travail}

\begin{methode}[La méthode du commentaire de texte (comentario de texto)]
\begin{enumerate}
\item \textbf{Présentation} : type de texte, thème, source. \esp{Este texto es un artículo periodístico que trata del mundo del trabajo.}
\item \textbf{Résumé bref} (3-4 lignes) : l'idée principale sans recopier. \esp{El autor denuncia... / describe... / defiende...}
\item \textbf{Analyse de la structure} : introduction, développement, conclusion du texte ; repérer la thèse et les arguments.
\item \textbf{Analyse du lexique et du style} : champ lexical dominant (\esp{empleo, paro, derechos}), ton (dénonciation, espoir, critique), procédés (chiffres, témoignages, oppositions).
\item \textbf{Opinion personnelle justifiée} : \esp{En mi opinión... / Me parece que... / Estoy de acuerdo con... porque...} C'est la partie la plus valorisée : donner son avis avec un argument et un exemple (par exemple la situation au Cameroun).
\item \textbf{Conclusion} : phrase de synthèse. \esp{En definitiva, el texto nos invita a reflexionar sobre...}
\end{enumerate}
\textbf{Plan type à retenir :} Presentación → Resumen → Análisis → Opinión → Conclusión.
\end{methode}

\begin{exoresolu}[Commentaire guidé : introduction et opinion]
\textbf{Énoncé :} Voici un extrait : \esp{La jubilación debería ser un descanso merecido después de años de trabajo, pero muchos jubilados viven con pensiones muy bajas y deben seguir trabajando para sobrevivir.} Rédigez : a) la présentation du texte (2 phrases) ; b) votre opinion personnelle (3 phrases) avec un exemple camerounais.
\tcblower
\textbf{Raisonnement :} a) identifier le type de texte et le thème ; b) donner un avis + un argument + un exemple, avec les connecteurs \esp{en mi opinión, porque, por ejemplo}.

\textbf{Réponse rédigée :}

a) \esp{Este texto es un fragmento de un artículo de opinión sobre el mundo del trabajo. El autor denuncia la difícil situación de los jubilados que reciben pensiones muy bajas.} « Ce texte est un fragment d'un article d'opinion sur le monde du travail. L'auteur dénonce la situation difficile des retraités qui reçoivent des pensions très basses. »

b) \esp{En mi opinión, es una injusticia que una persona trabaje toda su vida y no pueda descansar al final. Me parece que el Estado debe proteger más a los jubilados. Por ejemplo, en Camerún, muchos ancianos que trabajaron en el campo o en el sector informal no reciben ninguna pensión y dependen de sus hijos.} « À mon avis, c'est une injustice qu'une personne travaille toute sa vie et ne puisse pas se reposer à la fin. Il me semble que l'État doit davantage protéger les retraités. Par exemple, au Cameroun, beaucoup de personnes âgées qui ont travaillé dans les champs ou dans le secteur informel ne reçoivent aucune pension et dépendent de leurs enfants. »

\textbf{Justifications grammaticales :} \esp{es una injusticia que + subjuntivo} (\esp{trabaje, pueda}) : après une expression de jugement personnel, l'espagnol exige le subjonctif ; \esp{ninguna pensión} : adjectif indéfini négatif au féminin singulier ; \esp{ancianos} : accent non nécessaire (mot grave terminé par -s).

\textbf{Conclusion :} une opinion sans argument ni exemple ne rapporte presque rien ; la triade avis + parce que + exemple est la clé de la réussite.
\end{exoresolu}

\courssub{Méthode 5 : Version — traduire de l'espagnol vers le français}

\begin{methode}[Traduire un texte sur le travail sans calquer]
\begin{enumerate}
\item \textbf{Lire tout le passage} avant de traduire la première phrase : le sens global conditionne chaque choix.
\item \textbf{Identifier les temps espagnols} et choisir l'équivalent français : présent espagnol → présent français ; \esp{pretérito imperfecto} → imparfait ; \esp{pretérito indefinido} → passé composé ou passé simple.
\item \textbf{Se méfier des faux amis} : \esp{actualmente} = actuellement (et non « actuellement » au sens de « en réalité ») ; \esp{una entrevista de trabajo} = un entretien d'embauche ; \esp{largo} = long (et non « large ») ; \esp{asistir a} = assister à / être présent à.
\item \textbf{Traduire les structures impersonnelles} : \esp{se buscan empleados} = on recherche des employés ; \esp{hay que} = il faut.
\item \textbf{Relire en français} : la phrase traduite doit être naturelle en français, pas un mot-à-mot.
\end{enumerate}
\end{methode}

\begin{exoresolu}[Version : un passage sur le chômage]
\textbf{Énoncé :} Traduisez en français : \esp{El paro es uno de los problemas más graves de nuestra sociedad. Muchos jóvenes, después de terminar sus estudios, no encuentran trabajo y tienen que aceptar empleos mal pagados. Por eso, algunos deciden crear su propia empresa.}
\tcblower
\textbf{Analyse préalable :}
\begin{itemize}
\item \esp{el paro} : faux ami partiel — ici « le chômage », pas « l'arrêt ».
\item \esp{después de terminar} : gérondif de temps → « après avoir terminé ».
\item \esp{tienen que + infinitivo} : obligation → « doivent ».
\item \esp{mal pagados} : participe passé accordé → « mal payés ».
\item \esp{su propia empresa} : « leur propre entreprise ».
\end{itemize}
\textbf{Traduction :} « Le chômage est l'un des problèmes les plus graves de notre société. Beaucoup de jeunes, après avoir terminé leurs études, ne trouvent pas de travail et doivent accepter des emplois mal payés. C'est pourquoi certains décident de créer leur propre entreprise. »

\textbf{Justifications :} \esp{uno de los problemas más graves} : superlatif relatif → « l'un des problèmes les plus graves » ; \esp{por eso} : connecteur de conséquence → « c'est pourquoi » (et non « pour cela » qui serait lourd) ; \esp{no encuentran trabajo} : négation simple → « ne trouvent pas de travail ».

\textbf{Conclusion :} en version, la fidélité au sens prime sur la littéralité : chaque mot espagnol doit trouver son équivalent français naturel.
\end{exoresolu}

\begin{attention}[Les 5 erreurs qui coûtent des points au Bac]
\begin{enumerate}
\item \textbf{Oublier les accents} : \esp{profesión} (accent sur la ó), \esp{jubilación}, \esp{Camerún}, \esp{también}. Un mot sans accent est une faute : \esp{*profesion} est sanctionné.
\item \textbf{Les faux amis} : \esp{asistir a} = assister à (et non « aider » = \esp{ayudar}) ; \esp{actualmente} = actuellement ; \esp{una larga jornada} = une longue journée (et non « large ») ; \esp{el paro} = le chômage (et non « la panne »).
\item \textbf{Les calques du français} : « chercher un emploi » = \esp{buscar empleo} (sans \esp{por} : jamais \esp{*buscar por empleo}) ; « gagner sa vie » = \esp{ganarse la vida} (jamais \esp{*ganar su vida}) ; « être au chômage » = \esp{estar en (el) paro}.
\item \textbf{Les accords} : \esp{las trabajadoras}, \esp{vacaciones pagadas}, \esp{empleos mal pagados}. Le participe passé employé comme adjectif s'accorde toujours en genre et en nombre.
\item \textbf{Le choix du mode et du temps} : après \esp{es importante que / es una injusticia que / me parece que + négation}, il faut le \textbf{subjonctif} : \esp{Es importante que los jóvenes \textbf{encuentren} trabajo.} Et dans un récit au passé, distinguer \esp{indefinido} (action ponctuelle : \esp{firmó el contrato}) et \esp{imperfecto} (description, habitude : \esp{trabajaba en una fábrica}).
\end{enumerate}
\end{attention}

\newpage

\coursec{Exercices}

\courssub{Je vérifie mes connaissances}

\begin{exercice}[QCM : le monde du travail]
Entoure la bonne réponse.
\begin{enumerate}
\item \esp{El \cle{paro}} désigne :
\begin{enumerate}
\item une augmentation de salaire
\item la situation d'une personne sans emploi
\item une fête du travail
\end{enumerate}
\item \esp{El \cle{sueldo}} est :
\begin{enumerate}
\item le document signé à l'embauche
\item la rémunération reçue par le travailleur
\item le lieu de travail
\end{enumerate}
\item \esp{Un \cle{sindicato}} est :
\begin{enumerate}
\item une organisation qui défend les droits des travailleurs
\item une entreprise privée
\item un impôt sur le salaire
\end{enumerate}
\item \esp{La \cle{jubilación}} correspond à :
\begin{enumerate}
\item la première embauche
\item le stage professionnel
\item la cessation d'activité à la fin de la carrière
\end{enumerate}
\end{enumerate}
\end{exercice}

\begin{exercice}[Vrai ou faux ? Justifie]
Dis si chaque affirmation est \esp{verdadera} ou \esp{falsa} et justifie ta réponse en une phrase en espagnol.
\begin{enumerate}
\item \esp{Hacer huelga significa trabajar más horas de lo normal.}
\item \esp{Un contrato es un acuerdo escrito entre el trabajador y el empresario.}
\item \esp{Un oficio es siempre una actividad que exige estudios universitarios.}
\item \esp{Los trabajadores tienen derecho a un salario justo y a descansar.}
\end{enumerate}
\end{exercice}

\begin{exercice}[Vocabulaire : phrases à trous]
Complète chaque phrase avec le mot qui convient : \esp{paro -- sueldo -- contrato -- sindicato -- huelga -- jubilación}.
\begin{enumerate}
\item \esp{Mi padre cobra su \ldots\ldots\ldots a fin de mes.}
\item \esp{Los obreros están en \ldots\ldots\ldots porque piden mejores condiciones de trabajo.}
\item \esp{Después de cuarenta años de trabajo, mi abuelo disfruta de su \ldots\ldots\ldots.}
\item \esp{Antes de empezar a trabajar, firmé un \ldots\ldots\ldots de dos años.}
\item \esp{Desde que cerró la fábrica, muchos jóvenes están en el \ldots\ldots\ldots.}
\item \esp{El \ldots\ldots\ldots ha negociado un aumento salarial para todos los empleados.}
\end{enumerate}
\end{exercice}

\begin{exercice}[Conjugaison : à compléter]
Conjugue le verbe entre parenthèses au temps et à la personne indiqués par le sujet.
\begin{enumerate}
\item \esp{Los trabajadores (tener) \ldots\ldots\ldots derecho a vacaciones pagadas.}
\item \esp{Mi tío (estar) \ldots\ldots\ldots en el paro desde hace un año.}
\item \esp{El año pasado, nosotros (firmar) \ldots\ldots\ldots un contrato nuevo.}
\item \esp{Es importante que tú (conocer) \ldots\ldots\ldots tus derechos. (subjonctif présent)}
\item \esp{Cuando yo (ser) \ldots\ldots\ldots mayor, trabajaré como médico.}
\end{enumerate}
\end{exercice}

\courssub{Je m'entraîne}

\begin{exercice}[Association mot / définition]
Relie chaque mot à sa définition en espagnol.

\begin{center}
\begin{tabularx}{\linewidth}{|X|X|}
\hline
\rowcolor{popblueL} \textbf{Palabra} & \textbf{Definición} \\
\hline
\esp{1. oficio} & \esp{a) Cese colectivo del trabajo para reivindicar algo.} \\
\hline
\esp{2. profesión} & \esp{b) Actividad laboral que se aprende con la práctica, como carpintero o sastre.} \\
\hline
\esp{3. empleo} & \esp{c) Puesto de trabajo remunerado.} \\
\hline
\esp{4. jubilación} & \esp{d) Ocupación que exige estudios superiores, como abogado o profesor.} \\
\hline
\esp{5. huelga} & \esp{e) Retiro de la vida laboral al llegar a cierta edad.} \\
\hline
\end{tabularx}
\end{center}
\end{exercice}

\begin{exercice}[Traduction : thème]
Traduis en espagnol les phrases suivantes.
\begin{enumerate}
\item Les travailleurs ont le droit de se syndiquer et de faire grève.
\item Mon oncle est au chômage depuis un an.
\item Ma sœur a signé un contrat dans une banque de Douala.
\item Tous les employés reçoivent un salaire à la fin du mois.
\end{enumerate}
\end{exercice}

\begin{exercice}[Texte à trous]
Complète le texte suivant avec les mots de la liste : \esp{empleo -- paro -- contratos -- derechos -- sueldo -- formación}.

\begin{textebox}[El trabajo en mi barrio]
En mi barrio de Yaoundé, muchos jóvenes buscan \ldots\ldots\ldots pero no lo encuentran fácilmente. Algunos llevan dos años en el \ldots\ldots\ldots. Los que trabajan en las empresas de la ciudad tienen \ldots\ldots\ldots de seis meses o de un año. Mi vecina, que es costurera, gana un \ldots\ldots\ldots modesto, pero está contenta porque ejerce el oficio que le gusta. Todos los trabajadores deben conocer sus \ldots\ldots\ldots : el descanso semanal, la seguridad y un salario justo. Por eso, la \ldots\ldots\ldots profesional es tan importante para el futuro de los jóvenes cameruneses.
\end{textebox}
\end{exercice}

\begin{exercice}[Transformation de phrases]
Réécris chaque phrase en suivant la consigne entre parenthèses.
\begin{enumerate}
\item \esp{Los obreros piden un aumento de sueldo.} (passer au pluriel du sujet \esp{el obrero} et accorder tout)
\item \esp{Mi madre trabaja en un hospital.} (mettre au passé composé : \esp{pretérito perfecto})
\item \esp{El sindicato defiende a los trabajadores.} (mettre à la forme négative)
\item \esp{El joven firmó el contrato ayer.} (transformer en question avec \esp{¿Cuándo...?})
\end{enumerate}
\end{exercice}

\courssub{Je me prépare au Bac}

\begin{exercice}[Texte type Bac : compréhension et langue]
Lis le texte suivant, puis réponds aux questions.

\begin{textebox}[El paro juvenil, un reto para España]
En España, muchos jóvenes terminan sus estudios y no encuentran trabajo. El paro juvenil es uno de los problemas más graves del país: miles de chicos y chicas de veinte a treinta años buscan empleo cada día sin éxito. Algunos aceptan contratos de pocos meses, con sueldos bajos y sin seguridad. Otros deciden marcharse al extranjero para trabajar como camareros, enfermeros o ingenieros.

Sin embargo, hay motivos para la esperanza. Los sindicatos y las asociaciones de jóvenes piden más formación profesional y contratos estables. Además, muchos jóvenes crean sus propias empresas: tiendas en línea, talleres de reparación, granjas ecológicas. «No queremos limosnas, queremos oportunidades», declara una joven de Madrid. Para ella, todo trabajador tiene derecho a un sueldo digno, a vacaciones y a la jubilación después de una vida de esfuerzo.
\end{textebox}

\textbf{Questions de compréhension} (réponds en espagnol, avec des phrases complètes) :
\begin{enumerate}
\item \esp{¿Por qué es grave el problema del paro juvenil en España?}
\item \esp{¿Qué tipo de contratos aceptan algunos jóvenes?}
\item \esp{¿Qué hacen otros jóvenes para encontrar trabajo?}
\item \esp{¿Qué piden los sindicatos y las asociaciones de jóvenes?}
\item \esp{Según la joven de Madrid, ¿qué derechos tiene todo trabajador?}
\end{enumerate}

\textbf{Questions de langue} :
\begin{enumerate}
\item Trouve dans le texte un synonyme de \esp{joven} et un contraire de \esp{éxito}.
\item Conjugue \esp{encontrar} et \esp{decidir} au \esp{pretérito indefinido}, première personne du singulier.
\item Traduis en français la phrase : \esp{«No queremos limosnas, queremos oportunidades.»}
\end{enumerate}
\end{exercice}

\begin{exercice}[Thème et version]
\textbf{A. Thème.} Traduis en espagnol :
\begin{enumerate}
\item Depuis la fermeture de l'usine, beaucoup d'ouvriers sont au chômage.
\item Les employés ont fait grève hier pour demander un meilleur salaire.
\item Quand mon père prendra sa retraite, il ouvrira un petit atelier à Yaoundé.
\end{enumerate}

\textbf{B. Version.} Traduis en français :

\begin{textebox}[Los derechos del trabajador]
Todo trabajador tiene derechos, pero también deberes. Debe cumplir su horario, respetar a sus compañeros y hacer bien su trabajo. A cambio, tiene derecho a un contrato claro, a un sueldo justo, a días de descanso y a afiliarse a un sindicato. Conocer estos derechos es el primer paso para defenderlos.
\end{textebox}
\end{exercice}

\begin{exercice}[Expression écrite type Bac]
Rédige un texte en espagnol de \textbf{80 à 100 mots} sur le sujet suivant :

\esp{¿Cuál es el oficio de tus sueños? ¿Qué derechos debe tener todo trabajador?}

Consignes : présente le métier de tes rêves et explique pourquoi il t'attire ; cite ensuite au moins \textbf{trois droits} que tout travailleur doit avoir (utilise le lexique de la leçon : \esp{sueldo, contrato, descanso, sindicato, jubilación...}) ; termine par une phrase de conclusion personnelle.
\end{exercice}

\newpage

\coursec{Corrigés des exercices}

\begin{corrige}[Exercice 1 — QCM : le monde du travail]
\begin{enumerate}
\item Réponse \textbf{b)} : \esp{el paro} désigne la situation d'une personne sans emploi (le chômage). On dit aussi \esp{estar en el paro} (être au chômage).
\item Réponse \textbf{b)} : \esp{el sueldo} est la rémunération reçue par le travailleur (le salaire). Synonymes : \esp{el salario}, \esp{la paga}.
\item Réponse \textbf{a)} : \esp{un sindicato} est une organisation qui défend les droits des travailleurs (un syndicat).
\item Réponse \textbf{c)} : \esp{la jubilación} correspond à la cessation d'activité à la fin de la carrière (la retraite).
\end{enumerate}
\textbf{Point de vigilance} : ne pas confondre \esp{paro} (chômage) et \esp{huelga} (grève) ; ne pas traduire \esp{jubilación} par « jubilation » (faux ami).
\end{corrige}

\begin{corrige}[Exercice 2 — Vrai ou faux ? Justifie]
\begin{enumerate}
\item \textbf{Falsa.} \esp{Hacer huelga significa dejar de trabajar para reivindicar algo, no trabajar más horas.} (Faire grève, c'est cesser le travail, pas travailler davantage.)
\item \textbf{Verdadera.} \esp{Un contrato es un acuerdo escrito y firmado entre el trabajador y el empresario.} (Un contrat est un accord écrit entre le salarié et l'employeur.)
\item \textbf{Falsa.} \esp{Un oficio se aprende generalmente con la práctica, como carpintero o sastre; no exige estudios universitarios.} (Un métier manuel s'apprend par la pratique ; c'est la \esp{profesión} qui exige des études supérieures.)
\item \textbf{Verdadera.} \esp{Todos los trabajadores tienen derecho a un salario justo, al descanso semanal y a las vacaciones pagadas.}
\end{enumerate}
\textbf{Point de vigilance} : pour justifier, employer une phrase complète avec un verbe conjugué ; \esp{hacer huelga} se construit sans article.
\end{corrige}

\begin{corrige}[Exercice 3 — Vocabulaire : phrases à trous]
\begin{enumerate}
\item \esp{Mi padre cobra su \textbf{sueldo} a fin de mes.} (Mon père touche son salaire à la fin du mois.)
\item \esp{Los obreros están en \textbf{huelga} porque piden mejores condiciones de trabajo.}
\item \esp{Después de cuarenta años de trabajo, mi abuelo disfruta de su \textbf{jubilación}.}
\item \esp{Antes de empezar a trabajar, firmé un \textbf{contrato} de dos años.}
\item \esp{Desde que cerró la fábrica, muchos jóvenes están en el \textbf{paro}.}
\item \esp{El \textbf{sindicato} ha negociado un aumento salarial para todos los empleados.}
\end{enumerate}
\textbf{Points de vigilance} : on dit \esp{estar en huelga} (sans article) mais \esp{estar en el paro} (avec article) ; \esp{cobrar el sueldo} = toucher son salaire.
\end{corrige}

\begin{corrige}[Exercice 4 — Conjugaison : à compléter]
\begin{enumerate}
\item \esp{Los trabajadores \textbf{tienen} derecho a vacaciones pagadas.} — présent de l'indicatif, 3\up{e} personne du pluriel de \esp{tener} (irrégulier : radical \esp{tien-}).
\item \esp{Mi tío \textbf{está} en el paro desde hace un año.} — présent, 3\up{e} personne du singulier de \esp{estar} ; accent obligatoire sur \esp{está}.
\item \esp{El año pasado, nosotros \textbf{firmamos} un contrato nuevo.} — \esp{pretérito indefinido} ; attention, la forme \esp{firmamos} est identique au présent : c'est le complément de temps \esp{el año pasado} qui impose le passé.
\item \esp{Es importante que tú \textbf{conozcas} tus derechos.} — subjonctif présent après \esp{es importante que} ; \esp{conocer} : 1\up{re} personne \esp{conozco}, d'où subjonctif \esp{conozca, conozcas, conozca...}.
\item \esp{Cuando yo \textbf{sea} mayor, trabajaré como médico.} — subjonctif présent de \esp{ser} après \esp{cuando} suivi d'un futur (action non encore réalisée).
\end{enumerate}
\textbf{Points de vigilance} : après \esp{cuando} + futur en français, l'espagnol exige le subjonctif ; ne jamais écrire \esp{cuando seré mayor}.
\end{corrige}

\begin{corrige}[Exercice 5 — Association mot / définition]
\begin{center}
\begin{tabularx}{\linewidth}{|X|X|}
\hline
\rowcolor{popblueL} \textbf{Palabra} & \textbf{Definición} \\
\hline
\esp{1. oficio} & \esp{b) Actividad laboral que se aprende con la práctica.} \\
\hline
\esp{2. profesión} & \esp{d) Ocupación que exige estudios superiores.} \\
\hline
\esp{3. empleo} & \esp{c) Puesto de trabajo remunerado.} \\
\hline
\esp{4. jubilación} & \esp{e) Retiro de la vida laboral al llegar a cierta edad.} \\
\hline
\esp{5. huelga} & \esp{a) Cese colectivo del trabajo para reivindicar algo.} \\
\hline
\end{tabularx}
\end{center}
Réponses : \textbf{1-b ; 2-d ; 3-c ; 4-e ; 5-a.}
\textbf{Point de vigilance} : la distinction \esp{oficio} / \esp{profesión} est classique au Bac : l'\esp{oficio} s'apprend par la pratique (menuisier, tailleur), la \esp{profesión} exige des études supérieures (avocat, médecin).
\end{corrige}

\begin{corrige}[Exercice 6 — Traduction : thème]
\begin{enumerate}
\item \esp{Los trabajadores tienen (el) derecho de sindicarse y de hacer huelga.} — \esp{tener derecho a / de} + infinitif ; \esp{sindicarse} = se syndiquer (verbe pronominal).
\item \esp{Mi tío está en el paro desde hace un año.} — \esp{estar} (et non \esp{ser}) pour une situation ; \esp{desde hace un año} = depuis un an.
\item \esp{Mi hermana ha firmado un contrato en un banco de Duala.} — \esp{pretérito perfecto} pour une action récente liée au présent ; on accepte aussi \esp{firmó} (indefinido). « Banque » se traduit ici \esp{banco} (établissement financier).
\item \esp{Todos los empleados reciben (cobran) un sueldo a fin de mes.} — \esp{a fin de mes} sans article ; \esp{cobrar} est plus idiomatique que \esp{recibir} pour un salaire.
\end{enumerate}
\textbf{Points de vigilance} : « être au chômage » = \esp{estar en el paro} (jamais \esp{ser}) ; « faire grève » = \esp{hacer huelga} ; ne pas calquer « depuis un an » par \esp{desde un año}.
\end{corrige}

\begin{corrige}[Exercice 7 — Texte à trous]
\begin{enumerate}
\item \esp{...muchos jóvenes buscan \textbf{empleo} pero no lo encuentran fácilmente.}
\item \esp{Algunos llevan dos años en el \textbf{paro}.}
\item \esp{...tienen \textbf{contratos} de seis meses o de un año.}
\item \esp{...gana un \textbf{sueldo} modesto...}
\item \esp{Todos los trabajadores deben conocer sus \textbf{derechos}...}
\item \esp{...la \textbf{formación} profesional es tan importante...}
\end{enumerate}
\textbf{Justifications} : \esp{buscar empleo} (chercher du travail) ; \esp{en el paro} avec article ; \esp{contratos} au pluriel (accord avec \esp{de seis meses o de un año}) ; \esp{ganar un sueldo} (gagner un salaire) ; \esp{conocer sus derechos} (connaître ses droits) ; \esp{la formación profesional} (la formation professionnelle).
\end{corrige}

\begin{corrige}[Exercice 8 — Transformation de phrases]
\begin{enumerate}
\item \esp{El obrero pide un aumento de sueldo.} — au singulier : \esp{los obreros} devient \esp{el obrero}, et le verbe s'accorde : \esp{piden} $\rightarrow$ \esp{pide}.
\item \esp{Mi madre ha trabajado en un hospital.} — \esp{pretérito perfecto} : auxiliaire \esp{haber} au présent (\esp{ha}) + participe passé \esp{trabajado} (invariable).
\item \esp{El sindicato no defiende a los trabajadores.} — la négation se fait avec \esp{no} placé devant le verbe conjugué.
\item \esp{¿Cuándo firmó el joven el contrato?} — question introduite par \esp{¿Cuándo...?} avec inversion du sujet ; ne pas oublier le point d'interrogation ouvrant \esp{¿}.
\end{enumerate}
\textbf{Points de vigilance} : le participe passé employé avec \esp{haber} ne s'accorde jamais (\esp{mi madre ha trabajado}, sans accord) ; en espagnol, toute question s'ouvre par \esp{¿}.
\end{corrige}

\begin{corrige}[Exercice 9 — Texte type Bac : compréhension et langue]
\textbf{Questions de compréhension} (réponses modèles, en phrases complètes) :
\begin{enumerate}
\item \esp{Es grave porque miles de jóvenes de veinte a treinta años terminan sus estudios y buscan empleo cada día sin éxito.}
\item \esp{Algunos jóvenes aceptan contratos de pocos meses, con sueldos bajos y sin seguridad.}
\item \esp{Otros jóvenes deciden marcharse al extranjero para trabajar como camareros, enfermeros o ingenieros.}
\item \esp{Piden más formación profesional y contratos estables.}
\item \esp{Según ella, todo trabajador tiene derecho a un sueldo digno, a vacaciones y a la jubilación después de una vida de esfuerzo.}
\end{enumerate}
\textbf{Questions de langue} :
\begin{enumerate}
\item Synonyme de \esp{joven} : \esp{chico / chica} ; contraire de \esp{éxito} : \esp{sin éxito} (l'antonyme nominal est \esp{fracaso}).
\item \esp{Encontrar} : \esp{yo encontré} ; \esp{decidir} : \esp{yo decidí}. Attention à l'accent sur la dernière syllabe : \esp{encontré}, \esp{decidí}.
\item Traduction : « Nous ne voulons pas la charité (l'aumône), nous voulons des opportunités (des chances). »
\end{enumerate}
\textbf{Barème indicatif (sur 20)} : compréhension 10 pts (2 pts par question : contenu 1 pt, correction linguistique 1 pt) ; questions de langue 6 pts (2 pts par item) ; présentation générale 4 pts.
\textbf{Points de vigilance} : répondre par des phrases complètes en reprenant les termes de la question ; ne pas recopier tout le texte ; soigner les accents (\esp{encontré}, \esp{decidí}, \esp{jubilación}) ; \esp{limosna} = aumône, charité (ne pas traduire par « limousin » !).
\end{corrige}

\begin{corrige}[Exercice 10 — Thème et version]
\textbf{A. Thème.}
\begin{enumerate}
\item \esp{Desde el cierre de la fábrica, muchos obreros están en el paro.} — « fermeture » = \esp{cierre} (nom, pas le verbe) ; \esp{estar en el paro}.
\item \esp{Los empleados hicieron huelga ayer para pedir un sueldo mejor (un mejor sueldo).} — \esp{indefinido} de \esp{hacer} : \esp{hicieron} (irrégulier) ; \esp{para} + infinitif exprime le but.
\item \esp{Cuando mi padre se jubile, abrirá un pequeño taller en Yaundé.} — subjonctif après \esp{cuando} (futur) : \esp{se jubile} ; futur simple : \esp{abrirá} ; « atelier » = \esp{taller} (faux ami : \esp{taller} n'est pas « taille »).
\end{enumerate}
\textbf{B. Version.}

« Tout travailleur a des droits, mais aussi des devoirs. Il doit respecter son horaire (ses horaires), respecter ses collègues et bien faire son travail. En échange, il a droit à un contrat clair, à un salaire juste, à des jours de repos et à s'affilier à un syndicat. Connaître ces droits est le premier pas pour les défendre. »

\textbf{Barème indicatif (sur 20)} : thème 10 pts (environ 3 pts par phrase : lexique 1 pt, grammaire 1 pt, orthographe/accents 1 pt) ; version 10 pts (fidélité au sens 5 pts, qualité du français 5 pts).
\textbf{Points de vigilance} : \esp{a cambio} = en échange ; \esp{afiliarse a un sindicato} = s'affilier à un syndicat ; \esp{cumplir su horario} = respecter ses horaires ; en version, éviter le mot à mot et rédiger un français naturel.
\end{corrige}

\begin{corrige}[Exercice 11 — Expression écrite type Bac]
\textbf{Proposition de rédaction modèle} (environ 95 mots) :

\begin{textebox}[El oficio de mis sueños]
El oficio de mis sueños es ser carpintero, como mi abuelo. Me atrae porque me gusta trabajar con las manos y crear objetos útiles para la gente de mi barrio. Además, un buen carpintero siempre encuentra clientes en Yaoundé.

Sin embargo, pienso que todo trabajador debe tener derechos. Primero, tiene derecho a un sueldo justo y a un contrato claro. Segundo, merece días de descanso y vacaciones pagadas. Tercero, debe poder afiliarse a un sindicato para defender sus intereses. Finalmente, después de una vida de esfuerzo, tiene derecho a una jubilación digna.

En conclusión, el trabajo dignifica, pero solo si se respetan los derechos de los trabajadores.
\end{textebox}

\textbf{Traduction} : « Le métier de mes rêves, c'est d'être menuisier, comme mon grand-père. Il m'attire parce que j'aime travailler avec mes mains et créer des objets utiles pour les gens de mon quartier. De plus, un bon menuisier trouve toujours des clients à Yaoundé. Cependant, je pense que tout travailleur doit avoir des droits. D'abord, il a droit à un salaire juste et à un contrat clair. Ensuite, il mérite des jours de repos et des congés payés. Enfin, il doit pouvoir s'affilier à un syndicat pour défendre ses intérêts. Pour finir, après une vie d'efforts, il a droit à une retraite digne. En conclusion, le travail rend digne, mais seulement si l'on respecte les droits des travailleurs. »

\textbf{Barème indicatif (sur 20)} : respect du sujet et de la consigne (métier + trois droits + conclusion) 5 pts ; richesse et correction du lexique du travail 5 pts ; correction grammaticale (accords, temps, accents) 6 pts ; organisation et connecteurs (\esp{primero, segundo, además, sin embargo, en conclusión}) 2 pts ; respect du nombre de mots 2 pts.

\textbf{Points de vigilance} :
\begin{itemize}
\item Citer au moins \textbf{trois droits} distincts avec le lexique de la leçon (\esp{sueldo justo, contrato claro, descanso, vacaciones pagadas, sindicato, jubilación}).
\item Employer \esp{tener derecho a} + nom ou infinitif, et non un calque du français.
\item Soigner les accents : \esp{jubilación, además, después, conclusión, digno}.
\item Structurer en trois paragraphes : présentation du métier, droits, conclusion personnelle.
\item Compter les mots : entre 80 et 100 ; une production trop courte est pénalisée.
\end{itemize}
\end{corrige}

\newpage

\coursec{Fiche bilan}

\begin{popfigure}
\begin{tikzpicture}[mindmap, grow cyclic, every node/.style={concept, text=white, font=\small}, concept color=popblue, level 1/.append style={level distance=3.2cm, sibling angle=51}, root concept/.append style={font=\small\bfseries}]
\node[concept] {El mundo del trabajo}
  child[concept color=poporange] { node {Oficios y profesiones} }
  child[concept color=popgreen] { node {Empleo y paro} }
  child[concept color=poppurple] { node {Sueldo y contrato} }
  child[concept color=poppink] { node {Sindicato y huelga} }
  child[concept color=popgold] { node {Derechos del trabajador} }
  child[concept color=popdark] { node {Deberes del trabajador} }
  child[concept color=popblue!70] { node {Comentario de texto} };
\end{tikzpicture}
\legende{Carte mentale de la leçon : le monde du travail, ses acteurs, ses droits et la méthode du commentaire.}
\end{popfigure}

\begin{bilan}
\textbf{Lexique clé à connaître par cœur}

\begin{center}
\begin{tabularx}{\linewidth}{|X|X|X|X|}
\hline
\rowcolor{popblueL} Español & Français & Español & Français \\
\hline
el oficio & le métier (manuel) & el paro / el desempleo & le chômage \\
\hline
la profesión & la profession & el sueldo / el salario & le salaire \\
\hline
el empleo & l'emploi & el contrato & le contrat \\
\hline
el/la trabajador(a) & le/la travailleur(se) & el sindicato & le syndicat \\
\hline
el/la empleado(a) & l'employé(e) & la huelga & la grève \\
\hline
el/la jefe(a) & le/la chef(fe) & la jubilación & la retraite \\
\hline
la empresa & l'entreprise & el despido & le licenciement \\
\hline
el horario & l'horaire & las vacaciones & les congés \\
\hline
\end{tabularx}
\end{center}

\textbf{Expressions utiles}

\begin{itemize}
\item \esp{trabajar a tiempo completo / a tiempo parcial} : travailler à temps plein / partiel.
\item \esp{buscar empleo / encontrar trabajo} : chercher un emploi / trouver du travail.
\item \esp{estar en paro / estar en huelga} : être au chômage / en grève.
\item \esp{firmar un contrato} : signer un contrat ; \esp{jubilarse} : prendre sa retraite.
\item \esp{reivindicar mejores condiciones laborales} : revendiquer de meilleures conditions de travail.
\end{itemize}

\textbf{Grammaire à revoir pour ce thème}

\begin{center}
\begin{tabularx}{\linewidth}{|X|X|}
\hline
\rowcolor{popblueL} Point de grammaire & À retenir \\
\hline
\esp{ser} / \esp{estar} & \esp{ser} + profession : \esp{Soy profesora.} (sans article) ; \esp{estar} + situation temporaire : \esp{Está en paro.} \\
\hline
Verbes en \esp{-ar} au présent & \esp{trabajar} : trabajo, trabajas, trabaja, trabajamos, trabajáis, trabajan. \\
\hline
Futur simple & \esp{trabajaré, trabajarás, trabajará...} : projets professionnels. \\
\hline
Subjonctif après \esp{es importante que} & \esp{Es importante que los trabajadores conozcan sus derechos.} \\
\hline
\esp{hay que} + infinitif & obligation générale : \esp{Hay que respetar el horario.} \\
\hline
\end{tabularx}
\end{center}

\textbf{Droits et devoirs : idées forces}

\begin{itemize}
\item \textbf{Derechos} : \esp{un salario justo, un contrato, vacaciones pagadas, seguridad en el trabajo, afiliarse a un sindicato, la jubilación.}
\item \textbf{Deberes} : \esp{cumplir el horario, respetar a los compañeros, trabajar con responsabilidad, cuidar el material, guardar el secreto profesional.}
\end{itemize}

\textbf{Méthode express du commentaire de texte}

\begin{enumerate}
\item Lire deux fois ; présenter : type de texte, thème, source.
\item Dégager la \cle{idea principal} et les idées secondaires.
\item Repérer le lexique du travail et les procédés (exemples, chiffres, opposition derechos/deberes).
\item Donner son avis argumenté : \esp{En mi opinión... / Estoy de acuerdo porque...}
\item Conclure en ouvrant sur le contexte camerounais ou africain.
\end{enumerate}
\end{bilan}

\begin{aretenir}[Checklist avant le Bac]
\begin{itemize}
\item[$\square$] Je connais le lexique : \esp{oficio, profesión, empleo, paro, sueldo, contrato, sindicato, huelga, jubilación}.
\item[$\square$] Je distingue \esp{oficio} (métier manuel) et \esp{profesión} (profession qualifiée).
\item[$\square$] Je dis une profession avec \esp{ser} sans article : \esp{Mi padre es carpintero.}
\item[$\square$] J'emploie \esp{estar en paro}, \esp{estar en huelga} sans faute.
\item[$\square$] Je cite au moins trois droits et trois devoirs des travailleurs en espagnol.
\item[$\square$] Je conjugue \esp{trabajar} au présent et au futur sans erreur.
\item[$\square$] J'utilise \esp{hay que} + infinitif pour exprimer l'obligation.
\item[$\square$] Je sais employer le subjonctif après \esp{es importante que}.
\item[$\square$] Je respecte les accents : \esp{profesión, jubilación, sindicato}.
\item[$\square$] Je sais structurer un commentaire : présentation, idées, avis personnel, conclusion.
\item[$\square$] J'évite le faux ami : \esp{el paro} = le chômage (et non « la panne »).
\item[$\square$] Je peux relier le texte au monde du travail au Cameroun (agriculteurs, comerciantes, funcionarios).
\end{itemize}
\end{aretenir}

\findecours

\end{document}
